% Présenter son CV dans le cadre d’un entretien d’embauche — Français Tle Tech — Cours signé OnBuch+
% Programme officiel MINESEC. Compiler avec : tectonic N24-presenter-son-cv-cadre-entretien-embauche.tex
\newcommand{\DOCMATIERE}{Français}
\newcommand{\DOCNIVEAU}{Tle Tech}
\newcommand{\DOCMODULE}{Français – classe de Terminale technique}
\newcommand{\DOCLECON}{N24}
\newcommand{\DOCTITRE}{Présenter son CV dans le cadre d’un entretien d’embauche}
% OnBuch+ — préambule « cours signés » ENRICHI (compilé avec tectonic / XeLaTeX)
% Système visuel « Pop » d'OnBuch+ : fond crème (jamais blanc), bordures encre
% épaisses, ombres « dures » décalées, titres Archivo Black en capitales.
% Version MATHÉMATIQUES : graphes (pgfplots/tikz), boîtes pédagogiques
% (définition, propriété, méthode, attention, le savais-tu, exercice résolu,
% exercice, TP, bilan), page de garde et signature OnBuch+ ancrée en bas.
%
% Macros attendues dans main.tex AVANT \input{preamble} :
%   \DOCMATIERE  \DOCNIVEAU  \DOCMODULE  \DOCLECON (numéro)  \DOCTITRE
\documentclass[11pt]{article}

\usepackage{fontspec}
\usepackage[french]{babel}
\usepackage[a4paper,margin=2cm,top=3.4cm,bottom=2.8cm,headheight=1.4cm,footskip=1.3cm]{geometry}
\usepackage{amsmath,amssymb,amsfonts}
\usepackage{mathtools} % \xmapsto, \coloneqq... souvent utilisés par les modèles
\usepackage{cancel} % simplifications barrées (bilans de forces)
% \abs{z} (module, valeur absolue) et \norm{u} (norme), utilisés par les modèles
\providecommand{\abs}{}\let\abs\relax\DeclarePairedDelimiter{\abs}{\lvert}{\rvert}
\providecommand{\norm}{}\let\norm\relax\DeclarePairedDelimiter{\norm}{\lVert}{\rVert}
\usepackage[table]{xcolor} % [table] : fournit aussi \rowcolors (alternance de lignes)
\usepackage{pagecolor}
\usepackage{tikz}
\usetikzlibrary{arrows.meta,positioning,calc,shapes.geometric,shapes.misc,shapes.arrows,decorations.pathmorphing,decorations.pathreplacing,decorations.markings,patterns,shadows,mindmap,trees,fit,backgrounds,matrix,intersections,angles,quotes}
\usepackage{pgfplots}
\usepackage[version=4]{mhchem} % équations nucléaires \ce{^{235}_{92}U}
\usepackage[european,siunitx]{circuitikz} % schémas électriques
\pgfplotsset{compat=1.18}
\usepgfplotslibrary{fillbetween,groupplots} % aires entre courbes (calcul intégral)
\usepackage{enumitem}
\usepackage{fancyhdr}
% [most] charge toutes les bibliothèques tcolorbox (listings, minted, raster,
% documentation...) et gonfle inutilement la mémoire TeX sur un document long
% (~300 boîtes) : on ne charge que ce qui est réellement utilisé.
\usepackage[skins,breakable]{tcolorbox}
\usepackage{graphicx}
\usepackage{colortbl}
\usepackage{booktabs}
\usepackage{tabularx}
\usepackage{multirow}
\usepackage{diagbox} % case d'en-tête diagonale des tableaux à double entrée
% Colonnes extensibles centrée / gauche / droite (C, L, R), souvent utilisées par les modèles
\newcolumntype{C}{>{\centering\arraybackslash}X}
\newcolumntype{L}{>{\raggedright\arraybackslash}X}
\newcolumntype{R}{>{\raggedleft\arraybackslash}X}
\usepackage{array}
\usepackage{multicol}
\usepackage{siunitx}
% parse-numbers=false : accepte aussi \SI{2,5 \times 10^{-3}}{...}
\sisetup{locale=FR,output-decimal-marker={,},group-minimum-digits=4,parse-numbers=false}
\DeclareSIUnit\ohm{\ensuremath{\symup{\Omega}}} % U+2126 absent de la police
\DeclareSIUnit\division{div} % réglages d'oscilloscope (ms/div, V/div)
\DeclareSIUnit\tr{tr} % tours (vitesse de rotation, ex. \SI{1500}{\tr\per\minute})
\DeclareSIUnit\molar{M} % concentration molaire (ex. \SI{3}{\molar}, \SI{1}{\micro\molar})
\DeclareSIUnit\an{an} % année (le modèle confond parfois avec \year, un compteur TeX) — ex. \SI{5}{\kilo\meter\per\an}
\DeclareSIUnit\FCFA{FCFA} % franc CFA (ex. \SI{500}{\FCFA\per\kilo\gram})
\DeclareSIUnit\degreeBrix{\textdegree Brix} % teneur en sucre d'un sirop/jus (réfractométrie)
\DeclareSIUnit\month{mois} % le modèle utilise parfois \month, un compteur TeX, comme unité de durée
\DeclareSIUnit\microliter{\micro\liter} % le modèle écrit parfois \microliter en un seul token
\DeclareSIUnit\million{millions} % dénombrement (ex. \SI{15}{\million\per\mL} spermatozoïdes)
\usepackage{qrcode}
% \degree : commande fréquemment utilisée par les modèles pour un angle ou une
% température en dehors de \SI{}{} (non fournie par les paquets chargés).
\providecommand{\degree}{^{\circ}}
% \qed (fin de démonstration), utilisé par les modèles sans amsthm
\providecommand{\qed}{\hfill$\square$}
% \barre : double barre des valeurs interdites dans un tableau de variations (tkz-tab)
\providecommand{\barre}{\ensuremath{\|}}
% environnement proof (amsthm non chargé) utilisé par les modèles
\ifdefined\proof\else\newenvironment{proof}[1][Démonstration]{\par\noindent\textit{#1.}\ }{\qed\par}\fi
\usepackage{lastpage}
\usepackage{hyperref}
\hypersetup{hidelinks}

% ============================================================
% Palette « Pop » OnBuch+ — mêmes hex que la classe P (plus_ui.dart)
% ============================================================
\definecolor{poporange}{HTML}{F59321}
\definecolor{popdark}{HTML}{E07A0C}
\definecolor{popcream}{HTML}{FFF6E8}
\definecolor{popink}{HTML}{1C1714}
\definecolor{poppurple}{HTML}{7A5AE0}
\definecolor{popgreen}{HTML}{2E9E6A}
\definecolor{poppink}{HTML}{E0457B}
\definecolor{popblue}{HTML}{2F7DE1}
\definecolor{popgold}{HTML}{C98A12}
\definecolor{popmuted}{HTML}{8A7F76}
\definecolor{popline}{HTML}{EADFD2}
\definecolor{popgray}{HTML}{8A7F76}
\colorlet{popgrey}{popgray}
% Alias tolérants (le modèle nomme parfois les couleurs différemment)
\colorlet{popwhite}{white}
\colorlet{popblack}{popink}
\colorlet{popred}{poppink}
\colorlet{popyellow}{popgold}
% Teintes claires (fonds de tableaux/encadrés) — le modèle nomme parfois
% les couleurs avec un suffixe L (ex. popblueL) sans les avoir définies.
\colorlet{poporangeL}{poporange!20}
\colorlet{popdarkL}{popdark!20}
\colorlet{poppurpleL}{poppurple!20}
\colorlet{popgreenL}{popgreen!20}
\colorlet{poppinkL}{poppink!20}
\colorlet{popblueL}{popblue!20}
\colorlet{popgoldL}{popgold!20}
\colorlet{popredL}{popred!20}
\colorlet{popgrayL}{popgray!20}
% Teintes claires pour fonds de tableaux / zones
\colorlet{popblueL}{popblue!10!white}
\colorlet{poporangeL}{poporange!14!white}
\colorlet{popgreenL}{popgreen!12!white}
\colorlet{poppurpleL}{poppurple!12!white}
\colorlet{poppinkL}{poppink!12!white}
\colorlet{popgoldL}{popgold!14!white}

\pagecolor{popcream}

% ============================================================
% Typographie
% ============================================================
\defaultfontfeatures{Ligatures=TeX}
\setmainfont{PlusJakartaSans-Medium.ttf}[Path=../../../fonts/,BoldFont=PlusJakartaSans-Bold.ttf]
% Maths en sans-serif (Fira Math) avec chiffres et lettres droites de Plus
% Jakarta, pour que les formules se fondent dans le texte.
\usepackage[math-style=ISO,bold-style=ISO]{unicode-math}
\setmathfont{FiraMath-Regular.otf}[Path=../../../fonts/]
\setmathfont{PlusJakartaSans-Medium.ttf}[Path=../../../fonts/,range={up/num,up/{Latin,latin},\mathup}]
\usepackage{icomma} % virgule décimale sans espace en mode math
% Caractères mathématiques Unicode absents des polices de texte (Plus Jakarta,
% Archivo Black) : rendus par leur équivalent mathématique, en texte comme en
% formule — sinon ils s'affichent en carré vide (ex. « Arithmétique dans ℤ »).
\usepackage{newunicodechar}
\newunicodechar{ℝ}{\ensuremath{\mathbb{R}}}
\newunicodechar{ℤ}{\ensuremath{\mathbb{Z}}}
\newunicodechar{ℂ}{\ensuremath{\mathbb{C}}}
\newunicodechar{ℕ}{\ensuremath{\mathbb{N}}}
\newunicodechar{ℚ}{\ensuremath{\mathbb{Q}}}
\newunicodechar{𝔻}{\ensuremath{\mathbb{D}}}
\newunicodechar{α}{\ensuremath{\alpha}}
\newunicodechar{β}{\ensuremath{\beta}}
\newunicodechar{γ}{\ensuremath{\gamma}}
\newunicodechar{δ}{\ensuremath{\delta}}
\newunicodechar{ε}{\ensuremath{\varepsilon}}
\newunicodechar{θ}{\ensuremath{\theta}}
\newunicodechar{λ}{\ensuremath{\lambda}}
\newunicodechar{μ}{\ensuremath{\mu}}
\newunicodechar{σ}{\ensuremath{\sigma}}
\newunicodechar{τ}{\ensuremath{\tau}}
\newunicodechar{φ}{\ensuremath{\varphi}}
\newunicodechar{ψ}{\ensuremath{\psi}}
\newunicodechar{ω}{\ensuremath{\omega}}
\newunicodechar{Σ}{\ensuremath{\Sigma}}
% majuscules produites par \MakeUppercase dans les titres (α-aminés -> Α)
\newunicodechar{Α}{\ensuremath{\alpha}}
\newunicodechar{Β}{\ensuremath{\beta}}
\newunicodechar{Γ}{\ensuremath{\Gamma}}
\newunicodechar{Δ}{\ensuremath{\Delta}}
\newunicodechar{Ε}{\ensuremath{\varepsilon}}
\newunicodechar{Θ}{\ensuremath{\Theta}}
\newunicodechar{Λ}{\ensuremath{\Lambda}}
\newunicodechar{Μ}{\ensuremath{\mu}}
\newunicodechar{Τ}{\ensuremath{\tau}}
\newunicodechar{Φ}{\ensuremath{\Phi}}
\newunicodechar{Ψ}{\ensuremath{\Psi}}
\newunicodechar{Ω}{\ensuremath{\Omega}}
\newunicodechar{↦}{\ensuremath{\mapsto}}
\newunicodechar{∈}{\ensuremath{\in}}
\newunicodechar{∉}{\ensuremath{\notin}}
\newunicodechar{∧}{\ensuremath{\wedge}}
\newunicodechar{⇔}{\ensuremath{\Leftrightarrow}}
\newunicodechar{⟺}{\ensuremath{\Longleftrightarrow}}
\newunicodechar{⇒}{\ensuremath{\Rightarrow}}
\newunicodechar{⟹}{\ensuremath{\Longrightarrow}}
\newunicodechar{∘}{\ensuremath{\circ}}
\newunicodechar{∪}{\ensuremath{\cup}}
\newunicodechar{∩}{\ensuremath{\cap}}
\newunicodechar{≡}{\ensuremath{\equiv}}
\newunicodechar{⊂}{\ensuremath{\subset}}
\newunicodechar{⊥}{\ensuremath{\perp}}
\newunicodechar{∃}{\ensuremath{\exists}}
\newunicodechar{∀}{\ensuremath{\forall}}
\newunicodechar{∛}{\ensuremath{\sqrt[3]{\,}}}
\newunicodechar{∜}{\ensuremath{\sqrt[4]{\,}}}
\newunicodechar{⃗}{\ensuremath{{}^{\to}}}
\newunicodechar{ⁿ}{\ifmmode^{n}\else\textsuperscript{\ensuremath{n}}\fi}
\newunicodechar{ˣ}{\ifmmode^{x}\else\textsuperscript{\ensuremath{x}}\fi}
\newunicodechar{ᵇ}{\ifmmode^{b}\else\textsuperscript{\ensuremath{b}}\fi}
\newunicodechar{ʳ}{\ifmmode^{r}\else\textsuperscript{\ensuremath{r}}\fi}
\newunicodechar{ᵘ}{\ifmmode^{u}\else\textsuperscript{\ensuremath{u}}\fi}
\newunicodechar{ʸ}{\ifmmode^{y}\else\textsuperscript{\ensuremath{y}}\fi}
\newunicodechar{ᵃ}{\ifmmode^{a}\else\textsuperscript{\ensuremath{a}}\fi}
\newunicodechar{ᵉ}{\ifmmode^{e}\else\textsuperscript{\ensuremath{e}}\fi}
\newunicodechar{ᵏ}{\ifmmode^{k}\else\textsuperscript{\ensuremath{k}}\fi}
\newunicodechar{ᵖ}{\ifmmode^{p}\else\textsuperscript{\ensuremath{p}}\fi}
\newunicodechar{ᴺ}{\ifmmode^{N}\else\textsuperscript{\ensuremath{N}}\fi}
\newunicodechar{ᶜ}{\ifmmode^{c}\else\textsuperscript{\ensuremath{c}}\fi}
\newunicodechar{ʰ}{\ifmmode^{h}\else\textsuperscript{\ensuremath{h}}\fi}
\newunicodechar{ᵠ}{\ifmmode^{q}\else\textsuperscript{\ensuremath{q}}\fi}
\newunicodechar{ₙ}{\ifmmode_{n}\else\textsubscript{\ensuremath{n}}\fi}
\newunicodechar{ₐ}{\ifmmode_{a}\else\textsubscript{\ensuremath{a}}\fi}
\newunicodechar{ᵢ}{\ifmmode_{i}\else\textsubscript{\ensuremath{i}}\fi}
\newunicodechar{ᵣ}{\ifmmode_{r}\else\textsubscript{\ensuremath{r}}\fi}
\newunicodechar{ₖ}{\ifmmode_{k}\else\textsubscript{\ensuremath{k}}\fi}
\newunicodechar{ᵦ}{\ifmmode_{\beta}\else\textsubscript{\ensuremath{\beta}}\fi}
\newunicodechar{ₓ}{\ifmmode_{x}\else\textsubscript{\ensuremath{x}}\fi}
\newfontfamily\popdisplay{ArchivoBlack-Regular.ttf}[Path=../../../fonts/]
\newfontfamily\poplogo{SpaceGrotesk-Bold.ttf}[Path=../../../fonts/]
\newfontfamily\popsemi{PlusJakartaSans-SemiBold.ttf}[Path=../../../fonts/]
\newfontfamily\popxbold{PlusJakartaSans-ExtraBold.ttf}[Path=../../../fonts/]
\newfontfamily\popxboldls{PlusJakartaSans-ExtraBold.ttf}[Path=../../../fonts/,LetterSpace=3.0]

\setlist[enumerate]{leftmargin=1.7em,itemsep=5pt,topsep=4pt}
\setlist[itemize]{leftmargin=1.7em,itemsep=4pt,topsep=4pt,label={\color{poporange}\textbullet}}
\setlength{\parindent}{0pt}
\setlength{\parskip}{5pt}
\renewcommand{\arraystretch}{1.25}

% pgfplots : style Pop par défaut
\pgfplotsset{
  every axis/.append style={axis line style={popink,line width=1pt},tick style={popink},
    grid style={popline,line width=0.5pt},label style={font=\small},tick label style={font=\footnotesize}},
  popcurve/.style={poporange,line width=1.8pt,no markers,smooth},
}
% Même style disponible dans un \draw TikZ simple (hors pgfplots)
\tikzset{popcurve/.style={poporange,line width=1.8pt}}
\tikzset{popgrid/.style={popline,very thin}}
\colorlet{popcurve}{poporange} % le nom du style est parfois utilisé comme couleur

% ============================================================
% Pill / squiggle
% ============================================================
\newcommand{\poppill}[3]{%
  \tikz[baseline=-0.55ex]\node[draw=popink,line width=1.2pt,rounded corners=2.6pt,%
    fill=#2,inner xsep=6.5pt,inner ysep=3pt]{%
    \popxboldls\fontsize{7}{7}\selectfont\color{#3}\MakeUppercase{#1}};%
}
\newcommand{\popsquiggle}[1]{%
  \tikz[baseline=-0.4ex]{
    \draw[#1,line width=2.6pt,line cap=round]
      plot [smooth] coordinates {(0,0.09) (0.34,0.01) (0.68,0.21) (1.02,0.01) (1.36,0.10)};
  }%
}
% Mot-clé surligné (marqueur orange sous le texte)
\newcommand{\cle}[1]{{\popxbold\color{popdark}#1}}

% ============================================================
% En-tête / pied
% ============================================================
\pagestyle{fancy}
\fancyhf{}
\renewcommand{\headrulewidth}{0pt}
\renewcommand{\footrulewidth}{0pt}
\lhead{\raisebox{-0.15ex}{\poplogo\fontsize{13}{13}\selectfont\color{popink}OnBuch\color{poporange}+}}
\rhead{\poppill{\DOCMATIERE\ \textbullet\ \DOCNIVEAU\ \textbullet\ Leçon \DOCLECON}{white}{popink}}
\lfoot{\popxbold\fontsize{7.6}{7.6}\selectfont\color{popmuted}\MakeUppercase{Cours signé OnBuch+}\ \textbullet\ {\popsemi\DOCTITRE}}
\rfoot{\tikz[baseline=-0.6ex]\node[rounded corners=8.5pt,fill=popink,minimum height=17pt,minimum width=17pt,inner xsep=5pt,inner sep=0]{\popxbold\fontsize{8}{8}\selectfont\color{popcream}\thepage\,/\,\pageref*{LastPage}};}

% ============================================================
% Titres
% ============================================================
\newcommand{\chaptitre}[2]{%
  \begin{tcolorbox}[enhanced,colback=poporange,colframe=popink,boxrule=2.6pt,arc=11pt,
      drop shadow=popink,left=15pt,right=15pt,top=12pt,bottom=11pt,before skip=3pt,after skip=18pt]
    \poppill{#2}{popink}{popcream}\par\vspace{9pt}
    {\raggedright\hyphenpenalty=10000\exhyphenpenalty=10000\popdisplay\fontsize{22}{24}\selectfont\color{popcream}\MakeUppercase{#1}\par}
  \end{tcolorbox}
}
% Section numérotée : pastille orange avec le numéro + titre Archivo
\newcounter{popsec}
\newcommand{\coursec}[1]{%
  \stepcounter{popsec}\setcounter{popsub}{0}%
  \par\vspace{16pt}\noindent
  \begin{minipage}{\linewidth}
  \tikz[baseline=-0.7ex]\node[draw=popink,line width=1.6pt,fill=poporange,rounded corners=4pt,minimum size=22pt,inner sep=0]{\popdisplay\fontsize{12}{12}\selectfont\color{popcream}\thepopsec};%
  \hspace{8pt}{\popdisplay\fontsize{15}{17}\selectfont\color{popink}\MakeUppercase{#1}}\par
  \vspace{4pt}\noindent{\color{poporange}\rule{36pt}{3.6pt}}
  \end{minipage}\par\nopagebreak\vspace{8pt}
}
\newcounter{popsub}
\newcommand{\courssub}[1]{%
  \stepcounter{popsub}%
  \par\vspace{9pt}\noindent{\popxbold\fontsize{11.5}{13}\selectfont\color{popink}{\color{poporange}\thepopsec.\thepopsub}\hspace{6pt}#1}\par\nopagebreak\vspace{4pt}
}

% ============================================================
% Boîtes pédagogiques — même recette : fond blanc, bordure épaisse colorée,
% ombre dure encre, badge plein. Titre optionnel [#1].
% ============================================================
\tcbset{popbase/.style={breakable,enhanced,colback=white,boxrule=2.3pt,arc=9pt,
  shadow={3pt}{-3pt}{0pt}{popink},left=14pt,right=14pt,top=11pt,bottom=11pt,before skip=11pt,after skip=15pt,
  fontupper=\popsemi\fontsize{10.6}{14.5}\selectfont\color{popink},
  fontlower=\popsemi\fontsize{10.6}{14.5}\selectfont\color{popink}}}
% Coche dessinée (le glyphe ✓ n'existe pas dans les polices du document)
\newcommand{\popcheck}[1]{\tikz[baseline=-0.5ex]\draw[#1,line width=1.6pt,line cap=round,line join=round](-0.13,0.0)--(-0.04,-0.09)--(0.13,0.1);}
\providecommand{\coche}{\popcheck{popgreen}}
\providecommand{\resultat}[1]{\par\smallskip\noindent\fcolorbox{popgreen}{popgreenL}{\parbox{\dimexpr\linewidth-2\fboxsep-2\fboxrule\relax}{\textbf{Résultat.} #1}}\par\smallskip}
\newcommand{\popboxhead}[3]{\poppill{#1}{#2}{white}\ifx\relax#3\relax\else\hspace{6pt}{\popxbold\fontsize{10.6}{12}\selectfont\color{popink}#3}\fi\par\vspace{7pt}}

\newtcolorbox{aretenir}[1][]{popbase,colframe=popgreen,before upper={\popboxhead{À retenir}{popgreen}{#1}}}
\newtcolorbox{exemplebox}[1][]{popbase,colframe=poporange,before upper={\popboxhead{Exemple}{poporange}{#1}}}
\newtcolorbox{definition}[1][]{popbase,colframe=popblue,before upper={\popboxhead{Définition}{popblue}{#1}}}
\newtcolorbox{propriete}[1][]{popbase,colframe=poppurple,before upper={\popboxhead{Propriété}{poppurple}{#1}}}
\newtcolorbox{methode}[1][]{popbase,colframe=popgold,before upper={\popboxhead{Méthode}{popgold}{#1}}}
\newtcolorbox{attention}[1][]{popbase,colframe=poppink,before upper={\popboxhead{Attention}{poppink}{#1}}}
\newtcolorbox{experience}[1][]{popbase,colframe=popblue,colback=popblueL,before upper={\popboxhead{Expérience}{popblue}{#1}}}
% « Le savais-tu ? » : carte violette pleine (contexte, culture, Cameroun)
\newtcolorbox{savaistu}[1][]{popbase,colback=poppurple,colframe=popink,
  fontupper=\popsemi\fontsize{10.4}{14.2}\selectfont\color{white},
  before upper={\poppill{Le savais-tu\,?}{white}{poppurple}\ifx\relax#1\relax\else\hspace{6pt}{\popxbold\color{white}#1}\fi\par\vspace{7pt}}}
% Exercice résolu : énoncé en haut, \tcblower puis la solution sur fond crème
\newcounter{popexo}\newcounter{popexr}
\newtcolorbox{exoresolu}[1][]{popbase,colframe=poporange,colbacklower=popcream,
  segmentation style={popink,line width=1.2pt,dashed},
  before upper={\stepcounter{popexr}\popboxhead{Exercice résolu \thepopexr}{poporange}{#1}},
  before lower={\poppill{Solution}{popink}{popcream}\par\vspace{6pt}}}
% Exercice (non résolu en place) — la correction est dans la partie Corrigés
\newtcolorbox{exercice}[1][]{popbase,colframe=popink,
  before upper={\stepcounter{popexo}\popboxhead{Exercice \thepopexo}{popink}{#1}}}
% Corrigé
\newtcolorbox{corrige}[1][]{popbase,colframe=popgreen,colback=popgreenL,
  before upper={\popboxhead{Corrigé}{popgreen}{#1}}}
% Objectifs / prérequis / bilan
\newtcolorbox{objectifs}{popbase,colframe=popink,colback=poporangeL,
  before upper={\popboxhead{Objectifs de la leçon}{popink}{}}}
\newtcolorbox{prerequis}{popbase,colframe=popink,colback=popblueL,
  before upper={\popboxhead{Prérequis}{popblue}{}}}
\newtcolorbox{bilan}{popbase,colframe=popink,colback=white,boxrule=2.8pt,
  before upper={\popboxhead{Fiche bilan}{poporange}{L'essentiel à connaître pour l'examen}}}
% Figure encadrée avec légende
\newtcolorbox{popfigure}[1][]{enhanced,breakable=false,colback=white,colframe=popline,boxrule=1.4pt,arc=8pt,
  left=10pt,right=10pt,top=10pt,bottom=8pt,before skip=10pt,after skip=12pt,halign=center,
  lower separated=false,halign lower=center,
  fontlower=\popsemi\fontsize{9}{11.5}\selectfont\color{popmuted},
  #1}
\newcommand{\legende}[1]{\par\vspace{6pt}{\popsemi\fontsize{9}{11.5}\selectfont\color{popmuted}#1}}

% ============================================================
% Signature OnBuch+
% ============================================================
\newcommand{\poplogoinline}[1]{{\poplogo\fontsize{#1}{#1}\selectfont\color{popink}OnBuch\color{poporange}+}}
\newcommand{\signatureonbuch}{%
  \begin{tcolorbox}[enhanced,colback=popink,colframe=popink,boxrule=2.2pt,arc=10pt,
      drop shadow=poporange,left=14pt,right=14pt,top=10pt,bottom=10pt,before skip=0pt,after skip=0pt]
    \tikz[baseline=-0.7ex]\node[circle,fill=popgreen,draw=popcream,line width=1.3pt,minimum size=22pt,inner sep=0]{\popcheck{white}};%
    \hspace{9pt}\begin{minipage}[c]{0.62\linewidth}
      {\popxbold\fontsize{12}{13}\selectfont\color{popcream}Signé\ \textbullet\ {\poplogo OnBuch\color{poporange}+}}\par\vspace{2pt}
      {\raggedright\popsemi\fontsize{8}{10.5}\selectfont\color{popline}\MakeUppercase{Équipe pédagogique OnBuch+}\\\MakeUppercase{Conforme au programme officiel MINESEC}\par}
    \end{minipage}\hfill
    {\poplogo\fontsize{22}{22}\selectfont\color{popcream}OnBuch\color{poporange}+}
  \end{tcolorbox}%
}

% ============================================================
% Page de garde
% #1 titre · #2 matière · #3 niveau · #4 module · #5 n° leçon · #6 durée ·
% #7 description / objectifs (texte ou itemize)
% ============================================================
\newcommand{\pagegarde}[7]{%
  \thispagestyle{empty}
  \newgeometry{a4paper,margin=1.7cm,top=1.6cm,bottom=1.4cm}%
  \begin{tikzpicture}[remember picture,overlay]
    \fill[poporange!18] ([xshift=-3.2cm,yshift=-3.6cm]current page.north east) circle (4.6cm);
    \fill[poppurple!14] ([xshift=2.4cm,yshift=7.6cm]current page.south west) circle (3.8cm);
  \end{tikzpicture}%
  {\poplogo\fontsize{30}{30}\selectfont\color{popink}OnBuch\color{poporange}+}\hfill
  \raisebox{0.6ex}{\poppill{Cours signé \textbullet\ Premium}{popink}{popcream}}\par
  \vspace{1pt}\popsquiggle{poppurple}\par
  \vspace{8pt}
  \begin{tcolorbox}[enhanced,colback=poporange,colframe=popink,boxrule=2.8pt,arc=13pt,
      drop shadow=popink,left=18pt,right=18pt,top=12pt,bottom=13pt,after skip=10pt]
    \poppill{#2 \textbullet\ #3}{popink}{popcream}\hspace{5pt}\poppill{#4}{white}{popink}\par\vspace{8pt}
    {\popdisplay\fontsize{14}{14}\selectfont\color{popink}LEÇON #5}\par\vspace{4pt}
    {\raggedright\hyphenpenalty=10000\exhyphenpenalty=10000\popdisplay\fontsize{25}{27}\selectfont\color{popcream}\MakeUppercase{#1}\par}
    \vspace{8pt}
    {\popxbold\fontsize{9.5}{9.5}\selectfont\color{popink}\MakeUppercase{Durée conseillée : #6}}
  \end{tcolorbox}
  \begin{tcolorbox}[enhanced,colback=white,colframe=popink,boxrule=2.2pt,arc=10pt,
      drop shadow=popink,left=16pt,right=16pt,top=10pt,bottom=10pt,after skip=8pt]
    \poppill{Dans cette leçon}{poppurple}{white}\par\vspace{5pt}
    \popsemi\fontsize{9.3}{12.6}\selectfont\color{popink}#7
  \end{tcolorbox}
  \noindent\begin{minipage}[t]{0.487\linewidth}
    \begin{tcolorbox}[enhanced,colback=popgreen,colframe=popink,boxrule=2.2pt,arc=10pt,
        drop shadow=popink,left=11pt,right=11pt,top=10pt,bottom=10pt,height=3.1cm,valign=center]
      \tcbox[colback=white,colframe=popink,boxrule=1.3pt,arc=5pt,boxsep=0pt,nobeforeafter,
        left=3pt,right=3pt,top=3pt,bottom=3pt,valign=center]{\qrcode[height=1.9cm]{https://play.google.com/store/apps/details?id=com.luvvix.onbuch&hl=fr}}%
      \hspace{8pt}\begin{minipage}[c]{0.5\linewidth}
        {\popdisplay\fontsize{11}{12.5}\selectfont\color{white}\MakeUppercase{Télécharge OnBuch}}\par\vspace{4pt}
        {\raggedright\popsemi\fontsize{8.4}{11}\selectfont\color{white}Gratuit \textbullet\ Google Play\\Tuteur IA, annales, concours, résultats\par}
      \end{minipage}
    \end{tcolorbox}
  \end{minipage}\hfill
  \begin{minipage}[t]{0.487\linewidth}
    \begin{tcolorbox}[enhanced,colback=poppurple,colframe=popink,boxrule=2.2pt,arc=10pt,
        drop shadow=popink,left=11pt,right=11pt,top=10pt,bottom=10pt,height=3.1cm,valign=center]
      \tikz[baseline=-0.7ex]\node[circle,fill=white,draw=popink,line width=1.3pt,minimum size=1.6cm,inner sep=0]{\popdisplay\fontsize{24}{24}\selectfont\color{poppurple}+};%
      \hspace{8pt}\begin{minipage}[c]{0.6\linewidth}
        {\popdisplay\fontsize{11}{12.5}\selectfont\color{white}\MakeUppercase{Passe à OnBuch+}}\par\vspace{4pt}
        {\raggedright\popsemi\fontsize{8.4}{11}\selectfont\color{white}Tous les cours signés, Tuteur IA illimité, fiches PDF — onglet OnBuch+ de l'app\par}
      \end{minipage}
    \end{tcolorbox}
  \end{minipage}
  \vspace{8pt}
  \popcontenu
  \vfill
  \signatureonbuch
  \restoregeometry
  \newpage
}

% Rangée « Ce cours contient » (page de garde)
\newcommand{\poptile}[3]{% #1 couleur #2 gros texte #3 légende
  \begin{tcolorbox}[enhanced,colback=#1,colframe=popink,boxrule=2pt,arc=9pt,shadow={2.5pt}{-2.5pt}{0pt}{popink},
      left=6pt,right=6pt,top=9pt,bottom=9pt,halign=center,valign=center,height=1.75cm,width=\linewidth,nobeforeafter]
    {\popdisplay\fontsize{12.5}{13}\selectfont\color{white}\MakeUppercase{#2}}\par\vspace{3pt}
    {\centering\popsemi\fontsize{7.8}{9.5}\selectfont\color{white}#3\par}
  \end{tcolorbox}}
\newcommand{\popcontenu}{%
  \par\noindent{\popxboldls\fontsize{7.5}{7.5}\selectfont\color{popmuted}CE COURS SIGNÉ CONTIENT}\par\vspace{7pt}
  \noindent\begin{minipage}[t]{0.235\linewidth}\poptile{popblue}{Cours}{complet, illustré, pas à pas}\end{minipage}\hfill
  \begin{minipage}[t]{0.235\linewidth}\poptile{poporange}{Méthodes}{et exercices résolus}\end{minipage}\hfill
  \begin{minipage}[t]{0.235\linewidth}\poptile{poppink}{Exercices}{type examen + corrigés}\end{minipage}\hfill
  \begin{minipage}[t]{0.235\linewidth}\poptile{popgreen}{Fiche bilan}{l'essentiel à retenir}\end{minipage}\par}

% Sommaire
\newtcolorbox{sommaire}{enhanced,colback=white,colframe=popink,boxrule=2.2pt,arc=10pt,
  drop shadow=popink,left=18pt,right=18pt,top=13pt,bottom=13pt,before skip=6pt,after skip=20pt,
  before upper={\poppill{Sommaire}{poppurple}{white}\par\vspace{9pt}\popsemi\fontsize{10.6}{14.5}\selectfont\color{popink}}}

% Signature de fin de cours — poussée tout en bas de la dernière page
\newcommand{\findecours}{%
  \par\vfill\nopagebreak
  \begin{center}\popsquiggle{poporange}\end{center}\vspace{4pt}
  \signatureonbuch
}
\AtBeginDocument{\def\llbracket{\mathopen{[\mkern-3.5mu[}}\def\rrbracket{\mathclose{]\mkern-3.5mu]}}} % intervalles d'entiers (glyphe ⟦ absent de la police)
% Glyphes absents de Fira Math : redéfinis à partir de glyphes disponibles
\AtBeginDocument{%
  \def\setminus{\mathbin{\backslash}}\let\smallsetminus\setminus
  \def\vdots{\mathord{\vbox{\baselineskip3.2pt\lineskiplimit0pt\kern4pt\hbox{.}\hbox{.}\hbox{.}}}}%
  \def\ddots{\mathinner{\mkern1mu\raise6.5pt\hbox{.}\mkern2mu\raise3.6pt\hbox{.}\mkern2mu\raise0.7pt\hbox{.}\mkern1mu}}%
  \def\triangle{\mathord{\scalebox{0.9}{$\symup{\Delta}$}}}%
  \def\nabla{\mathord{\raisebox{\depth}{\scalebox{1}[-1]{$\symup{\Delta}$}}}}%
  \def\checkmark{\popcheck{popdark}}%
  \def\star{\mathbin{\ast}}\def\bigstar{\mathord{\ast}}%
  \def\diamond{\mathbin{\scalebox{0.6}{$\Diamond$}}}%
  \def\bigcup{\mathop{\mathchoice{\raisebox{-0.3em}{\scalebox{1.7}{$\cup$}}}{\scalebox{1.25}{$\cup$}}{\cup}{\cup}}\displaylimits}%
  \def\bigcap{\mathop{\mathchoice{\raisebox{-0.3em}{\scalebox{1.7}{$\cap$}}}{\scalebox{1.25}{$\cap$}}{\cap}{\cap}}\displaylimits}%
  \def\lesseqqgtr{\mathrel{\lessgtr}}\def\bigoplus{\mathop{\oplus}\displaylimits}%
}
\providecommand{\ssi}{si et seulement si} % abréviation fréquente
\AtBeginDocument{\providecommand{\wideparen}{\overparen}} % arc de cercle

%% FIGFIX-BEGIN v3 — correctifs globaux des figures (docs/onbuch-generation/tools/figfix)
\usepackage{adjustbox}\usepackage{etoolbox}
% toute figure TikZ/pgfplots plus large que la ligne est réduite à la largeur disponible
\BeforeBeginEnvironment{tikzpicture}{\begin{adjustbox}{max width=\linewidth}}
\AfterEndEnvironment{tikzpicture}{\end{adjustbox}}
% légendes pgfplots lisibles par-dessus les courbes
\pgfplotsset{every axis/.append style={legend style={fill=white,fill opacity=0.92,text opacity=1,draw=black!15,inner sep=3pt}}}
% étiquettes d'arêtes (syntaxe edge["..."]) avec fond blanc
\tikzset{every edge quotes/.append style={fill=white,inner sep=1.2pt}}
% bulles de cartes mentales : texte à la ligne dans la bulle, bulle un peu plus grande
\tikzset{every concept/.append style={text width=2.0cm,align=center,minimum size=2.3cm,inner sep=2pt}}
%% FIGFIX-END
\begin{document}

\pagegarde{Présenter son CV dans le cadre d’un entretien d’embauche}{Français}{Tle Tech}{Français – classe de Terminale technique}{N24}{9 séances (fiche de progression harmonisée)}{Cette leçon outille l'élève de Terminale technique pour concevoir, rédiger et présenter un curriculum vitae clair, structuré et convaincant. Elle mobilise la lecture de textes fonctionnels et iconiques, la maîtrise du subjonctif et du conditionnel, et culmine par une simulation d'entretien d'embauche devant un jury.
\vspace{4pt}
\begin{multicols}{2}\small
\begin{itemize}
\item À la fin de cette leçon, je sais identifier les éléments constitutifs d'un CV (état civil, formation, expérience, compétences, centres d'intérêt)
\item À la fin de cette leçon, je sais distinguer un CV efficace d'un CV défaillant à partir de modèles réels
\item À la fin de cette leçon, je sais interpréter des textes iconiques (photo d'identité, mise en page, pictogrammes) dans un CV
\item À la fin de cette leçon, je sais employer correctement le subjonctif et le conditionnel pour exprimer le souhait, l'hypothèse et la politesse dans un contexte professionnel
\item À la fin de cette leçon, je sais rédiger un CV complet, sans faute, adapté à une offre d'emploi précise
\item À la fin de cette leçon, je sais rédiger une lettre de motivation cohérente avec mon CV (complément utile au Bac)
\end{itemize}\end{multicols}}

\begin{sommaire}
\begin{multicols}{2}
\begin{enumerate}
\item Découverte : le CV, un texte fonctionnel
\item Les éléments constitutifs du CV
\item Textes iconiques et mise en page du CV
\item Langue : les valeurs du subjonctif et du conditionnel
\item Production guidée : rédiger son CV
\item Présenter son CV lors d'un entretien d'embauche
\item Intégration, compte rendu et remédiation
\item Activité d'intégration
\item Méthodes et exercices résolus
\item Exercices
\item Corrigés des exercices
\item Fiche bilan
\end{enumerate}
\end{multicols}
\end{sommaire}

\begin{objectifs}
\begin{itemize}
\item À la fin de cette leçon, je sais identifier les éléments constitutifs d'un CV (état civil, formation, expérience, compétences, centres d'intérêt)
\item À la fin de cette leçon, je sais distinguer un CV efficace d'un CV défaillant à partir de modèles réels
\item À la fin de cette leçon, je sais interpréter des textes iconiques (photo d'identité, mise en page, pictogrammes) dans un CV
\item À la fin de cette leçon, je sais employer correctement le subjonctif et le conditionnel pour exprimer le souhait, l'hypothèse et la politesse dans un contexte professionnel
\item À la fin de cette leçon, je sais rédiger un CV complet, sans faute, adapté à une offre d'emploi précise
\item À la fin de cette leçon, je sais rédiger une lettre de motivation cohérente avec mon CV (complément utile au Bac)
\item À la fin de cette leçon, je sais présenter oralement mon CV devant un jury dans le cadre d'un entretien d'embauche simulé
\item À la fin de cette leçon, je sais justifier mes choix de rédaction et de présentation devant un jury
\end{itemize}
\end{objectifs}

\begin{prerequis}
\begin{itemize}
\item Les types de textes (narratif, descriptif, fonctionnel) vus en classe de Première
\item La conjugaison des temps de l'indicatif (présent, passé composé, futur)
\item Les notions de base sur les modes (indicatif, impératif) vues au collège
\item La lettre administrative et la correspondance professionnelle (classe de Première)
\item Les techniques de prise de parole en public (exposés oraux)
\end{itemize}
\end{prerequis}

\newpage

\coursec{Découverte : le CV, un texte fonctionnel}

\courssub{Situation d'accroche et problématique}

À Douala, une entreprise de BTP publie une offre d'emploi de technicien en génie civil. Sur 250 candidatures reçues, le recruteur n'en retient que 12 après lecture des CV, puis convoque les candidats pour un entretien. Junior, titulaire d'un Baccalauréat technique, est éliminé dès la première étape : son CV, mal structuré et truffé de fautes, n'a pas convaincu en 30 secondes de lecture.

\begin{aretenir}[Problématique de la leçon]
Qu'est-ce qu'un bon CV ? Comment le présenter efficacement face à un jury lors d'un entretien d'embauche ?
\end{aretenir}

\courssub{Le CV : définition et nature de texte fonctionnel}

\begin{definition}[Texte fonctionnel]
Un \cle{texte fonctionnel} (ou texte utilitaire) est un écrit conçu pour remplir une fonction précise dans une situation de communication réelle : informer, demander, convaincre, ordonner, etc. Il répond à un besoin immédiat et vise une action concrète du destinataire.
\end{definition}

Le \cle{curriculum vitae} (CV) est l'archétype du texte fonctionnel à visée professionnelle. Étymologiquement, \emph{curriculum vitae} signifie « parcours de vie » en latin. C'est un document synthétique qui présente, de manière structurée et chronologique, le parcours scolaire, universitaire, professionnel et les compétences d'un candidat dans le but de postuler à un emploi, un stage ou une formation.

\begin{propriete}[Caractéristiques du CV comme texte fonctionnel]
\begin{itemize}
    \item \textbf{Finalité pragmatique :} obtenir un entretien d'embauche (pas l'emploi directement).
    \item \textbf{Brièveté :} une page maximum pour un jeune diplômé (deux pages acceptées pour profils expérimentés).
    \item \textbf{Structure codifiée :} rubriques obligatoires et ordre logique imposés par la convention professionnelle.
    \item \textbf{Objectivité :} faits vérifiables (diplômes, dates, noms d'entreprises), pas d'opinions subjectives.
    \item \textbf{Adaptabilité :} doit être personnalisé pour chaque offre ciblée.
\end{itemize}
\end{propriete}

\courssub{Le circuit de communication du CV}

Le CV s'inscrit dans une situation de communication asymétrique : le candidat (émetteur) s'adresse à un recruteur (récepteur) qu'il ne connaît pas, via un support écrit (le CV) qui doit déclencher une action (convocation à l'entretien).

\begin{popfigure}
\begin{tikzpicture}[
    node distance=1.8cm and 2.5cm,
    box/.style={rectangle, rounded corners=6pt, draw=popblue, thick, fill=popblueL, minimum width=2.8cm, minimum height=1.2cm, align=center, font=\small},
    arrow/.style={->, >=Stealth, thick, popdark},
    label/.style={font=\footnotesize, text=popmuted, align=center}
]
\node[box, align=center] (candidat){\textbf{CANDIDAT}\\(Émetteur)\\\textit{Parcours, compétences,}\\\textit{motivation}};
\node[box, right=of candidat, align=center] (cv){\textbf{CV}\\(Support écrit)\\\textit{Structuré, clair,}\\\textit{adapté, sans fautes}};
\node[box, right=of cv, align=center] (recruteur){\textbf{RECRUTEUR}\\(Récepteur)\\\textit{Analyse en 30--40 s}\\\textit{Critères : adéquation,}\\\textit{cohérence, crédibilité}};
\node[box, below=of cv, align=center] (entretien){\textbf{ENTRETIEN}\\(Objectif visé)\\\textit{Échange oral}\\\textit{Validation des compétences}\\\textit{Négociation}};
\draw[arrow] (candidat) -- node[above, label] {rédige / adapte} (cv);
\draw[arrow] (cv) -- node[above, label] {transmet / lit en 30--40 s} (recruteur);
\draw[arrow] (recruteur) -- node[right, label] {sélectionne / convoque} (entretien);
\draw[arrow, dashed] (entretien) -- node[left, label] {retour / décision} (candidat);
\end{tikzpicture}
\legende{Schéma du circuit de communication du CV : du candidat à l'entretien d'embauche. La flèche pointillée représente le retour (acceptation / refus).}
\end{popfigure}

\begin{aretenir}[Les trois fonctions du CV]
\begin{enumerate}
    \item \textbf{Informer :} donner au recruteur les faits bruts (identité, formation, expériences, compétences).
    \item \textbf{Convaincre :} montrer l'adéquation entre le profil du candidat et les besoins du poste.
    \item \textbf{Décrocher l'entretien :} le CV est une « clé » ; l'entretien est la « porte » vers l'embauche.
\end{enumerate}
\end{aretenir}

\courssub{Le temps de lecture : une contrainte majeure}

\begin{savaistu}[Donnée chiffrée clé]
Des études en recrutement (eye-tracking, enquêtes auprès de DRH) convergent : \textbf{un recruteur consacre en moyenne 30 à 40 secondes à la première lecture d'un CV}. Au Cameroun, dans les grands groupes (SNEC, CDC, banques, ministères, ONG), ce temps peut descendre à 20 secondes pour les postes très demandés.
\end{savaistu}

\begin{exemplebox}[Conséquences pratiques des 30 secondes]
En 30 secondes, le recruteur ne \emph{lit} pas : il \emph{scanne}. Il repère :
\begin{itemize}
    \item Le nom du dernier diplôme / poste (en haut à gauche ou au centre).
    \item Les mots-clés techniques du métier (ex. : \textit{AutoCAD, béton armé, DAO, gestion de chantier}).
    \item La cohérence chronologique (trous inexpliqués ?).
    \item La propreté visuelle (marges, police, fautes visibles).
\end{itemize}
Si l'un de ces signaux est négatif, le CV passe dans la pile « non retenu » sans seconde lecture.
\end{exemplebox}

\courssub{Lecture critique de deux CV contrastés}

Considérons l'offre réelle suivante, publiée à Yaoundé :

\begin{center}
\begin{tabularx}{\linewidth}{|X|}\hline
\rowcolor{popblueL}
\textbf{Offre d'emploi : Technicien en Génie Civil (BTP)} \\
\hline
\textbf{Entreprise :} SOTRACO BTP, Yaoundé \\
\textbf{Profil :} Baccalauréat technique (Génie Civil / BTP) ou BT/BTS. \\
\textbf{Compétences requises :} Lecture de plans (AutoCAD), suivi de chantier, connaissance matériaux (béton, acier), rédaction de rapports quotidiens, respect normes sécurité. \\
\textbf{Qualités :} Rigueur, esprit d'équipe, disponibilité immédiate. \\
\textbf{Dépôt :} CV + lettre de motivation (1 page max) par email. \\
\hline
\end{tabularx}
\end{center}

Deux candidats, \textbf{Alain} et \textbf{Junior}, tous deux titulaires du Baccalauréat technique Génie Civil, postulent. Voici l'analyse comparative de leurs CV.

\begin{popfigure}
\begin{center}
\begin{tabularx}{\linewidth}{|l|X|X|}\hline
\rowcolor{popblueL}
\textbf{Critère} & \textbf{CV d'Alain (retenu)} & \textbf{CV de Junior (rejeté)} \\ \hline
\textbf{Présentation générale} & Une page, police Calibri 11, marges 2 cm, interligne 1,15, titres en gras 14 pt, puces cohérentes. & Deux pages, trois polices différentes (Arial, Times, Comic Sans), marges irrégulières, texte justifié « en pavé », pas de puces. \\ \hline
\textbf{En-tête (identité)} & Nom, prénom, téléphone, email professionnel (\texttt{alain.n@email.cm}), LinkedIn, quartier/ville. Photo pro (optionnelle, mais soignée). & Nom seulement, email fantaisiste (\texttt{juniordu95@yahoo.fr}), pas de téléphone, pas de localisation, photo de vacances. \\ \hline
\textbf{Accroche / Titre} & \textit{Technicien Génie Civil — Baccalauréat Technique — Spécialité BTP — Disponible immédiatement} & \textit{Curriculum Vitae} (titre générique, inutile). \\ \hline
\textbf{Formation (antéchronologique)} & \textbf{2023--2024 :} Baccalauréat Technique Génie Civil, Lycée Technique de Nkolbisson (Mention AB).\newline \textbf{2020--2023 :} Cycle Secondaire, même établissement. & \textbf{2024 :} Bac G.C. (pas de mention, pas d'établissement).\newline \textbf{2021--2023 :} « Études secondaires » (vague, pas de série). \\ \hline
\textbf{Expériences / Stages} & \textbf{Stage 2024 (2 mois) :} Entreprise ETC BTP, Yaoundé.\newline -- Suivi quotidien avancement voirie (PK 0+000 à 2+500).\newline -- Contrôle qualité béton (affaissement, éprouvettes).\newline -- AutoCAD : reprise plans exécution voirie/assainissement.\newline -- Rédaction 30 rapports journaliers. & \textbf{Stage :} « J'ai fait un stage dans une boîte de BTP. J'ai appris beaucoup de choses. » (Pas de nom, dates, tâches précises, résultats). \\ \hline
\textbf{Compétences techniques} & \textbf{Logiciels :} AutoCAD 2D/3D (bon niveau), MS Project (notions), Suite Office (maîtrisée).\newline \textbf{Métier :} Lecture plans EXE/DAO, calculs béton/acier (BAEL/EC2 notions), topographie (station totale), normes sécurité chantier. & \textbf{Informatique :} Word, Excel, Internet.\newline \textbf{Autres :} « Travail manuel, courageux, honnête. » \\ \hline
\textbf{Langues} & Français : langue maternelle.\newline Anglais : niveau B1 (TOEIC 650 -- 2023). & Français : « bien ». Anglais : « scolaire ». \\ \hline
\textbf{Centres d'intérêt (pertinents)} & Trésorier association élèves ingénieurs (gestion budget 500 000 FCFA).\newline Volley-ball (esprit équipe, discipline). & « Lecture, musique, sortir entre amis. » \\ \hline
\textbf{Orthographe / Syntaxe} & Zéro faute. Phrases nominales, verbes d'action au passé composé / infinitif. & 7 fautes d'orthographe/grammaire repérées en 10 s. Phrases familières (« j'ai géré », « c'était cool »). \\ \hline
\textbf{Adéquation à l'offre} & Mots-clés de l'offre présents : AutoCAD, béton, rapports, sécurité, disponibilité. & Aucun mot-clé de l'offre. Compétences génériques. \\ \hline
\end{tabularx}
\end{center}
\legende{Tableau comparatif : CV d'Alain (retenu pour entretien) vs CV de Junior (rejeté à la lecture).}
\end{popfigure}

\begin{exoresolu}[Analyse guidée : pourquoi Alain passe et Junior échoue]
\textbf{Énoncé :} À partir du tableau comparatif, identifiez les 5 critères discriminants qui ont conduit le recruteur à retenir Alain et à rejeter Junior en moins de 40 secondes.

\textbf{Solution détaillée :}
\begin{enumerate}
    \item \textbf{Respect du format « une page » :} Alain tient en une page ; Junior en fait deux pour un jeune diplômé $\rightarrow$ signal d'incapacité à synthétiser.
    \item \textbf{Lisibilité visuelle (scannabilité) :} Alain utilise une hiérarchie visuelle claire (titres, puces, blanc) ; Junior présente un « pavé » illisible en diagonale.
    \item \textbf{Précision factuelle (vérifiabilité) :} Alain donne noms d'entreprises, dates, tâches chiffrées (30 rapports, PK 0+000 à 2+500) ; Junior reste vague (« une boîte », « beaucoup de choses »).
    \item \textbf{Adéquation mots-clés / offre :} Alain miroir l'offre (AutoCAD, béton, rapports, sécurité) ; Junior n'utilise aucun terme technique du métier.
    \item \textbf{Zéro faute vs fautes multiples :} 7 fautes en 30 secondes = incompétence perçue en communication écrite, compétence clé pour rédiger des rapports de chantier.
\end{enumerate}
\emph{Moralité :} Le CV n'est pas une biographie, c'est une \textbf{réponse ciblée} à une offre précise.
\end{exoresolu}

\courssub{Critères d'efficacité d'un CV : la check-list du recruteur}

\begin{methode}[Analyser / vérifier un CV avant envoi -- 10 points de contrôle]
\begin{enumerate}
    \item \textbf{Structure :} Les 5 rubriques obligatoires sont-elles présentes ? (Identité, Formation, Expériences, Compétences, Langues/Centres d'intérêt).
    \item \textbf{Ordre antéchronologique :} Le plus récent en premier dans chaque rubrique ?
    \item \textbf{Cohérence temporelle :} Pas de trous inexpliqués > 6 mois ? Dates alignées (mois/année) ?
    \item \textbf{Adéquation offre :} Les 3--5 mots-clés techniques de l'annonce apparaissent-ils dans le CV ?
    \item \textbf{Quantification :} Au moins 3 résultats chiffrés (budget géré, effectifs encadrés, \% d'avancement, nombre de plans, surface bétonnée, etc.) ?
    \item \textbf{Verbes d'action :} Chaque expérience commence par un verbe fort au passé composé ou à l'infinitif (réalisé, supervisé, calculé, rédigé, coordonné...) ?
    \item \textbf{Orthographe / Grammaire :} Relu à voix haute + correcteur + relecture par un tiers. Zéro faute tolérée.
    \item \textbf{Mise en page :} Une page, police unique (Calibri, Arial, Helvetica), tailles cohérentes, marges $\ge$ 1,5 cm, PDF (pas Word) pour l'envoi.
    \item \textbf{Honnêteté :} Chaque ligne est-elle vérifiable (diplôme, stage, certificat de travail, référence) ?
    \item \textbf{Identité numérique :} Email professionnel (prénom.nom@...), LinkedIn à jour, pas de photo de soirée.
\end{enumerate}
\end{methode}

\begin{attention}[Piège éliminatoire : le mensonge sur le CV]
\begin{itemize}
    \item \textbf{Mentir sur un diplôme, une mention, une expérience, un niveau de langue} est vérifiable en 1 appel téléphonique ou 1 email (établissement, ancien employeur, organisme de certification).
    \item Au Cameroun, les grands recruteurs (Fonction publique, banques, ONG, groupes industriels) font systématiquement la \emph{due diligence} (vérification des certificats de scolarité, attestations de stage, casier judiciaire).
    \item \textbf{Conséquence :} Élimination immédiate, inscription sur liste noire interne, atteinte à la réputation professionnelle durable.
    \item \textbf{Alternative éthique :} Valoriser ses \emph{vrais} atouts (projets scolaires, bénévolat, auto-formation, soft skills) avec des verbes d'action et des résultats concrets.
\end{itemize}
\end{attention}

\courssub{Le CV comme texte de communication : analyse énonciative}

\begin{propriete}[Triangle communicationnel appliqué au CV]
\begin{center}
\begin{tabularx}{\linewidth}{|l|X|X|}\hline
\rowcolor{popblueL}
\textbf{Composante} & \textbf{Dans le CV} & \textbf{Enjeu pour le candidat} \\ \hline
\textbf{Émetteur} & Candidat (locuteur) & Crédibilité, honnêteté, professionnalisme perçus à travers la forme et le fond. \\ \hline
\textbf{Récepteur} & Recruteur / DRH / Jury (destinataire) & Lecture rapide, critères objectifs, comparaison entre candidats, décision binaire (entretien / poubelle). \\ \hline
\textbf{Message} & Le CV (support + contenu) & Clarté, concision, pertinence, absence d'ambiguïté, adaptation au poste visé. \\ \hline
\textbf{Code} & Français standard (écrit formel) + codes visuels (mise en page) + codes métier (vocabulaire technique) & Registre soutenu, zéro faute, terminologie précise du secteur (BTP, informatique, secrétariat, etc.). \\ \hline
\textbf{Contact} & Email / plateforme / dépôt physique & Canal fiable, format PDF, nom de fichier professionnel (\texttt{Nom\_Prenom\_CV\_Poste.pdf}). \\ \hline
\textbf{Fonction} & \emph{Conative} (inciter à l'action : convocation) + \emph{Référentielle} (informations factuelles) & Déclencher l'entretien ; le CV ne « vend » pas, il « ouvre la porte ». \\ \hline
\end{tabularx}
\end{center}
\end{propriete}

\begin{savaistu}[Spécificité camerounaise : la règle de la page unique]
Au Cameroun, pour les jeunes diplômés (Bac, BT, BTS, Licence), \textbf{la très grande majorité des recruteurs publics et privés exigent un CV strictement tenu sur une seule page A4}. Cela vaut pour la Fonction publique (concours, recrutements directs), les banques (Afriland, BICEC, UBA), les entreprises du secteur pétrolier (SNPC, Perenco), les ONG internationales et les grands groupes (Bolloré, Castel, Guinness). Un CV de deux pages pour un primo-demandeur d'emploi est perçu comme un manque de synthèse et de maturité professionnelle.
\end{savaistu}

\courssub{Synthèse de la découverte}

\begin{aretenir}[Ce qu'il faut retenir de cette section]
\begin{itemize}
    \item Le CV est un \textbf{texte fonctionnel} à finalité professionnelle : informer, convaincre, obtenir un entretien.
    \item Il s'inscrit dans un \textbf{circuit de communication asymétrique} : candidat $\rightarrow$ CV $\rightarrow$ recruteur (30--40 s) $\rightarrow$ entretien.
    \item L'efficacité d'un CV repose sur \textbf{5 piliers} : clarté visuelle, concision (1 page), honnêteté, adéquation à l'offre, correction linguistique.
    \item L'analyse comparative de deux CV réels montre que \textbf{la structure, la précision factuelle, les mots-clés métier et l'absence de fautes} sont discriminants dès la première lecture.
    \item Toute fausse information est \textbf{vérifiable et éliminatoire} ; l'éthique professionnelle commence par le CV.
\end{itemize}
\end{aretenir}

\begin{exercice}[Application : audit express d'un CV]
Vous trouverez ci-dessous un extrait de CV d'un candidat au poste de \textit{Technicien de maintenance industrielle} (offre exigeant : lecture plans électriques, dépannage automatismes, rédaction rapports, anglais technique).

\begin{quote}
\textbf{MOI} \\
\textit{jean.pierre@hotmail.fr} -- 6XX XX XX XX \\
\textbf{Formation} \\
2022 : BAC F3 (élec) \\
2019--2022 : Lycée technique \\
\textbf{Expérience} \\
Stage 2022 : Usine SOSUCAM. J'ai réparé des machines. J'ai appris l'automatisme. \\
\textbf{Compétences} \\
Électricité, mécanique, informatique (Word, Excel), anglais lu écrit parlé. \\
\textbf{Loisirs} \\
Football, cinéma, amis.
\end{quote}

\textbf{Consignes :} En vous appuyant sur la check-list de la méthode (10 points), relevez \textbf{au moins 7 défauts majeurs} qui feront rejeter ce CV en 30 secondes. Pour chaque défaut, proposez une correction concrète.
\end{exercice}

\begin{corrige}[Exercice : audit express d'un CV]
\textbf{7 défauts majeurs et corrections :}
\begin{enumerate}
    \item \textbf{En-tête incomplet / email non pro :} « MOI », pas de nom/prénom, email hotmail. \textit{Correction :} « Jean Pierre NGONO -- 6XX XX XX XX -- jean.pierre.ngono@email.cm -- Douala, Bonapriso -- LinkedIn : linkedin.com/in/jpngono ».
    \item \textbf{Titre / accroche absent :} Pas de cible visible. \textit{Correction :} Ajouter en haut : « Technicien Maintenance Industrielle -- Bac F3 Électrotechnique -- Spécialité Automatismes -- Disponible ».
    \item \textbf{Formation : mention / établissement / notes manquantes.} \textit{Correction :} « 2022 : Baccalauréat F3 Électrotechnique, Lycée Technique de Douala-Bassa (Mention AB, 14,5/20 en Spécialité) ».
    \item \textbf{Expérience : vague, non chiffrée, pas de verbes d'action, pas de résultats.} \textit{Correction :} « Stage 2022 (3 mois) -- SOSUCAM, Nkoteng -- Maintenance préventive/curative ligne conditionnement (30 machines) -- Dépannage automates Siemens S7-300 (5 interventions/semaine) -- Lecture schémas électriques -- Rédaction 12 fiches d'intervention -- Respect consignes sécurité (0 accident) ».
    \item \textbf{Compétences : génériques, pas de niveau, pas de mots-clés de l'offre (automatismes, anglais technique).} \textit{Correction :} « Automatismes : Siemens S7-300/1200 (programmation STEP 7, notions TIA Portal) -- Électrotechnique : lecture schémas puissance/commande, câblage armoires, variateurs de vitesse -- Anglais technique : lecture datasheets/manuels (niveau B1 technique, TOEIC 600) -- Bureautique : Excel avancé (TCD, macros), Word, Outlook ».
    \item \textbf{Mise en page :} pas de structure visuelle, pas de puces, police unique non respectée. \textit{Correction :} Appliquer la charte une page, puces, titres gras, alignement dates à droite.
    \item \textbf{Centres d'intérêt :} génériques, non transférables. \textit{Correction :} « Capitaine équipe football entreprise (leadership, gestion conflits) -- Bénévolat : maintenance matériel informatique ONG locale (autonomie, sens service) ».
\end{enumerate}
\emph{Bilan :} Ce CV corrigé passe le test des 30 secondes ; l'original ne le passe pas.
\end{corrige}

\coursec{Les éléments constitutifs du CV}

Dans la section précédente, nous avons découvert que le CV est un \cle{texte fonctionnel} : il ne cherche ni à amuser ni à émouvoir, mais à obtenir un résultat concret — décrocher un entretien d'embauche. Or, pour remplir cette fonction, un texte fonctionnel obéit à une architecture précise. De même qu'une lettre administrative possède un en-tête, un objet et une formule de politesse, le CV possède des \cle{rubriques} obligatoires, ordonnées selon une logique que le recruteur attend. Cette section présente chacun de ces blocs constitutifs : l'en-tête, le titre, la formation, l'expérience professionnelle, les compétences, les centres d'intérêt et les références.

\courssub{Vue d'ensemble : l'architecture d'un CV}

Un recruteur consacre en moyenne moins d'une minute à la première lecture d'un CV. Il ne lit pas : il \emph{balaye} le document du regard, à la recherche de repères fixes. Si ces repères manquent ou sont mal placés, le CV est écarté. L'architecture du CV répond donc à une exigence de \cle{lisibilité} : chaque information doit se trouver là où le lecteur l'attend.

\begin{exemplebox}[Un CV en chiffres]
Un CV de jeune diplômé compte en moyenne \textbf{5 rubriques} (en-tête et titre, formation, expérience, compétences, centres d'intérêt) et tient sur \textbf{une seule page A4}. Au-delà d'une page, un candidat débutant donne l'impression de ne pas savoir sélectionner l'essentiel — ce qui est en soi une mauvaise compétence professionnelle.
\end{exemplebox}

\begin{popfigure}
\begin{center}
\begin{tikzpicture}[scale=1]
% Bloc 1 : en-tête
\draw[fill=popblueL,draw=popblue,thick,rounded corners=2pt] (0,8.2) rectangle (10,9.6);
\node[font=\bfseries] at (5,9.1) {1. EN-T\^ETE + TITRE};
\node[font=\small] at (5,8.55) {Nom, pr\'enoms, t\'el\'ephone, courriel, ville + poste vis\'e};
% Bloc 2 : formation
\draw[fill=popgreenL,draw=popgreen,thick,rounded corners=2pt] (0,6.2) rectangle (10,8.0);
\node[font=\bfseries] at (5,7.5) {2. FORMATION};
\node[font=\small] at (5,6.9) {Dipl\^omes du plus r\'ecent au plus ancien};
\node[font=\small] at (5,6.5) {\'etablissements, ann\'ees d'obtention};
% Bloc 3 : expérience
\draw[fill=poporangeL,draw=poporange,thick,rounded corners=2pt] (0,4.2) rectangle (10,6.0);
\node[font=\bfseries] at (5,5.5) {3. EXP\'ERIENCE PROFESSIONNELLE};
\node[font=\small] at (5,4.9) {Stages et emplois : dates, structure, t\^aches};
\node[font=\small] at (5,4.5) {verbe d'action + t\^ache + r\'esultat};
% Bloc 4 : compétences
\draw[fill=poppurpleL,draw=poppurple,thick,rounded corners=2pt] (0,2.4) rectangle (10,4.0);
\node[font=\bfseries] at (5,3.5) {4. COMP\'ETENCES};
\node[font=\small] at (5,2.9) {Techniques (informatique, conduite) et linguistiques};
% Bloc 5 : centres d'intérêt
\draw[fill=popgoldL,draw=popgold,thick,rounded corners=2pt] (0,0.8) rectangle (10,2.2);
\node[font=\bfseries] at (5,1.7) {5. CENTRES D'INT\'ER\^ET (+ R\'EF\'ERENCES)};
\node[font=\small] at (5,1.15) {Activit\'es valorisantes et v\'eridiques};
% Annotations latérales
\draw[->,thick,popdark] (10.4,8.9) -- (11.6,8.9) node[right,font=\small\itshape,text width=3.2cm] {qui je suis, ce que je cherche};
\draw[->,thick,popdark] (10.4,7.1) -- (11.6,7.1) node[right,font=\small\itshape,text width=3.2cm] {ce que j'ai appris};
\draw[->,thick,popdark] (10.4,5.1) -- (11.6,5.1) node[right,font=\small\itshape,text width=3.2cm] {ce que j'ai d\'ej\`a fait};
\draw[->,thick,popdark] (10.4,3.2) -- (11.6,3.2) node[right,font=\small\itshape,text width=3.2cm] {ce que je sais faire};
\draw[->,thick,popdark] (10.4,1.5) -- (11.6,1.5) node[right,font=\small\itshape,text width=3.2cm] {qui je suis hors du travail};
\end{tikzpicture}
\end{center}
\legende{Sch\'ema annot\'e d'un CV type en cinq blocs. L'ordre des rubriques n'est pas arbitraire : il va de l'identit\'e vers les preuves (formation, exp\'erience), puis vers les compl\'ements (comp\'etences, centres d'int\'er\^et).}
\end{popfigure}

\courssub{L'en-tête : se présenter et se rendre joignable}

L'\cle{en-tête} occupe le haut de la page. C'est la première chose que lit le recruteur ; elle doit lui permettre de répondre à deux questions : \emph{qui est ce candidat ?} et \emph{comment le contacter ?}

L'en-tête comprend obligatoirement :
\begin{itemize}
\item le \textbf{nom} (en majuscules, c'est un usage apprécié : MBALLA) et le ou les \textbf{prénoms} ;
\item le \textbf{numéro de téléphone} (un numéro sur lequel on peut réellement vous joindre) ;
\item l'\textbf{adresse électronique}, de préférence professionnelle dans sa forme ;
\item la \textbf{ville} de résidence (Douala, Bafoussam, Ngaoundéré...), qui renseigne le recruteur sur votre mobilité.
\end{itemize}

\begin{attention}[L'adresse électronique fantaisiste]
Erreur fréquente et éliminatoire : utiliser une adresse créée à l'adolescence du type \emph{beaugoss237@...} ou \emph{princesse\_du\_mfoundi@...}. Une telle adresse détruit immédiatement l'image de sérieux que le CV cherche à construire. La règle est simple : créer une adresse sobre de la forme \textbf{nom.prenom@...} (par exemple \emph{mballa.jean@...}), réservée si possible aux démarches professionnelles. Vérifier aussi que le répondeur du téléphone ne comporte pas de message plaisantin.
\end{attention}

\courssub{Le titre du CV : annoncer clairement le poste visé}

Sous l'en-tête, le \cle{titre du CV} indique en une ligne l'\textbf{intitulé du poste visé} ou le \textbf{diplôme} du candidat. Il fonctionne comme l'objet d'une lettre : il oriente immédiatement la lecture.

\begin{exemplebox}[Titres de CV efficaces]
\begin{itemize}
\item \emph{Technicien en froid et climatisation}
\item \emph{Secrétaire comptable — BTS Gestion des PME}
\item \emph{Électromécanicien, titulaire du Probatoire technique F2}
\end{itemize}
Un titre vague comme \emph{« À la recherche d'un emploi »} est à proscrire : il n'apporte aucune information et suggère que le candidat ne sait pas lui-même ce qu'il cherche.
\end{exemplebox}

\courssub{La formation : les diplômes du plus récent au plus ancien}

La rubrique \cle{Formation} présente les diplômes obtenus ou en préparation, avec l'établissement et l'année. Son principe d'organisation est l'\textbf{ordre antichronologique}.

\begin{definition}[Ordre antichronologique]
L'\cle{ordre antichronologique} est un mode de présentation qui range les événements \textbf{du plus récent au plus ancien} (préfixe \emph{anti-} : contraire ; \emph{chrono-} : temps). C'est l'inverse de l'ordre chronologique du récit. Le CV l'adopte parce que l'information la plus importante pour le recruteur est votre situation \emph{actuelle} : votre dernier diplôme, votre dernier emploi.
\end{definition}

\begin{exemplebox}[Rubrique « Formation » bien ordonnée]
\begin{center}
\begin{tabularx}{\linewidth}{|l|X|l|}\hline
\rowcolor{popblueL} \textbf{Ann\'ee} & \textbf{Dipl\^ome} & \textbf{\'Etablissement} \\ \hline
2025 & Probatoire technique F3 (\'Electronique) & Lyc\'ee technique de Douala-Bassa \\ \hline
2024 & BEPC & Coll\`ege Bilingue de Bonab\'eri \\ \hline
2022 & CAPI (Certificat d'Aptitude Professionnelle Industrielle) & CETIC de De\"ido \\ \hline
\end{tabularx}
\end{center}
On remarque que 2025 vient \emph{avant} 2022 : c'est l'ordre antichronologique.
\end{exemplebox}

\begin{popfigure}
\begin{center}
\begin{tikzpicture}[scale=1]
% Frise
\draw[->,very thick,popdark] (0,0) -- (13,0) node[right,font=\small\itshape] {temps};
% Graduations
\foreach \x/\a in {1/2019, 3.4/2021, 5.8/2022, 8.2/2024, 10.6/2025} {
  \draw[thick,popdark] (\x,-0.15) -- (\x,0.15);
  \node[below,font=\small] at (\x,-0.2) {\a};
}
% Ordre chronologique (récit)
\draw[->,thick,popmuted] (1,0.6) -- (10.6,0.6);
\node[above,font=\small\itshape,popmuted] at (5.8,0.6) {ordre chronologique (r\'ecit) : du plus ancien au plus r\'ecent};
% Ordre antichronologique (CV)
\draw[->,very thick,poporange] (10.6,-1.1) -- (1,-1.1);
\node[below,font=\small\bfseries,poporange] at (5.8,-1.2) {ordre antichronologique (CV) : du plus r\'ecent au plus ancien};
% Jalons
\node[circle,fill=popgreen,inner sep=1.5pt] at (1,0) {};
\node[circle,fill=popgreen,inner sep=1.5pt] at (3.4,0) {};
\node[circle,fill=popgreen,inner sep=1.5pt] at (5.8,0) {};
\node[circle,fill=popblue,inner sep=1.5pt] at (8.2,0) {};
\node[circle,fill=poporange,inner sep=1.5pt] at (10.6,0) {};
\node[above,font=\scriptsize,text width=1.7cm,align=center] at (1,0.1) {CAPI};
\node[above,font=\scriptsize,text width=1.7cm,align=center] at (5.8,0.1) {BEPC};
\node[above,font=\scriptsize,text width=1.7cm,align=center] at (8.2,0.1) {stage ENEO};
\node[above,font=\scriptsize,text width=1.9cm,align=center] at (10.6,0.1) {Probatoire F3};
\end{tikzpicture}
\end{center}
\legende{Frise chronologique : le r\'ecit lit la frise de gauche \`a droite ; le CV la lit de droite \`a gauche. La rubrique « Formation » commence donc par le Probatoire (2025) et remonte vers les dipl\^omes plus anciens.}
\end{popfigure}

\courssub{L'expérience professionnelle : prouver par les faits}

La rubrique \cle{Expérience professionnelle} recense les \textbf{stages} et les \textbf{emplois} occupés. Pour un jeune diplômé, les stages de fin de cycle, les petits boulots et même les activités familiales à caractère technique (atelier, boutique, ferme) comptent comme de l'expérience. Chaque expérience mentionne :
\begin{enumerate}
\item les \textbf{dates} (mois et année de début et de fin) ;
\item la \textbf{structure} d'accueil (entreprise, administration, atelier) ;
\item le \textbf{poste} occupé et les \textbf{tâches accomplies}.
\end{enumerate}

C'est dans la formulation des tâches que se joue la qualité du CV. Une tâche mal formulée est passive et vague ; une tâche bien formulée est active et mesurable.

\begin{methode}[Rédiger une ligne d'expérience : verbe d'action + tâche + résultat]
Pour chaque expérience, suivre trois étapes :
\begin{enumerate}
\item Choisir un \textbf{verbe d'action} au participe passé ou à l'infinitif : \emph{assurer, installer, réparer, gérer, accueillir, saisir, contrôler}.
\item Nommer précisément la \textbf{tâche} accomplie : \emph{la maintenance des installations frigorifiques}.
\item Ajouter si possible un \textbf{résultat chiffré} : \emph{sur 15 installations}.
\end{enumerate}
Formule complète : \emph{« Assuré la maintenance de 15 installations frigorifiques »}. Comparer avec la formulation faible : \emph{« J'étais chargé de la maintenance »}, qui ne dit ni ce qui a été fait ni avec quel résultat.
\end{methode}

\begin{exoresolu}[Transformer une expérience vague en ligne de CV efficace]
Énoncé : Aïcha a effectué un stage de deux mois dans un cybercafé de Garoua. Elle y accueillait les clients, saisissait des documents et a réduit les erreurs de facturation. Rédiger la ligne « expérience » de son CV.
\tcblower
Solution détaillée :
\begin{enumerate}
\item Dates et structure : \emph{Juillet--août 2025 — Cybercafé Le Sahel, Garoua}.
\item Verbe d'action + tâche : \emph{Accueilli la clientèle et saisi des documents administratifs}.
\item Résultat chiffré : \emph{réduit les erreurs de facturation de moitié}.
\end{enumerate}
Ligne finale : \emph{« Juillet--août 2025 — Stagiaire, Cybercafé Le Sahel (Garoua) : accueilli la clientèle, saisi des documents administratifs, réduit les erreurs de facturation de moitié. »}
\end{exoresolu}

\courssub{Les compétences : techniques et linguistiques}

La rubrique \cle{Compétences} répond à la question : \emph{que savez-vous faire concrètement ?} Elle se divise en deux familles.

\begin{itemize}
\item \textbf{Compétences techniques} : maîtrise de l'informatique (traitement de texte, tableur, messagerie), conduite d'engins ou de véhicules (permis B, C), utilisation de logiciels ou de machines propres à la spécialité (logiciel de comptabilité, oscilloscope, tour, soudeuse).
\item \textbf{Compétences linguistiques} : français et anglais (préciser le niveau : lu, écrit, parlé ; courant, scolaire), ainsi que les \textbf{langues nationales} (ewondo, fulfuldé, duala...), qui constituent un atout réel pour les postes au contact de la clientèle.
\end{itemize}

\begin{attention}[N'affirmer que ce que l'on peut prouver]
Le recruteur peut vérifier chaque compétence pendant l'entretien : il suffit qu'il passe à l'anglais ou qu'il demande une démonstration sur machine. Une compétence exagérée (\emph{« anglais courant »} alors qu'on peine à se présenter) décrédibilise l'ensemble du CV. Mieux vaut écrire honnêtement \emph{« anglais scolaire »}.
\end{attention}

\courssub{Les centres d'intérêt : montrer une personnalité}

La rubrique \cle{Centres d'intérêt} humanise le candidat. Elle doit mentionner des activités \textbf{valorisantes} et \textbf{véridiques} :
\begin{itemize}
\item le \textbf{sport} (football en club, athlétisme), qui suggère l'esprit d'équipe et l'endurance ;
\item l'\textbf{engagement associatif} (membre d'une coopérative, animateur de jeunesse), qui suggère le sens des responsabilités ;
\item la \textbf{lecture}, la musique, l'informatique créative, qui suggèrent la curiosité.
\end{itemize}
Deux pièges : les intérêts trop vagues (\emph{« cinéma, voyages »} n'engagent à rien) et les mensonges (le recruteur, amateur de lecture, peut demander quel livre vous avez aimé récemment).

\courssub{Les références : des témoins de votre sérieux}

Les \cle{références} sont des personnes (anciens maîtres de stage, enseignants, employeurs) qui acceptent de \textbf{témoigner du sérieux} du candidat auprès du recruteur. Cette rubrique est \textbf{facultative} mais appréciée, surtout pour un jeune diplômé dont l'expérience est courte. On indique le nom, la fonction et le téléphone de deux personnes au maximum — après leur avoir demandé l'autorisation. Une formule admise : \emph{« Références disponibles sur demande »}.

\courssub{Ce qu'il faut éviter}

Un CV n'est pas une fiche d'état civil. Certaines informations doivent en être exclues :
\begin{itemize}
\item les \textbf{informations discriminatoires inutiles} : religion, opinion politique, ethnie, situation matrimoniale détaillée, nombre d'enfants — aucune ne concerne l'aptitude au poste ;
\item les \textbf{détails personnels excessifs} : taille, poids, groupe sanguin, noms des parents ;
\item les \textbf{mensonges} sur les diplômes ou les dates, vérifiables et éliminatoires ;
\item les \textbf{fautes d'orthographe}, qui contredisent à elles seules toute la rubrique « compétences ».
\end{itemize}

\begin{savaistu}[La photo : obligatoire ou non ?]
Contrairement à une idée répandue, la \textbf{photo n'est pas obligatoire} sur un CV. Aucun texte ne l'impose, et beaucoup de recruteurs préfèrent les CV sans photo pour éviter toute discrimination. Si le candidat choisit d'en insérer une, elle doit être de type \textbf{photo d'identité} : récente, nette, sur fond clair, en tenue professionnelle — jamais un autoportrait de téléphone pris entre amis.
\end{savaistu}

\begin{aretenir}[L'essentiel de la section]
Un CV complet comprend : un \textbf{en-tête} (nom, prénoms, téléphone, courriel sobre, ville), un \textbf{titre} (poste visé ou diplôme), la \textbf{formation} et l'\textbf{expérience} en ordre antichronologique (du plus récent au plus ancien), les \textbf{compétences} techniques et linguistiques, les \textbf{centres d'intérêt} valorisants et, facultativement, des \textbf{références}. Chaque ligne d'expérience suit la formule \emph{verbe d'action + tâche + résultat}. Sont à bannir : adresse électronique fantaisiste, informations discriminatoires, détails personnels excessifs et tout mensonge.
\end{aretenir}

\begin{exercice}[À vous de jouer]
Voici les informations concernant Rodrigue, titulaire d'un CAP en menuiserie métallique : il a obtenu son CAP en 2025 au Lycée technique d'Edéa ; son BEPC en 2022 ; il a fait un stage de trois mois en 2024 à la SABC de Douala où il a assemblé des châssis en aluminium ; il parle français et bassa ; il joue au handball en club. Rédiger les rubriques « Formation », « Expérience » et « Compétences » de son CV, dans le bon ordre.
\end{exercice}

\begin{corrige}[Exercice]
\textbf{Formation} (ordre antichronologique) : 2025 — CAP Menuiserie métallique, Lycée technique d'Edéa ; 2022 — BEPC. \par
\textbf{Expérience} : 2024 (trois mois) — Stagiaire, SABC, Douala : assemblé des châssis en aluminium (verbe d'action + tâche). \par
\textbf{Compétences} : techniques — assemblage de châssis aluminium ; linguistiques — français (lu, écrit, parlé), bassa (courant). \par
Le handball en club relève des centres d'intérêt, non des compétences : c'est le piège de l'énoncé.
\end{corrige}

\coursec{Textes iconiques et mise en page du CV}

Dans la section précédente, nous avons identifié les \cle{éléments constitutifs} du CV : l'état civil, les coordonnées, la formation, l'expérience professionnelle, les compétences et les centres d'intérêt. Mais un contenu correct ne suffit pas : deux CV contenant exactement les mêmes informations peuvent produire des effets très différents sur un recruteur. La différence tient à la \cle{forme} : la photographie, les pictogrammes, la mise en page, la typographie. C'est ce que nous étudions maintenant : comment l'image et l'organisation visuelle du document participent au sens et à l'efficacité du CV, conformément aux attendus du programme (Réception des textes : Textes iconiques ; Production des textes : Le CV : présentation).

\courssub{Le texte iconique : quand l'image porte le sens}

\textbf{Observation préalable.} Prenez deux documents : une page de journal composée uniquement de phrases, et la page d'accueil d'un site d'annonces d'emploi. Dans le second cas, vous comprenez immédiatement où cliquer, où chercher, grâce aux icônes (une loupe pour « rechercher », une enveloppe pour « contact »), aux couleurs et à la disposition des blocs. Vous avez lu le document \emph{avant même de lire les phrases} : c'est l'image qui vous a guidé.

\begin{definition}[Texte iconique]
Un \cle{texte iconique} est un document dans lequel l'\cle{image} (photographie, pictogramme, graphique, schéma) associée à l'\cle{écrit} porte une partie du sens. L'image n'est pas un simple décor : elle informe, oriente la lecture, remplace parfois des mots. Le CV, la carte d'identité, le panneau de signalisation, la notice illustrée sont des textes iconiques.
\end{definition}

Dans un CV, trois éléments iconiques jouent un rôle essentiel :
\begin{itemize}
\item la \cle{photographie d'identité}, qui personnalise le document ;
\item les \cle{pictogrammes} (téléphone, enveloppe, localisation), qui signalent les coordonnées sans phrases inutiles ;
\item la \cle{mise en page} elle-même (blocs, alignements, filets, titres), qui structure visuellement l'information.
\end{itemize}

\begin{exemplebox}[L'effet des dix premières secondes]
Des études sur le comportement des recruteurs montrent qu'un CV bien structuré permet de repérer les rubriques essentielles (formation, expérience, coordonnées) en \textbf{moins de 10 secondes}. À l'inverse, un CV confus est souvent écarté avant même d'être lu. La forme n'est donc pas un luxe : c'est une condition pour que le contenu soit lu.
\end{exemplebox}

\courssub{La photographie d'identité}

La photo n'est pas obligatoire sur un CV, mais lorsqu'elle est présente, elle doit respecter des règles précises, car elle donne la première impression du candidat. Au Cameroun, comme dans de nombreux pays, une photo soignée est souvent attendue pour les postes de techniciens et d'employés de contact.

\begin{propriete}[Règles de la photographie d'identité sur un CV]
\begin{itemize}
\item \textbf{Cadrage} : le visage et le haut des épaules, de face ou légèrement de trois quarts ; le regard dirigé vers l'objectif.
\item \textbf{Tenue} : sobre et professionnelle (chemise, chemisier, tenue adaptée au métier visé).
\item \textbf{Arrière-plan} : neutre et uni (mur clair), sans objet ni personne visible.
\item \textbf{Expression} : naturelle, sérieuse mais avenante ; un léger sourire est recommandé.
\item \textbf{Qualité} : image nette, bien éclairée, récente ; jamais de photo de fête, de selfie ni de photo floue.
\end{itemize}
\end{propriete}

\begin{attention}[Piège fréquent]
Beaucoup de candidats insèrent une photo de mauvaise qualité (photo découpée sur un cliché de groupe, arrière-plan chargé, grimace). Une mauvaise photo dessert le candidat plus qu'une absence de photo. En cas de doute, mieux vaut ne pas en mettre.
\end{attention}

\courssub{Les pictogrammes : économie et lisibilité}

Un \cle{pictogramme} est un petit dessin symbolique qui représente une idée de manière immédiatement compréhensible. Sur un CV, les pictogrammes usuels sont :
\begin{itemize}
\item un \textbf{téléphone} devant le numéro (inutile d'écrire « Numéro de téléphone : ») ;
\item une \textbf{enveloppe} devant l'adresse électronique ;
\item un \textbf{repère de localisation} devant l'adresse ou la ville ;
\item éventuellement une \textbf{mallette} pour l'expérience, un \textbf{diplôme} pour la formation.
\end{itemize}

Leur fonction est double : l'\cle{économie} (ils remplacent des étiquettes verbales) et la \cle{lisibilité} (l'œil repère instantanément la nature de chaque information).

\begin{popfigure}
\begin{center}
\begin{tikzpicture}[scale=1]
% Téléphone
\draw[thick,popblue,fill=popblueL,rounded corners=2pt] (0.0,0.1) -- (0.15,0.55) -- (0.45,0.55) -- (0.6,0.1) -- (0.45,0.0) -- (0.15,0.0) -- cycle;
\node[below,font=\small] at (0.3,-0.1) {Téléphone};
% Enveloppe
\begin{scope}[xshift=2.2cm]
\draw[thick,poporange,fill=poporangeL] (0.0,0.0) rectangle (0.9,0.55);
\draw[thick,poporange] (0.0,0.55) -- (0.45,0.2) -- (0.9,0.55);
\node[below,font=\small] at (0.45,-0.1) {Courriel};
\end{scope}
% Localisation
\begin{scope}[xshift=4.4cm]
\draw[thick,popgreen,fill=popgreenL] (0.3,0.75) circle (0.22);
\draw[thick,popgreen,fill=popgreenL] (0.12,0.62) -- (0.3,0.0) -- (0.48,0.62) -- cycle;
\fill[popgreen] (0.3,0.75) circle (0.07);
\node[below,font=\small] at (0.3,-0.1) {Adresse};
\end{scope}
% Mallette
\begin{scope}[xshift=6.4cm]
\draw[thick,poppurple,fill=poppurpleL,rounded corners=1pt] (0.0,0.0) rectangle (0.9,0.5);
\draw[thick,poppurple] (0.3,0.5) -- (0.3,0.68) -- (0.6,0.68) -- (0.6,0.5);
\draw[thick,poppurple] (0.0,0.25) -- (0.9,0.25);
\node[below,font=\small] at (0.45,-0.1) {Expérience};
\end{scope}
% Diplôme
\begin{scope}[xshift=8.6cm]
\draw[thick,popgold,fill=popgoldL] (0.05,0.15) -- (0.45,0.6) -- (0.85,0.15) -- (0.45,-0.05) -- cycle;
\draw[thick,popgold] (0.45,-0.05) -- (0.45,-0.25);
\fill[popgold] (0.45,-0.28) circle (0.05);
\node[below,font=\small] at (0.45,-0.35) {Formation};
\end{scope}
\end{tikzpicture}
\end{center}
\legende{Planche de pictogrammes usuels d'un CV : téléphone, courriel, adresse, expérience professionnelle, formation. Chaque symbole remplace une étiquette verbale et accélère la lecture.}
\end{popfigure}

\courssub{La mise en page : aérer, aligner, hiérarchiser, uniformiser}

\textbf{Observation.} Photocopiez deux fois le même texte : une fois en un seul bloc compact, une fois découpé en paragraphes titrés et espacés. Montrez les deux versions à un camarade pendant cinq secondes, puis demandez-lui de retrouver une information précise. Il la trouvera toujours plus vite dans la version aérée. La mise en page est donc un \emph{outil de lecture}, pas une décoration.

\begin{methode}[Réussir la mise en page d'un CV]
\begin{enumerate}
\item \textbf{Aérer} : laisser des marges régulières (environ \SI{2}{cm}), espacer les rubriques, éviter les blocs de texte compacts. Le blanc guide l'œil.
\item \textbf{Aligner} : les titres, les dates et les contenus doivent suivre des axes verticaux constants (dates à gauche, descriptions à droite, par exemple). Rien ne doit « flotter » au hasard.
\item \textbf{Hiérarchiser} : le nom du candidat en grand en haut, les titres de rubriques en gras et d'une taille intermédiaire, le contenu en taille normale. L'œil doit pouvoir balayer le document du plus important au plus détaillé.
\item \textbf{Uniformiser} : une seule police, un seul style de titre, un seul format de date (ex. « 09/2024 » partout), des puces identiques dans tout le document.
\end{enumerate}
\end{methode}

\begin{propriete}[Choix typographique]
On utilise \textbf{une seule police} lisible et sobre (par exemple Arial ou Calibri), en corps \textbf{10 à 12} pour le texte courant, un peu plus grand pour les titres. Le \textbf{gras} est réservé aux titres et aux éléments à mettre en évidence ; l'\emph{italique} peut signaler les intitulés de postes ou d'établissements. Tout effet (gras, italique, couleur) s'emploie \textbf{avec parcimonie}.
\end{propriete}

\begin{attention}[Erreur fréquente]
Multiplier les couleurs et les polices nuit gravement à la lisibilité et donne une impression de désordre. Un CV professionnel se limite à \textbf{une police} et à \textbf{une ou deux couleurs} au maximum (par exemple le noir et un bleu sobre pour les titres). De même, on évite les tableaux à bordures visibles et les cadres multiples.
\end{attention}

\begin{popfigure}
\begin{center}
\begin{tikzpicture}[scale=0.9]
% Gabarit désordonné
\node[font=\bfseries] at (2.2,4.6) {Version désordonnée};
\draw[thick,popmuted] (0.0,0.0) rectangle (4.4,4.2);
\fill[popmuted!40] (0.3,3.6) rectangle (2.6,4.0);
\fill[popmuted!40] (0.9,2.7) rectangle (4.1,3.4);
\fill[popmuted!40] (0.2,1.9) rectangle (3.4,2.5);
\fill[popmuted!40] (1.2,0.9) rectangle (4.2,1.7);
\fill[popmuted!40] (0.4,0.2) rectangle (2.2,0.7);
\node[font=\small,popmuted] at (2.2,-0.4) {blocs de tailles et d'axes différents};
% Gabarit aligné
\begin{scope}[xshift=6cm]
\node[font=\bfseries] at (2.2,4.6) {Version aérée et alignée};
\draw[thick,popblue] (0.0,0.0) rectangle (4.4,4.2);
\fill[popblueL] (0.3,3.55) rectangle (4.1,4.0);
\draw[popblue,thick] (0.3,3.35) -- (4.1,3.35);
\fill[popblue!25] (0.3,2.45) rectangle (1.3,3.15);
\fill[popblue!10] (1.5,2.45) rectangle (4.1,3.15);
\fill[popblue!25] (0.3,1.55) rectangle (1.3,2.25);
\fill[popblue!10] (1.5,1.55) rectangle (4.1,2.25);
\fill[popblue!25] (0.3,0.65) rectangle (1.3,1.35);
\fill[popblue!10] (1.5,0.65) rectangle (4.1,1.35);
\draw[popblue,thick] (0.3,0.45) -- (4.1,0.45);
\fill[popblueL] (0.3,0.05) rectangle (4.1,0.35);
\node[font=\small,popblue] at (2.2,-0.4) {mêmes marges, dates à gauche, contenu à droite};
\end{scope}
\end{tikzpicture}
\end{center}
\legende{Comparaison de deux gabarits de mise en page pour un même contenu. À gauche : blocs déséquilibrés, non alignés, espacements irréguliers. À droite : marges constantes, colonne des dates alignée, rubriques séparées par des filets et des espaces réguliers.}
\end{popfigure}

\courssub{Lecture de documents : analyser deux mises en page d'un même contenu}

\begin{experience}[Analyse comparative guidée]
\textbf{Matériel} : deux versions d'un même CV (mêmes rubriques, mêmes informations) : la version A en un seul bloc, sans titres visibles, avec trois polices et des dates dispersées ; la version B aérée, alignée, avec une seule police et des rubriques titrées.

\textbf{Protocole} :
\begin{enumerate}
\item Chronométrez le temps nécessaire pour retrouver dans chaque version : le diplôme le plus récent, le numéro de téléphone, la dernière expérience professionnelle.
\item Notez pour chaque version les éléments iconiques présents (photo, pictogrammes, filets, puces).
\item Relevez les défauts de la version A en les classant selon les quatre critères : aération, alignement, hiérarchie, uniformité.
\end{enumerate}

\textbf{Observations attendues} : les informations sont repérées beaucoup plus vite dans la version B ; la version A fatigue l'œil et donne une impression négative du candidat, alors même que son contenu est identique.

\textbf{Interprétation} : la mise en page agit comme un \emph{paratexte visuel} : elle oriente la lecture et construit l'image du candidat avant toute lecture du contenu.
\end{experience}

\begin{exoresolu}[Diagnostiquer une mise en page défaillante]
\textbf{Énoncé.} Voici la description du CV de Marlyse : « Le document tient sur une page. Le nom est écrit en taille 10 comme le reste. On compte quatre polices différentes, des titres en rouge, vert et bleu, les dates tantôt à gauche, tantôt à droite, et aucune marge en bas de page. La photo est un selfie pris dans un salon. » Relevez cinq défauts et proposez une correction pour chacun.
\tcblower
\textbf{Solution détaillée.}
\begin{enumerate}
\item \emph{Nom en taille 10} : la hiérarchie est absente. Correction : écrire le nom en taille 16 à 20, en gras, en haut du document.
\item \emph{Quatre polices} : violation de l'uniformité. Correction : une seule police (Calibri ou Arial) pour tout le document.
\item \emph{Titres en trois couleurs} : surcharge visuelle. Correction : une seule couleur sobre (ou le noir) pour tous les titres.
\item \emph{Dates tantôt à gauche, tantôt à droite} : alignement incohérent. Correction : aligner toutes les dates sur le même axe vertical, idéalement à gauche.
\item \emph{Pas de marge en bas} : défaut d'aération. Correction : marges régulières d'environ \SI{2}{cm} sur les quatre côtés.
\item (Bonus) \emph{Selfie en photo} : photo non professionnelle. Correction : photo d'identité sur fond neutre, ou pas de photo du tout.
\end{enumerate}
\end{exoresolu}

\courssub{TP guidé : refonte de la mise en page d'un CV défaillant}

\begin{experience}[Travaux pratiques : refondre un CV]
\textbf{Document fourni} : le CV « défaillant » de M. Etoundi, candidat à un poste de technicien de maintenance (texte compact, trois polices, titres non hiérarchisés, coordonnées noyées dans un paragraphe, dates désordonnées, photo inadaptée).

\textbf{Tâche} : produire une maquette de refonte (à main levée ou à l'ordinateur) en appliquant la méthode.

\textbf{Étapes} :
\begin{enumerate}
\item \textbf{Repérage} : soulignez au crayon les rubriques présentes (état civil, formation, expérience, compétences, centres d'intérêt) et entourez les défauts de forme.
\item \textbf{Plan de page} : dessinez le gabarit cible : nom en haut au centre ou à gauche en grand ; coordonnées avec pictogrammes sous le nom ; deux colonnes (dates à gauche, contenu à droite) ; rubriques séparées par des espaces et des filets discrets.
\item \textbf{Refonte} : reportez chaque information dans le gabarit, en uniformisant le format des dates (ex. « 2023 – 2024 ») et les puces.
\item \textbf{Vérification} : contrôlez les quatre critères (aération, alignement, hiérarchie, uniformité) et faites le test des dix secondes avec un camarade : doit-il chercher longtemps votre diplôme ?
\end{enumerate}

\textbf{Critères de réussite} : une seule police ; marges régulières ; titres de rubriques en gras et de même taille ; coordonnées repérables en moins de trois secondes ; photo correcte ou absente.
\end{experience}

\begin{savaistu}[Les logiciels de tri de CV (ATS)]
Dans les grandes entreprises (SONARA, CDC, banques, administrations), les CV sont d'abord lus par des logiciels appelés \cle{ATS} (\emph{Applicant Tracking Systems}), qui extraient automatiquement les informations (nom, diplômes, compétences) avant même qu'un recruteur humain ne voie le document. Or ces logiciels lisent mal les \textbf{tableaux complexes}, les \textbf{colonnes multiples}, les \textbf{en-têtes et pieds de page} et les \textbf{images contenant du texte}. Un CV trop sophistiqué graphiquement peut donc être rejeté par la machine malgré un excellent contenu. La sobriété de la mise en page est aussi une stratégie numérique.
\end{savaistu}

\begin{aretenir}[L'essentiel de la section]
\begin{itemize}
\item Un \cle{texte iconique} associe image et écrit pour produire du sens ; le CV en est un exemple.
\item La photo d'identité : cadrage du visage, tenue sobre, fond neutre, expression naturelle — ou pas de photo.
\item Les pictogrammes (téléphone, enveloppe, localisation) assurent économie et lisibilité.
\item La mise en page obéit à quatre principes : \textbf{aérer, aligner, hiérarchiser, uniformiser}.
\item Typographie : une seule police lisible (Arial, Calibri), corps 10 à 12, gras et italique avec parcimonie, une ou deux couleurs maximum.
\item Un recruteur repère les rubriques d'un CV bien structuré en moins de dix secondes : la forme conditionne la lecture du contenu.
\end{itemize}
\end{aretenir}

\begin{exercice}[À faire en classe ou à la maison]
1. Définissez « texte iconique » et citez trois exemples de textes iconiques autres que le CV.

2. Relevez dans un CV réel (ou dans un modèle fourni par l'enseignant) tous les éléments iconiques présents et précisez la fonction de chacun.

3. Rédigez un court paragraphe (8 à 10 lignes) expliquant pourquoi la mise en page d'un CV peut être décisive lors d'une recherche d'emploi, en vous appuyant sur les quatre critères de la méthode.
\end{exercice}

\begin{corrige}[Exercice]
1. Un texte iconique est un document où l'image, associée à l'écrit, porte une partie du sens. Exemples : le panneau de signalisation, la carte d'identité, la notice de montage illustrée, l'affiche publicitaire, la une d'un journal.

2. Éléments attendus : la photo (personnaliser, donner une première impression) ; les pictogrammes des coordonnées (économie, repérage rapide) ; les filets et puces (séparer et ordonner les rubriques) ; le gras des titres (hiérarchiser) ; les blocs et marges (aérer, guider l'œil).

3. Le paragraphe doit mobiliser les quatre critères : un CV aéré et aligné se lit vite ; une hiérarchie claire (nom en grand, titres en gras) permet au recruteur de repérer les rubriques en quelques secondes ; l'uniformité (une police, un format de dates) donne une impression de sérieux et de rigueur, qualités attendues d'un technicien. À contenu égal, la forme peut donc décider de la convocation à un entretien.
\end{corrige}

\coursec{Langue : les valeurs du subjonctif et du conditionnel}

Dans la section précédente, nous avons appris à soigner la \cle{mise en page} du CV : photo, titres, couleurs, lisibilité. Mais un beau document ne suffit pas. Lors de l'entretien d'embauche et dans la lettre de motivation qui accompagne le CV, le recruteur juge aussi la \emph{langue} du candidat. Or deux modes grammaticaux reviennent sans cesse dans le discours professionnel : le \cle{subjonctif} («~Je souhaite que vous reteniez ma candidature~») et le \cle{conditionnel} («~Je serais honoré de rejoindre votre entreprise~»). Mal les employer, c'est risquer de paraître négligent ; bien les employer, c'est montrer sa maîtrise et sa politesse. Cette section nous donne les outils pour les former et les employer correctement.

\courssub{Le mode grammatical : l'attitude du locuteur}

\begin{definition}[Le mode grammatical]
Le \cle{mode} est une forme du verbe qui exprime l'\textbf{attitude du locuteur} face à son énoncé : il indique si le fait est présenté comme réel (\cle{indicatif}), comme un ordre (\cle{impératif}), comme un souhait, une volonté ou un doute (\cle{subjonctif}), ou comme une hypothèse, une éventualité (\cle{conditionnel}).
\end{definition}

Pour bien comprendre, observons quatre phrases qui parlent du même événement — votre arrivée à l'entretien — mais avec quatre attitudes différentes :

\begin{itemize}
\item \emph{Tu arrives à l'heure.} (indicatif : le fait est constaté, présenté comme réel) ;
\item \emph{Arrive à l'heure !} (impératif : le fait est ordonné) ;
\item \emph{Je veux que tu arrives à l'heure.} (subjonctif : le fait est souhaité, pas encore réalisé) ;
\item \emph{Tu arriverais à l'heure si tu partais tôt.} (conditionnel : le fait dépend d'une condition).
\end{itemize}

Le \emph{procès} (l'action d'arriver) ne change pas ; c'est le \emph{regard} du locuteur sur ce procès qui change. Le mode est donc une sorte de «~filtre~» posé sur la réalité.

\begin{popfigure}
\begin{tikzpicture}[
  grow=down,
  level distance=1.6cm,
  level 1/.style={sibling distance=4.6cm},
  level 2/.style={sibling distance=2.2cm},
  every node/.style={draw, rounded corners=2pt, align=center, font=\small, inner sep=3pt},
  edge from parent/.style={draw=popmuted, -{Stealth[length=2mm]}}
]
\node[fill=popblue, text=white, font=\small\bfseries] {Les modes du verbe}
  child { node[fill=popblueL, align=center]{Indicatif\\ \emph{fait réel}}
    child { node[fill=white, align=center]{présent, imparfait,\\ passé composé\ldots} } }
  child { node[fill=poporangeL, align=center]{Impératif\\ \emph{ordre, conseil}}
    child { node[fill=white, align=center]{présent\\ (sans sujet)} } }
  child { node[fill=popgreenL, align=center]{Subjonctif\\ \emph{volonté, doute,}\\\emph{ sentiment}}
    child { node[fill=white] {présent} } }
  child { node[fill=poppinkL, align=center]{Conditionnel\\ \emph{hypothèse,}\\\emph{ politesse}}
    child { node[fill=white] {présent} } };
\end{tikzpicture}
\legende{Arbre des modes personnels du verbe français et de leurs valeurs principales. L'indicatif présente le fait comme certain ; l'impératif comme un ordre ; le subjonctif comme un fait souhaité, douté ou ressenti ; le conditionnel comme un fait soumis à condition ou atténué par politesse. Au programme de Terminale technique, on travaille le subjonctif et le conditionnel au présent.}
\end{popfigure}

\begin{aretenir}[L'essentiel sur les modes]
Le mode n'indique pas \emph{quand} l'action se passe (c'est le rôle du temps), mais \emph{comment} le locuteur la considère : réelle, ordonnée, souhaitée ou hypothétique. Un même temps (le présent, par exemple) existe dans plusieurs modes.
\end{aretenir}

\courssub{Le subjonctif présent : formation}

\textbf{Le principe de formation.} Pour former le subjonctif présent de la plupart des verbes, on prend le \cle{radical de la 3\textsuperscript{e} personne du pluriel} du présent de l'indicatif, auquel on ajoute les terminaisons \textbf{-e, -es, -e, -ions, -iez, -ent}.

Prenons le verbe \emph{finir} : à l'indicatif présent, \emph{ils finissent}. On retire \emph{-ent} : il reste le radical \emph{finiss-}. On obtient : que je finiss\emph{e}, que tu finiss\emph{es}, qu'il finiss\emph{e}, que nous finiss\emph{ions}, que vous finiss\emph{iez}, qu'ils finiss\emph{ent}.

\begin{exemplebox}[Vérification avec «~vendre~»]
Indicatif présent : \emph{ils vendent} $\rightarrow$ radical \emph{vend-}. Subjonctif présent : que je \textbf{vende}, que tu \textbf{vendes}, qu'il \textbf{vende}, que nous \textbf{vendions}, que vous \textbf{vendiez}, qu'ils \textbf{vendent}.
\end{exemplebox}

\textbf{Les subjonctifs irréguliers usuels.} Quelques verbes très fréquents échappent à cette règle et doivent être mémorisés, car ils reviennent constamment dans les écrits professionnels.

\begin{center}
\begin{tabularx}{\linewidth}{|l|X|X|}
\hline
\rowcolor{popblueL} \textbf{Infinitif} & \textbf{Subjonctif présent (formes clés)} & \textbf{Exemple professionnel} \\
\hline
être & que je \textbf{sois}, que nous \textbf{soyons} & «~Je veux que mon dossier \textbf{soit} complet.~» \\
\hline
avoir & que j'\textbf{aie}, que nous \textbf{ayons} & «~Il faut que j'\textbf{aie} mon diplôme en main.~» \\
\hline
aller & que j'\textbf{aille}, que nous \textbf{allions} & «~Il faut que j'\textbf{aille} à cet entretien.~» \\
\hline
faire & que je \textbf{fasse} & «~L'entreprise demande que je \textbf{fasse} un essai.~» \\
\hline
pouvoir & que je \textbf{puisse} & «~Je cherche un poste où je \textbf{puisse} progresser.~» \\
\hline
savoir & que je \textbf{sache} & «~Il est important que je \textbf{sache} utiliser ce logiciel.~» \\
\hline
vouloir & que je \textbf{veuille}, que nous \textbf{voulions} & «~Nous souhaitons que vous \textbf{vouliez} bien nous répondre.~» \\
\hline
falloir & qu'il \textbf{faille} & «~Il le faut qu'il \textbf{faille} peu de temps pour répondre.~» (rare au présent, fréquent à l'imparfait : \emph{fallût}) \\
\hline
\end{tabularx}
\end{center}

\begin{exemplebox}[Un chiffre à retenir]
Dans les lettres de motivation et les échanges professionnels, \textbf{8 verbes irréguliers} (être, avoir, aller, faire, pouvoir, savoir, vouloir, falloir) concentrent l'essentiel des emplois du subjonctif. Celui qui maîtrise ces huit formes élimine la quasi-totalité des risques d'erreur.
\end{exemplebox}

\courssub{Les emplois du subjonctif}

Le subjonctif apparaît presque toujours dans une \cle{subordonnée} introduite par \emph{que}. On l'emploie :

\begin{enumerate}
\item après les verbes de \textbf{volonté} et de \textbf{souhait} : \emph{vouloir, désirer, souhaiter, exiger, demander que}\ldots\ Exemple : «~Je désire que vous examiniez ma candidature.~» ;
\item après les verbes d'\textbf{obligation} et de \textbf{nécessité} : \emph{il faut que, il est nécessaire que, il est essentiel que}\ldots\ Exemple : «~Il faut que je sois ponctuel.~» ;
\item après les verbes de \textbf{sentiment} : \emph{craindre que, avoir peur que, regretter que, être content que, être surpris que}\ldots\ Exemple : «~Je crains que mon expérience ne soit insuffisante.~» ;
\item après certaines \textbf{conjonctions} : \emph{afin que, pour que, bien que, quoique, avant que, à condition que, pourvu que}\ldots\ Exemple : «~Je joins mon CV afin que vous disposiez de toutes les informations.~»
\end{enumerate}

\begin{methode}[Choisir entre subjonctif et indicatif]
Pour décider du mode d'une subordonnée, procéder en trois étapes :
\begin{enumerate}
\item \textbf{Repérer} le verbe ou la conjonction qui introduit la subordonnée (\emph{je veux que}, \emph{il faut que}, \emph{bien que}, \emph{je constate que}\ldots) ;
\item \textbf{Classer} cet introducteur : exprime-t-il la volonté, le souhait, l'obligation, le sentiment, le doute ? $\rightarrow$ subjonctif. Exprime-t-il un fait constaté, une certitude (\emph{je sais que, je constate que, il est vrai que, il est certain que}) ? $\rightarrow$ indicatif ;
\item \textbf{Conjuguer} en conséquence et vérifier la forme (attention aux irréguliers : \emph{sois, aie, fasse, puisse, sache, veuille, aille, faille}).
\end{enumerate}
Exemple d'application : «~Je sais que vous \emph{êtes} (indicatif, certitude) très occupé, mais je souhaite que vous \emph{receviez} (subjonctif, souhait) ma demande.~»
\end{methode}

\courssub{Le conditionnel présent : formation}

Le conditionnel présent se forme sur le \cle{radical du futur simple}, auquel on ajoute les \cle{terminaisons de l'imparfait} : \textbf{-ais, -ais, -ait, -ions, -iez, -aient}.

\begin{itemize}
\item \emph{parler} : futur \emph{je parlerai} $\rightarrow$ conditionnel \emph{je parlerais, tu parlerais, il parlerait, nous parlerions, vous parleriez, ils parleraient} ;
\item \emph{être} : radical de futur irrégulier \emph{ser-} $\rightarrow$ \emph{je serais} ;
\item \emph{avoir} : radical \emph{aur-} $\rightarrow$ \emph{j'aurais} ;
\item \emph{faire} : radical \emph{fer-} $\rightarrow$ \emph{je ferais} ;
\item \emph{aller} : radical \emph{ir-} $\rightarrow$ \emph{j'irais} ; \emph{pouvoir} : \emph{je pourrais} ; \emph{vouloir} : \emph{je voudrais} ; \emph{voir} : \emph{je verrais} ; \emph{envoyer} : \emph{j'enverrais}.
\end{itemize}

\begin{popfigure}
\begin{tikzpicture}[
  every node/.style={font=\small},
  pers/.style={draw=popmuted, fill=popcream, rounded corners=2pt, minimum width=1.15cm, minimum height=0.62cm, align=center},
  head/.style={font=\small\bfseries, text=popdark}
]
% En-têtes de colonnes
\node[head] at (-2.3,2.6) {Personne} ;
\node[head] at (0.4,2.6) {être} ;
\node[head] at (3.1,2.6) {avoir} ;
\node[head] at (5.8,2.6) {faire} ;
% Lignes subjonctif (haut) et conditionnel (bas) séparées
\node[head, text=popgreen] at (-4.6,1.3) {Subjonctif} ;
\node[head, text=popgreen] at (-4.6,0.9) {présent} ;
\node[head, text=poppink] at (-4.6,-1.7) {Conditionnel} ;
\node[head, text=poppink] at (-4.6,-2.1) {présent} ;
% Personnes
\foreach \y/\p in {1.9/{je}, 1.3/{tu}, 0.7/{il/elle}, 0.1/{nous}, -0.5/{vous}, -1.1/{ils}} {
  \node at (-2.3,\y) {\p} ;
}
% Subjonctif : être
\node[pers] at (0.4,1.9) {sois} ;
\node[pers] at (0.4,1.3) {sois} ;
\node[pers] at (0.4,0.7) {soit} ;
\node[pers] at (0.4,0.1) {soyons} ;
\node[pers] at (0.4,-0.5) {soyez} ;
\node[pers] at (0.4,-1.1) {soient} ;
% Subjonctif : avoir
\node[pers] at (3.1,1.9) {aie} ;
\node[pers] at (3.1,1.3) {aies} ;
\node[pers] at (3.1,0.7) {ait} ;
\node[pers] at (3.1,0.1) {ayons} ;
\node[pers] at (3.1,-0.5) {ayez} ;
\node[pers] at (3.1,-1.1) {aient} ;
% Subjonctif : faire
\node[pers] at (5.8,1.9) {fasse} ;
\node[pers] at (5.8,1.3) {fasses} ;
\node[pers] at (5.8,0.7) {fasse} ;
\node[pers] at (5.8,0.1) {fassions} ;
\node[pers] at (5.8,-0.5) {fassiez} ;
\node[pers] at (5.8,-1.1) {fassent} ;
% Conditionnel
\foreach \y/\p in {-1.5/{je}, -2.1/{tu}, -2.7/{il/elle}, -3.3/{nous}, -3.9/{vous}, -4.5/{ils}} {
  \node at (-2.3,\y) {\p} ;
}
\node[pers, fill=poppinkL] at (0.4,-1.5) {serais} ;
\node[pers, fill=poppinkL] at (0.4,-2.1) {serais} ;
\node[pers, fill=poppinkL] at (0.4,-2.7) {serait} ;
\node[pers, fill=poppinkL] at (0.4,-3.3) {serions} ;
\node[pers, fill=poppinkL] at (0.4,-3.9) {seriez} ;
\node[pers, fill=poppinkL] at (0.4,-4.5) {seraient} ;
\node[pers, fill=poppinkL] at (3.1,-1.5) {aurais} ;
\node[pers, fill=poppinkL] at (3.1,-2.1) {aurais} ;
\node[pers, fill=poppinkL] at (3.1,-2.7) {aurait} ;
\node[pers, fill=poppinkL] at (3.1,-3.3) {aurions} ;
\node[pers, fill=poppinkL] at (3.1,-3.9) {auriez} ;
\node[pers, fill=poppinkL] at (3.1,-4.5) {auraient} ;
\node[pers, fill=poppinkL] at (5.8,-1.5) {ferais} ;
\node[pers, fill=poppinkL] at (5.8,-2.1) {ferais} ;
\node[pers, fill=poppinkL] at (5.8,-2.7) {ferait} ;
\node[pers, fill=poppinkL] at (5.8,-3.3) {ferions} ;
\node[pers, fill=poppinkL] at (5.8,-3.9) {feriez} ;
\node[pers, fill=poppinkL] at (5.8,-4.5) {feraient} ;
% Cadres de séparation
\draw[popgreen, thick, rounded corners=4pt] (-3.1,2.35) rectangle (7.0,-1.45) ;
\draw[poppink, thick, rounded corners=4pt] (-3.1,-1.45) rectangle (7.0,-4.85) ;
\end{tikzpicture}
\legende{Tableau de conjugaison des verbes \emph{être}, \emph{avoir} et \emph{faire} au subjonctif présent (cadre vert) et au conditionnel présent (cadre rose). Au subjonctif, les radicaux sont irréguliers ; au conditionnel, on retrouve le radical du futur (\emph{ser-}, \emph{aur-}, \emph{fer-}) suivi des terminaisons de l'imparfait.}
\end{popfigure}

\courssub{Les valeurs du conditionnel}

Le conditionnel présent possède quatre valeurs principales, toutes utiles dans le contexte professionnel.

\begin{enumerate}
\item \textbf{L'hypothèse} : le fait dépend d'une condition. «~Si j'avais plus d'expérience, je \emph{postulerais} à ce poste.~» C'est la valeur la plus fréquente.
\item \textbf{Le souhait} : «~Je \emph{voudrais} travailler dans votre entreprise.~» Le conditionnel rend le désir moins brutal qu'à l'indicatif («~je veux~»).
\item \textbf{La politesse atténuée} : «~Je \emph{souhaiterais} vous rencontrer. Pourriez-vous me recevoir ?~» Le conditionnel adoucit la demande et marque le respect.
\item \textbf{L'information non confirmée} (fréquent dans la presse) : «~L'entreprise \emph{recruterait} cinquante techniciens.~» Le journaliste signale qu'il ne garantit pas le fait.
\end{enumerate}

\begin{propriete}[La concordance avec «~si~»]
Dans une phrase exprimant une condition (hypothèse non réalisée au présent ou avenir), la subordonnée introduite par \emph{si} se met à l'\textbf{imparfait de l'indicatif}, et la proposition principale au \textbf{conditionnel présent} :
\[ \text{Si} + \text{imparfait} \quad \longrightarrow \quad \text{conditionnel présent} \]
Exemple : «~Si j'\emph{avais} (imparfait) un entretien demain, je \emph{préparerais} (conditionnel) mes réponses ce soir.~» Cette règle se mémorise telle quelle ; aucune démonstration n'est exigible.
\end{propriete}

\begin{attention}[L'erreur à ne jamais commettre]
Après \emph{si} exprimant une condition, on n'emploie \textbf{jamais} le conditionnel. Dire «~si j'\emph{aurais} su, je serais venu~» est une faute grave ; la forme correcte est «~si j'\emph{avais} su, je serais venu~». De même : «~si je pourrais~» est incorrect ; on dit «~si je \emph{pouvais}~». Astuce de vérification : le conditionnel ne se trouve que dans la principale, jamais dans la subordonnée en \emph{si}.
\end{attention}

\courssub{Application au contexte professionnel}

Pourquoi consacrer une leçon entière à ces deux modes ? Parce qu'ils structurent le discours de la recherche d'emploi. Observons quelques formules types :

\begin{itemize}
\item Lettre de motivation : «~Je vous serais reconnaissant de bien vouloir examiner ma candidature.~» (conditionnel de politesse) ;
\item Lettre de motivation : «~Je souhaite que vous reteniez mon profil pour ce poste.~» (subjonctif après un verbe de souhait) ;
\item Entretien : «~Je serais honoré de faire partie de votre équipe.~» (conditionnel de politesse et de souhait) ;
\item Entretien : «~Si je devais décrire ma principale qualité, je dirais la rigueur.~» (si + imparfait $\rightarrow$ conditionnel) ;
\item Entretien : «~Bien que je sois jeune diplômé, j'ai déjà réalisé plusieurs stages.~» (subjonctif après \emph{bien que}).
\end{itemize}

\begin{savaistu}[Le conditionnel, passeport de la politesse]
Dans la culture professionnelle francophone, le conditionnel de politesse est une marque de respect très attendue lors d'un entretien d'embauche. Dire «~Je \emph{voudrais} vous poser une question~» plutôt que «~Je \emph{veux} vous poser une question~» peut modifier l'impression laissée au recruteur. Les spécialistes du recrutement le confirment : au-delà des compétences techniques, c'est souvent la qualité de l'expression — dont l'usage correct des modes — qui départage deux candidats de niveau proche.
\end{savaistu}

\begin{exoresolu}[Corriger et compléter des phrases d'entretien]
\textbf{Énoncé.} Pour chaque phrase, indique le mode attendu et corrige si nécessaire.
\begin{enumerate}
\item «~Si j'aurais le poste, je donnerais le meilleur de moi-même.~»
\item «~Je veux que vous \emph{savez} (savoir) que je suis motivé.~»
\item «~Il faut que je \emph{sois} (être) ponctuel le jour de l'entretien.~»
\item «~Je \emph{souhaite} (souhaiter, politesse) vous rencontrer prochainement.~»
\end{enumerate}
\tcblower
\textbf{Solution détaillée.}
\begin{enumerate}
\item Faute classique : après \emph{si} de condition, pas de conditionnel. Correction : «~Si j'\textbf{avais} (imparfait) le poste, je \textbf{donnerais} (conditionnel dans la principale) le meilleur de moi-même.~»
\item Après \emph{vouloir que} (volonté), le subjonctif s'impose : «~Je veux que vous \textbf{sachiez} que je suis motivé.~» (subjonctif irrégulier de \emph{savoir}).
\item Correct : après \emph{il faut que} (obligation), le subjonctif \emph{sois} est bien employé.
\item Pour exprimer la politesse atténuée, le conditionnel convient mieux que l'indicatif : «~Je \textbf{souhaiterais} vous rencontrer prochainement.~» (radical du futur \emph{souhaiter-} + terminaison \emph{-ais}).
\end{enumerate}
\end{exoresolu}

\begin{exercice}[À vous de jouer]
Mets les verbes entre parenthèses au mode et au temps qui conviennent, puis justifie ton choix (volonté, obligation, condition, politesse\ldots).
\begin{enumerate}
\item «~Si vous (accepter) ma candidature, je (être) très honoré.~»
\item «~Je joins mon CV afin que vous (pouvoir) étudier mon parcours.~»
\item «~Bien que je (être) débutant, j'apprends vite.~»
\item «~Je (vouloir) vous remercier de votre attention.~» (formule de politesse)
\item «~Il est essentiel qu'un technicien (savoir) lire un schéma.~»
\end{enumerate}
\end{exercice}

\begin{corrige}[Exercice : modes au travail]
\begin{enumerate}
\item «~Si vous \textbf{acceptiez} (imparfait, condition) ma candidature, je \textbf{serais} (conditionnel, principale) très honoré.~»
\item «~\ldots afin que vous \textbf{puissiez} étudier mon parcours.~» (subjonctif après la conjonction \emph{afin que}).
\item «~Bien que je \textbf{sois} débutant\ldots~» (subjonctif après \emph{bien que}, concession).
\item «~Je \textbf{voudrais} vous remercier\ldots~» (conditionnel de politesse, plus élégant que «~je veux~»).
\item «~\ldots qu'un technicien \textbf{sache} lire un schéma.~» (subjonctif après \emph{il est essentiel que}, nécessité).
\end{enumerate}
\end{corrige}

\begin{aretenir}[Bilan de la section]
\begin{itemize}
\item Le \cle{mode} exprime l'attitude du locuteur : fait (indicatif), ordre (impératif), souhait ou doute (subjonctif), hypothèse ou politesse (conditionnel).
\item Subjonctif présent : radical de la 3\textsuperscript{e} personne du pluriel + \emph{-e, -es, -e, -ions, -iez, -ent} ; irréguliers à mémoriser : \emph{sois, aie, aille, fasse, puisse, sache, veuille, faille}.
\item Conditionnel présent : radical du futur + terminaisons de l'imparfait (\emph{-ais, -ais, -ait, -ions, -iez, -aient}).
\item Règle d'or : \emph{si} + imparfait $\rightarrow$ conditionnel dans la principale ; jamais de conditionnel après \emph{si}.
\item Dans la lettre de motivation et l'entretien, le subjonctif et le conditionnel de politesse sont des marques de maîtrise et de respect.
\end{itemize}
\end{aretenir}

Ces outils grammaticaux acquis, nous sommes prêts pour la section suivante : la \emph{production guidée} du CV, où chaque phrase rédigée pourra mobiliser à bon escient le subjonctif et le conditionnel.

\coursec{Présenter son CV dans le cadre d'un entretien d'embauche}

\courssub{Découverte : le CV, un texte fonctionnel}

Le \cle{curriculum vitae} (CV) est un \cle{texte fonctionnel} à visée professionnelle : il informe un recruteur sur le parcours, les compétences et le potentiel d'un candidat en vue d'un entretien d'embauche. Au même titre que la lettre de motivation, le rapport de stage ou le compte rendu technique, il obéit à des normes codifiées (structure, concision, objectivité) et mobilise des compétences langagières précises (morphosyntaxe, lexique professionnel, mise en page iconique).

\begin{definition}[Textes fonctionnels et iconiques]
Un \textbf{texte fonctionnel} remplit une fonction sociale précise (postuler, informer, vendre, administrer). Le CV en est l'archétype : il doit être \textbf{lisible en moins de 30 secondes} (lecture diagonale). Un \textbf{texte iconique} associe verbal (mots) et non-verbal (mise en page, police, puces, espaces, photo) pour guider l'œil du lecteur. Maîtriser le CV, c'est donc maîtriser la \textbf{grammaire visuelle} autant que la grammaire syntaxique.
\end{definition}

\begin{savaistu}[Le CV au Cameroun : contexte du Probatoire/Bac Technique]
Dans l'enseignement technique camerounais, le CV est une épreuve certificative (épreuve écrite de production, épreuve orale de présentation). Les jurys attendent : \textbf{rigueur administrative} (exactitude des dates, titres officiels des diplômes : Probatoire, Baccalauréat technique, BT, BTS), \textbf{pertinence technique} (stages en entreprise, maîtrise des logiciels métiers : Sage, AutoCAD, Suite Office, machines-outils) et \textbf{qualité langagière} (zéro faute, registres soutenus, modes subjonctif/conditionnel pour la politesse).
\end{savaistu}

\courssub{Les éléments constitutifs du CV}

Tout CV conforme aux standards professionnels (et aux attendus du MINESEC) comporte cinq rubriques \textbf{obligatoires} et deux \textbf{facultatives}.

\begin{center}
\begin{tabularx}{\linewidth}{|l|X|c|}
\hline
\rowcolor{popblueL} \textbf{Rubrique} & \textbf{Contenu attendu} & \textbf{Statut} \\
\hline
\cle{En-tête / Identité} & Nom, Prénom, Âge / Date de naissance, Lieu de naissance, Nationalité, Adresse, Tél., Email pro, (Photo d'identité) & Obligatoire \\
\hline
\cle{Titre / Objectif} & Intitulé du poste visé ou spécialité (ex. : « Technicien en Maintenance Industrielle — Option Froid ») & Obligatoire \\
\hline
\cle{Formation} & Diplômes du plus récent au plus ancien : intitulé exact, établissement, année, mention éventuelle & Obligatoire \\
\hline
\cle{Expériences professionnelles} & Stages, emplois, jobs d'été, bénévolat : structure, dates, poste, \textbf{3 à 5 missions chiffrées} & Obligatoire \\
\hline
\cle{Compétences} & Techniques (métier), Informatiques (logiciels), Linguistiques (niveau CECRL), Transverselles (permis, premiers secours) & Obligatoire \\
\hline
\cle{Centres d'intérêt} & 2 à 3 activités valorisantes (sport collectif, association, lecture spécialisée) & Facultatif \\
\hline
\cle{Références} & Noms + contacts de 1 à 2 personnes (ancien maître de stage, enseignant) acceptant de recommander & Facultatif \\
\hline
\end{tabularx}
\end{center}

\begin{attention}[Titres officiels des diplômes camerounais]
Écrire \textbf{« Baccalauréat technique série G2 (Comptabilité) »} et non « Bac G2 ». Écrire \textbf{« Probatoire technique série F4 (Électrotechnique) »}. La précision est un gage de sérieux pour le recruteur et le jury d'examen.
\end{attention}

\courssub{Textes iconiques et mise en page du CV}

La mise en page n'est pas décorative : elle structure l'information et guide la lecture diagonale du recruteur.

\begin{methode}[Règles d'or de la mise en page (Grammaire visuelle)]
\begin{enumerate}
\item \textbf{Format} : Une page A4 recto, marges 2 cm, fond blanc, impression laser/jet d'encre qualité.
\item \textbf{Police} : Une seule famille sans empattement (Arial, Calibri, Helvetica, Roboto). Deux tailles max : 14--16 pt pour le nom/titres, 11--12 pt pour le corps.
\item \textbf{Hiérarchie} : Nom en \textbf{très grand} en haut (centré ou gauche). Rubriques en \textbf{gras} + \textbf{MAJUSCULES} ou \textbf{Petites Capitales}. Sous-rubriques en gras simple.
\item \textbf{Aération} : Interligne 1,15 ou 1,5. Espace blanc \textbf{avant} chaque rubrique (6--12 pt). Pas de blocs compacts.
\item \textbf{Puces} : Utiliser des puces simples (●, ■, –) pour les missions. Éviter les tirets longs ou les astérisques.
\item \textbf{Alignement} : Dates alignées à gauche ou droite (cohérence sur tout le doc). Texte justifié ou aligné gauche.
\item \textbf{Photo} : Facultative. Si présente : format identité (3,5 × 4,5 cm), fond neutre, tenue professionnelle, sourire léger, coin haut droit ou gauche. \textbf{Jamais} de photo de vacances, selfie, ou photo découpée.
\item \textbf{Fichier final} : Enregistrer en \textbf{PDF} (nom du fichier : \texttt{Nom\_Prenom\_CV\_Poste.pdf}).
\end{enumerate}
\end{methode}

\begin{popfigure}
\begin{center}
\begin{tikzpicture}[
  node distance=0.4cm and 0.5cm,
  box/.style={rectangle, rounded corners=2pt, draw=popblue, fill=popblue!10, text width=5.5cm, align=center, inner sep=4pt, font=\small},
  title/.style={rectangle, rounded corners=2pt, draw=popdark, fill=popgold, text width=11.5cm, align=center, inner sep=4pt, font=\small\bfseries},
  arrow/.style={-{Stealth[length=2mm]}, thick, poporange}
]
\node[title] (t) {Structure visuelle type d'un CV technique (une page)};
\node[box, below=of t, xshift=-3cm, align=center] (b1){\textbf{EN-TÊTE}\\\small Nom Prénom \\ Coordonnées \\ Photo (opt.)};
\node[box, below=of t, xshift=3cm, align=center] (b2){\textbf{TITRE CIBLE}\\\small Poste visé / Spécialité};
\node[box, below=of b1, align=center] (b3){\textbf{FORMATION}\\\small Diplômes récents $\to$ anciens \\ Mentions, établissements};
\node[box, below=of b2, align=center] (b4){\textbf{EXPÉRIENCES}\\\small Stages / Emplois \\ Missions + Chiffres};
\node[box, below=of b3, align=center] (b5){\textbf{COMPÉTENCES}\\\small Techniques / Informatiques \\ Langues / Transverses};
\node[box, below=of b4, align=center] (b6){\textbf{INTÉRÊTS / RÉF.}\\\small 2--3 items véridiques \\ Références (opt.)};
\draw[arrow] (t) -- (b1);
\draw[arrow] (t) -- (b2);
\draw[arrow] (b1) -- (b3);
\draw[arrow] (b2) -- (b4);
\draw[arrow] (b3) -- (b5);
\draw[arrow] (b4) -- (b6);
\end{tikzpicture}
\end{center}
\legende{Architecture visuelle standard : lecture en Z (haut gauche $\to$ haut droit $\to$ bas gauche $\to$ bas droit). L'en-tête et le titre captent l'attention en 3 secondes.}
\end{popfigure}

\courssub{Langue : les valeurs du subjonctif et du conditionnel}

Ces deux modes sont indispensables dans le CV (accroche, titre) et surtout dans la lettre de motivation qui l'accompagne. Leur maîtrise est évaluée au critère « Langue » (5 pts sur 20).

\begin{propriete}[Valeurs du conditionnel (présent et passé) au service de la candidature]
\begin{itemize}
\item \textbf{Politesse / Atténuation} : « Je \textbf{souhaiterais} vous rencontrer », « Je \textbf{serais} honoré de rejoindre votre équipe ». (Remplace l'impératif ou l'indicatif trop direct).
\item \textbf{Hypothèse / Projection} : « Une collaboration \textbf{serait} mutuellement bénéfique », « Mes compétences \textbf{permettraient} d'optimiser vos stocks ».
\item \textbf{Condition} : « Si mon profil vous intéresse, je \textbf{resterais} à votre disposition pour un entretien ».
\end{itemize}
\end{propriete}

\begin{propriete}[Valeurs du subjonctif (présent) au service de la candidature]
Le subjonctif s'impose après les verbes de \textbf{volonté}, \textbf{souhait}, \textbf{nécessité}, \textbf{émotion}, \textbf{doute} (que, pour que, il faut que, il est important que...).
\begin{itemize}
\item \textbf{Volonté / Souhait} : « Je \textbf{veux} que ma candidature \textbf{retienne} votre attention », « Je \textbf{souhaite} que vous \textbf{examiniez} mon dossier ».
\item \textbf{Nécessité / Obligation} : « Il \textbf{faut} que je \textbf{puisse} mettre mes compétences à l'épreuve », « Il est \textbf{indispensable} que le stage \textbf{soit} validé ».
\item \textbf{Finalité} : « Je vous adresse mon CV \textbf{afin que} vous \textbf{puissiez} évaluer mon profil ».
\end{itemize}
\end{propriete}

\begin{attention}[Pièges grammaticaux fréquents au Bac]
\begin{itemize}
\item \textbf{Confusion Indicatif/Subjonctif} : *\"Je veux que vous \textbf{examinez} mon dossier\" $\to$ \textbf{examiniez} (subjonctif).
\item \textbf{Conditionnel au lieu du Subjonctif} : *\"Il faut que je \textbf{serais} prudent\" $\to$ \textbf{sois} (subjonctif).
\item \textbf{Accord du participe passé} : \"Les factures que j'ai \textbf{saisies}\" (COD placé avant). \"Le stage que j'ai \textbf{effectué}\" (COD après).
\item \textbf{Conditionnel passé vs Subjonctif passé} : \"J'aurais aimé que vous \textbf{ayez reçu} ma lettre\" (subjonctif passé), pas *\"que vous \textbf{auriez reçu}\".
\end{itemize}
\end{attention}

\courssub{Production guidée : rédiger son CV}

Rédiger un CV suit une démarche processuelle en cinq étapes. Aucune ne peut être sautée.

\begin{popfigure}
\begin{center}
\begin{tikzpicture}[
  node distance=0.5cm,
  etape/.style={rectangle, rounded corners=3pt, draw=popblue, fill=popblue!15, text width=11cm, align=left, inner sep=6pt, font=\small},
  fleche/.style={-{Stealth[length=2.5mm]}, thick, poporange}
]
\node[etape] (e1) {\textbf{Étape 1 — Collecte} : Remplir sa \textbf{fiche d'identité professionnelle} (état civil, diplômes officiels, expériences datées, compétences techniques/logicielles/linguistiques, centres d'intérêt).};
\node[etape, below=of e1] (e2) {\textbf{Étape 2 — Choix du plan} : \textbf{Chronologique} (parcours régulier, formation $\to$ stages) \textbf{OU} \textbf{Par compétences} (peu d'expérience, reconversion, trous dans le parcours).};
\node[etape, below=of e2] (e3) {\textbf{Étape 3 — Rédaction ciblée} : \textbf{Verbes d'action} (infinitif pour poste actuel, participe passé pour passé), \textbf{données chiffrées}, \textbf{mots-clés de l'offre}. Adapter \textbf{chaque CV} à \textbf{chaque offre}.};
\node[etape, below=of e3] (e4) {\textbf{Étape 4 — Relecture experte} : Orthographe lexicale, accords (sujet/verbe, participe passé), \textbf{modes} (conditionnel politesse, subjonctif volonté), ponctuation, cohérence dates/titres. Faire relire par un tiers.};
\node[etape, below=of e4] (e5) {\textbf{Étape 5 — Mise en page finale} : Application règles iconiques (1 page, police unique, hiérarchie, puces, aération, PDF, nom de fichier pro). Photo d'identité si choisie.};
\draw[fleche] (e1) -- (e2);
\draw[fleche] (e2) -- (e3);
\draw[fleche] (e3) -- (e4);
\draw[fleche] (e4) -- (e5);
\end{tikzpicture}
\end{center}
\legende{Processus de production du CV : itératif (on peut revenir à l'étape 3 après la relecture).}
\end{popfigure}

\courssub{Étape 1 : Collecte — La fiche d'identité professionnelle}

\begin{methode}[Remplir sa fiche d'identité professionnelle (Brouillon documentaire)]
\begin{enumerate}
\item \textbf{État civil} : Nom, Prénom, Date/Lieu de naissance, Nationalité, Adresse, Tél., Email (\texttt{prenom.nom@domaine.pro}), LinkedIn (optionnel, si à jour).
\item \textbf{Formation} : \textbf{Intitulé exact} (ex. \emph{Baccalauréat Technique Série G2 — Comptabilité et Gestion}), Établissement, Année, Mention. Ordre antéchronologique.
\item \textbf{Expériences} : Pour chaque ligne : Structure, Ville, Dates (Mois/Année), Intitulé poste, \textbf{3--5 missions} commençant par verbe d'action + résultat chiffré.
\item \textbf{Compétences} : \textit{Techniques} (ex. \emph{Tenue comptabilité fournisseurs/clients, Rapprochement bancaire, Déclaration TVA}), \textit{Informatiques} (ex. \emph{Excel TCD, Sage 100, Word, PowerPoint}), \textit{Langues} (ex. \emph{Français C1, Anglais B1, Bassa A2}), \textit{Permis} (B, C...).
\item \textbf{Centres d'intérêt} : Sélectifs et véridiques (ex. \emph{Football (capitaine équipe lycée), Lecture presse économique, Bénévolat Croix-Rouge}).
\end{enumerate}
\end{methode}

\courssub{Étape 2 : Choix du plan — Chronologique vs Par compétences}

\begin{center}
\begin{tabularx}{\linewidth}{|l|X|X|}
\hline
\rowcolor{popblueL} \textbf{Critère de choix} & \textbf{Plan chronologique (antéchronologique)} & \textbf{Plan par compétences (fonctionnel)} \\
\hline
\textbf{Logique} & Temps : du plus récent au plus ancien & Thèmes : regroupement par savoir-faire \\
\hline
\textbf{Profil type} & Élève/Étudiant, parcours linéaire, stages cohérents & Autodidacte, reconversion, trous, expérience dispersée \\
\hline
\textbf{Avantage Bac} & Rassurant pour le jury : lisibilité immédiate du parcours scolaire & Masque les faiblesses chronologiques, met en avant le \textbf{savoir-faire} \\
\hline
\textbf{Risque} & Met en lumière les périodes d'inactivité & Peut frustrer le recruteur qui cherche les dates/entreprises \\
\hline
\textbf{Recommandation Tle Tech} & \textbf{Par défaut : CHRONOLOGIQUE.} C'est le standard attendu pour un primo-demandeur d'emploi. & Uniquement si le chronologique dessert manifestement le candidat. \\
\hline
\end{tabularx}
\end{center}

\courssub{Étape 3 : Rédaction ciblée — Verbes d'action et adaptation}

\begin{exemplebox}[Transformer une expérience banale en compétence prouvée]
\begin{tabularx}{\linewidth}{|l|X|}
\hline
\rowcolor{poporangeL} \textbf{Faible (Passif / Vague)} & \textbf{Fort (Actif / Précis / Chiffré)} \\
\hline
« J'ai fait un stage en comptabilité. » & « \textbf{Tenu} la comptabilité fournisseurs (50 factures/semaine) et \textbf{réalisé} les rapprochements bancaires mensuels sous Sage 100. » \\
\hline
« Je connais Excel. » & « \textbf{Maîtrise} Excel : tableaux croisés dynamiques, RECHERCHEV, macros simples pour automatiser les saisies. » \\
\hline
« J'ai réparé des machines. » & « \textbf{Diagnostiqué} et \textbf{réparé} 15 pannes hydrauliques sur presse à injecter, \textbf{réduit} le temps d'arrêt de 20\%. » \\
\hline
\end{tabularx}
\end{exemplebox}

\begin{methode}[Adapter son CV à une offre (Mots-clés)]
\begin{enumerate}
\item \textbf{Analyser l'offre} : Surligner 5--7 compétences/mots-clés exigés (ex. \emph{Sage, TVA, Rigueur, Anglais technique, Permis B}).
\item \textbf{Mirroir} : Dans sa fiche d'identité, cocher les éléments qui y répondent.
\item \textbf{Prioriser} : Placer ces éléments \textbf{en haut} des rubriques Compétences / Expériences. Reprendre le vocabulaire exact de l'annonce.
\item \textbf{Epurer} : Supprimer ou réduire ce qui est sans lien (ex. stage boulangerie pour poste comptable $\to$ garder seulement si compétence transférable : gestion caisse, rigueur).
\end{enumerate}
\end{methode}

\begin{attention}[Erreur éliminatoire : Le CV « passe-partout »]
Envoyer le même CV à 20 offres différentes signale au recruteur : « Je n'ai pas lu votre annonce, je ne sais pas ce que vous cherchez, je manque de rigueur. » \textbf{Un CV = Une offre.} Le temps d'adaptation (15 min) conditionne l'obtention de l'entretien.
\end{attention}

\courssub{Étape 4 : Relecture — La check-list « Zéro Faute »}

\begin{center}
\begin{tabularx}{\linewidth}{|c|X|}
\hline
\rowcolor{popblueL} \textbf{Point de contrôle} & \textbf{Action concrète} \\
\hline
\textbf{1. Orthographe lexicale} & Vérifier les pièges : \emph{accueil, expérience, secrétariat, hiérarchie, pièce jointe, pièce justificative, échéancier, bénéficiaire}. Dictionnaire obligatoire. \\
\hline
\textbf{2. Accords} & Participe passé (COD avant ?), Sujet/Verbe (attention aux collectifs : \emph{la majorité des dossiers \textbf{sont} traités}). \\
\hline
\textbf{3. Modes (Subj./Cond.)} & Accroche/Lettre : \emph{Je \textbf{souhaiterais} (Cond.) que vous \textbf{examiniez} (Subj.) mon dossier}. Pas d'indicatif après \emph{il faut que, pour que, bien que}. \\
\hline
\textbf{4. Ponctuation / Typographie} & Pas de point final aux titres/puces. Deux-points avant énumération. Espace insécable avant \og ; : ? ! \fg{}. Guillemets français. \\
\hline
\textbf{5. Cohérence factuelle} & Dates continues ? Titres exacts (Probatoire, Bac Tech) ? Noms d'entreprises orthographiés ? Numéros tél/email valides ? \\
\hline
\end{tabularx}
\end{center}

\begin{methode}[Technique de relecture croisée (Binôme)]
\begin{enumerate}
\item Imprimer le CV (on voit mieux les fautes sur papier).
\item Lire \textbf{à l'envers} (dernière ligne $\to$ première) pour casser la lecture anticipée et voir l'orthographe.
\item Échanger avec un camarade : lui lit \textbf{à voix haute} (entend les lourdeurs), vous corrigez sur votre exemplaire.
\item Valider ensemble la grille d'évaluation (voir fin de leçon).
\end{enumerate}
\end{methode}

\courssub{Étape 5 : Mise en page finale — Export PDF}

\begin{exemplebox}[Modèle de structure finale (Texte brut pour mise en page)]
\textbf{FOPA NGUEFACK Steve Arnaud} \hfill \textit{Technicien Supérieur en Maintenance Industrielle — Option Froid}\\
Douala, Bonabéri — +237 6 90 00 00 00 — steve.fopa@email.pro — \textit{Permis B} \\[4pt]
\textbf{FORMATION}
\begin{itemize}
\item 2025--2026 : \textbf{Baccalauréat Technique F4 (Maintenance Industrielle)}, Lycée Technique de Douala-Bassa. \textit{En préparation}.
\item 2024 : \textbf{Probatoire Technique F4}, Mention Assez Bien.
\item 2022 : \textbf{BEPC}, Collège Polyvalent de Bonabéri.
\end{itemize}
\textbf{EXPÉRIENCES PROFESSIONNELLES}
\begin{itemize}
\item \textbf{Juil.--Août 2025} — \textbf{Stagiaire Maintenance}, \textit{SOCAVER (Douala)}. \textbf{Missions} : Maintenance préventive niveau 1 compresseurs frigorifiques (12 unités) ; Diagnostic pannes électriques (multimètre, schémas) ; Rédigé 8 fiches d'intervention ; Respect strict consignes QHSE.
\item \textbf{Juil.--Août 2024} — \textbf{Aide-technicien}, \textit{Entreprise Froid Service (Douala)}. Pose/gaine cuivre, tirage câbles, test étanchéité azote.
\end{itemize}
\textbf{COMPÉTENCES TECHNIQUES}
\begin{itemize}
\item Froid/Clim : Cycle frigorifique, fluides (R134a, R404A, R290), brasage fort/brasage tendre, mise en service, recherche fuite.
\item Électricité : Lecture schémas unifilaires/multifilaires, câblage armoires, variateurs de fréquence, automatismes (API Siemens LOGO!).
\item Outillage : Multimètre, pince ampèremétrique, station de brasage, pompe à vide, manifolds.
\end{itemize}
\textbf{COMPÉTENCES INFORMATIQUES \& LANGUES}
\begin{itemize}
\item Bureautique : Word, Excel (suivi maintenance), PowerPoint. GMAO : Initiation \textit{Maintenance Assistant}.
\item Français : Courant (C1) — Anglais : Technique (B1 - lecture plans/notice) — Bassa : Notions.
\end{itemize}
\textbf{CENTRES D'INTÉRÊT}
\begin{itemize}
\item Football (Milieu terrain, équipe lycée — esprit équipe, endurance).
\item Bricolage électronique (réparation petits appareils, veille tech YouTube).
\end{itemize}
\end{exemplebox}

\courssub{TP Intégré : Produire son CV réponse à une offre réelle}

\begin{exoresolu}[Cas pratique : Offre « Technicien de Maintenance Froid » — Douala]
\textbf{Énoncé.} L'entreprise \textit{ClimService Douala} recrute : « Technicien maintenance froid/clim. Bac Tech F4 ou BT Froid. Maîtrise cycle frigorifique, brasage, lecture plans, multimètre. Permis B. Rigueur, autonomie. Stage SOCAVER/SEMC apprécié. » Produire le CV adapté.
\tcblower
\textbf{Solution commentée.}
\textbf{1. Analyse offre (Mots-clés) :} F4/BT Froid, Cycle frigorifique, Brasage, Lecture plans, Multimètre, Permis B, Rigueur, Autonomie, Stage SOCAVER.
\textbf{2. Fiche identité :} Le candidat (Steve, ex. ci-dessus) possède \textbf{tous} les mots-clés.
\textbf{3. Plan :} Chronologique (parcours scolaire régulier + stages cohérents).
\textbf{4. Rédaction ciblée :}
\begin{itemize}
\item \textit{Titre} : Reprend l'intitulé exact de l'offre.
\item \textit{Formation} : Bac Tech F4 mis en premier (critère diplôme).
\item \textit{Expérience SOCAVER} : Détaillée en premier (critère « stage apprécié »), verbes d'action + chiffrage (12 unités, 8 fiches).
\item \textit{Compétences} : Regroupées par thèmes \textbf{Froid}, \textbf{Électricité}, \textbf{Outillage} $\to$ miroir exact des exigences techniques.
\item \textit{Permis B} : Remonté dans l'en-tête (critère éliminatoire souvent).
\item \textit{Centres d'intérêt} : Football (autonomie/équipe), Bricolage électronique (passion technique, autonomie).
\end{itemize}
\textbf{5. Relecture :} Vérification accords (\emph{fiches rédigées}, \emph{pannes diagnostiquées}), modes (lettre : \emph{Je souhaiterais intégrer...}), cohérence dates.
\textbf{6. Mise en page :} Police Calibri 11/14, puces ●, interligne 1,15, PDF \texttt{Fopa\_Steve\_CV\_TechnicienFroid.pdf}.
\end{exoresolu}

\begin{exercice}[Votre tour : CV « Sur mesure »]
\begin{enumerate}
\item Trouvez une offre réelle (presse \emph{Mutations}, \emph{Cameroon Tribune}, sites \emph{Emploi.cm}, \emph{Jobibo}, LinkedIn) correspondant à \textbf{votre spécialité technique}.
\item Remplissez votre fiche d'identité professionnelle (brouillon).
\item Choisissez votre plan (chronologique par défaut).
\item Rédigez le CV complet (1 page) en appliquant la méthode « Mots-clés / Verbes d'action / Chiffres ».
\item Faites la relecture croisée avec un camarade (grille ci-dessous).
\item Déposez le PDF final sur la plateforme / remettez version papier.
\end{enumerate}
\end{exercice}

\courssub{Présenter son CV lors d'un entretien d'embauche}

Le CV écrit ouvre la porte ; la présentation orale (entretien) la franchit. Au Bac Technique, l'oral de présentation du CV est une épreuve certificative (souvent 10--15 min : 5 min exposé, 5--10 min questions).

\begin{methode}[Structure de l'exposé oral de présentation (5 minutes -- « Pitch »)]
\begin{enumerate}
\item \textbf{Accroche (30 s)} : « Bonjour, je suis \textbf{Nom Prénom}, âgé de \textbf{X ans}, titulaire du \textbf{Baccalauréat Technique Série [Spécialité]} obtenu en \textbf{Année} au \textbf{Lycée [Nom]} avec mention \textbf{[Mention]}. Je me présente aujourd'hui pour le poste de \textbf{[Intitulé exact du poste]}. »
\item \textbf{Parcours \& Motivation (1 min 30)} : « Mon parcours allie \textbf{théorie} (modules clés : \emph{citer 2--3 matières phares}) et \textbf{pratique} (stage majeur chez \emph{Entreprise} où j'ai \emph{mission principale + résultat chiffré}). Ce qui m'attire chez vous, c'est \emph{valeur/projet entreprise} et la perspective de \emph{mission technique précise}. »
\item \textbf{Compétences « Preuve » (2 min)} : « Trois atouts majeurs pour ce poste : \textbf{1.} \emph{Compétence technique dure} (ex. \emph{Maîtrise Sage 100 / Brasage R290 / Programmation API}) prouvée par \emph{stage/projet}. \textbf{2.} \emph{Compétence transversale} (ex. \emph{Rigueur comptable / Respect QHSE / Anglais technique}) illustrée par \emph{exemplebox}. \textbf{3.} \emph{Savoir-être} (ex. \emph{Autonomie, esprit d'équipe, capacité d'apprentissage}) validé par \emph{responsabilité associative/sportive}. »
\item \textbf{Projection \& Disponibilité (30 s)} : « Je suis \textbf{immédiatement disponible}, \textbf{mobile} sur [Zone], titulaire du \textbf{Permis B}. Je \textbf{souhaiterais} (Conditionnel) m'investir durablement dans votre structure et \textbf{contribuer} à \emph{objectif cité dans l'offre}. Je vous remercie de votre attention et reste à votre disposition pour vos questions. »
\end{enumerate}
\end{methode}

\begin{popfigure}
\begin{center}
\begin{tikzpicture}[
  mindmap, concept color=popblue, font=\small,
  level 1/.style={level distance=3.5cm, sibling angle=90},
  level 2/.style={level distance=2.8cm, sibling angle=45},
  every node/.style={concept, execute at begin node=\bfseries}
]
\node[concept color=popblue, align=center]{Entretien\\Réussi}
  child[concept color=popgreen] {node {Préparation}
    child {node {Offre analysée}}
    child {node {CV/Offre alignés}}
    child {node {Questions types répétées}}
  }
  child[concept color=poporange] {node {Présentation (Pitch)}
    child {node {Accroche claire}}
    child {node {Preuves chiffrées}}
    child {node {Projection (Nous)}}
  }
  child[concept color=poppurple] {node {Attitude Pro}
    child {node {Tenue soignée}}
    child {node {Contact visuel}}
    child {node {Écoute active}}
    child {node {Langage corporel ouvert}}
  }
  child[concept color=poppink] {node {Gestion Questions}
    child {node {STAR (Situation, Tâche, Action, Résultat)}}
    child {node {Honnêteté (\"Je ne sais pas mais je saurai\")}}
    child {node {Salaire : fourchette réaliste}}
  };
\end{tikzpicture}
\end{center}
\legende{Carte mentale des 4 piliers de la réussite à l'entretien d'embauche technique.}
\end{popfigure}

\begin{exemplebox}[Méthode STAR pour répondre aux questions comportementales]
\textbf{Q : « Parlez-moi d'une difficulté technique rencontrée en stage et comment vous l'avez résolue. »}
\begin{itemize}
\item \textbf{S (Situation)} : « Lors de mon stage chez SOCAVER, un compresseur Bitzer déclenchait sans cesse le pressostat HP. »
\item \textbf{T (Tâche)} : « Mon tuteur m'a confié le diagnostic et la remise en service sous 2h. »
\item \textbf{A (Action)} : « J'ai vérifié le condenseur (encrassé), nettoyé les ailettes, contrôlé le ventilateur (condensateur HS), remplacé le condensateur, purgé le circuit, rechargé R134a selon notice. »
\item \textbf{R (Résultat)} : « Pression HP stabilisée à 14 bar, machine tournée 48h sans alerte. Tuteur validé fiche intervention. J'ai appris l'importance de la maintenance préventive planning. »
\end{itemize}
\end{exemplebox}

\begin{attention}[Erreurs fatales en entretien]
\begin{itemize}
\item \textbf{Arriver en retard} (sans prévenir) / Tenue inadaptée (baskets sales, t-shirt, casquette).
\item \textbf{Lire son CV} (le recruteur l'a déjà lu). Il faut le \textbf{raconter} et l'\textbf{incarner}.
\item \textbf{Critiquer} ancien employeur, école, collègues.
\item \textbf{Mensonge} sur compétences (test technique possible immediate : Excel, brasage, lecture plan).
\item \textbf{Ne pas poser de question} à la fin (signe de désintérêt). Préparer 2 questions : « Comment se passe l'intégration des jeunes techniciens ? », « Quels sont les défis techniques actuels du service ? ».
\end{itemize}
\end{attention}

\courssub{Intégration, compte rendu et remédiation}

\begin{exercice}[Situation d'intégration finale : « Du CV à l'embauche »]
\textbf{Contexte :} Vous répondez à l'offre trouvée à l'exercice précédent (ou à une offre fournie par l'enseignant).
\begin{enumerate}
\item \textbf{Écrit (2h)} : Produire le \textbf{CV adapté} (1 page PDF) + la \textbf{Lettre de motivation} (structure « Vous, Moi, Nous », 1 page, modes conditionnel/subjonctif maîtrisés).
\item \textbf{Oral (10 min / candidat)} : Simuler l'entretien devant un jury (enseignant + 2 élèves observateurs).
\begin{itemize}
\item 5 min : Pitch présentation (méthode ci-dessus).
\item 5 min : Questions jury (technique, motivation, situations STAR, salaire, disponibilité).
\end{itemize}
\item \textbf{Auto-évaluation (Grille)} : Chaque candidat remplit la grille ci-dessous \textbf{immédiatement} après son passage. L'enseignant complète sa grille. Comparaison $\to$ Remédiation ciblée.
\end{enumerate}
\end{exercice}

\begin{exemplebox}[Grille d'auto-évaluation / évaluation enseignant (Sur 20)]
\begin{center}
\begin{tabularx}{\linewidth}{|l|X|c|c|}
\hline
\rowcolor{popblueL} \textbf{Critère} & \textbf{Indicateurs de réussite} & \textbf{Auto /5} & \textbf{Ens. /5} \\
\hline
\textbf{CV : Structure \& Contenu} & Rubriques complètes, plan adapté, infos exactes, verbes d'action, chiffrage, adaptation offre & & \\
\hline
\textbf{CV : Langue \& Présentation} & Zéro faute, accords, modes (subj/cond), mise en page aérée, police pro, PDF nommage correct & & \\
\hline
\textbf{Lettre : Structure \& Contenu} & Structure « Vous, Moi, Nous », cohérence CV, motivation sincère, politesse conditionnel, subjonctif correct & & \\
\hline
\textbf{Oral : Pitch \& Preuves} & Accroche pro, 3 atouts prouvés (STAR), projection « Nous », temps respecté (5 min), aisance & & \\
\hline
\textbf{Oral : Attitude \& Réponses} & Tenue, regard, voix, écoute, STAR sur questions comportementales, honnêteté, questions finales pertinentes & & \\
\hline
\textbf{TOTAL} & & \textbf{/25} & \textbf{/25} \\
\hline
\end{tabularx}
\end{center}
\emph{Note : Le total est sur 25 pour permettre une pondération (CV+Lettre = 10, Oral = 15) ou ramené sur 20 selon modalités établissement.}
\end{exemplebox}

\begin{popfigure}
\begin{center}
\begin{tikzpicture}
\begin{axis}[
  ybar, bar width=0.6cm,
  width=0.9\linewidth, height=6cm,
  ymin=0, ymax=5.5,
  symbolic x coords={CV Struct/Contenu, CV Langue/Pres, Lettre, Oral Pitch, Oral Attitude},
  xtick=data,
  xticklabel style={rotate=45, anchor=east, font=\tiny, align=center},
  ylabel={Note sur 5},
  ytick={0,1,2,3,4,5},
  nodes near coords,
  every node near coord/.append style={font=\scriptsize},
  axis line style={popline},
  tick label style={font=\small},
  legend style={at={(0.5,-0.35)}, anchor=north, legend columns=2, font=\small}
]
\addplot[fill=popblue, draw=popdark] coordinates {(CV Struct/Contenu,4.5) (CV Langue/Pres,4) (Lettre,4.5) (Oral Pitch,3.5) (Oral Attitude,4)};
\addlegendentry{Candidat Modèle (Objectif)}
\addplot[fill=poporange, draw=popdark, opacity=0.8] coordinates {(CV Struct/Contenu,3) (CV Langue/Pres,2.5) (Lettre,3) (Oral Pitch,2) (Oral Attitude,2.5)};
\addlegendentry{Candidat Type (Besoin remédiation Oral)}
\end{axis}
\end{tikzpicture}
\end{center}
\legende{Profil type de notes : les élèves de Terminale Technique sont souvent plus à l'aise à l'écrit (CV/Lettre) qu'à l'oral. La remédiation doit cibler l'entraînement oral (simulations, enregistrement vidéo, feedback pairs).}
\end{popfigure}

\begin{aretenir}[Synthèse de la Leçon 24 : Présenter son CV]
\begin{enumerate}
\item \textbf{CV = Texte fonctionnel + Iconique} : 1 page, structure codifiée (Identité, Titre, Formation, Expériences, Compétences), mise en page rigoureuse (police, puces, aération, PDF).
\item \textbf{Production en 5 étapes} : Collecte (fiche identité) $\to$ Choix plan (Chronologique par défaut) $\to$ Rédaction ciblée (Verbes d'action + Chiffres + Mots-clés offre) $\to$ Relecture experte (Zéro faute, Modes) $\to$ Mise en page finale.
\item \textbf{Langue : Subjonctif / Conditionnel} : Conditionnel = politesse / projection (\emph{je souhaiterais, ce serait}). Subjonctif = après volonté/nécessité (\emph{je veux que vous examiniez, il faut que je puisse}).
\item \textbf{Lettre de motivation} : Structure « \textbf{Vous, Moi, Nous} ». Cohérence absolue avec le CV. Complément, pas répétition.
\item \textbf{Entretien Oral} : Pitch 5 min (Accroche, Parcours/Preuves STAR, Projection). Attitude pro (tenue, regard, écoute). Méthode STAR pour questions comportementales. Questions finales préparées.
\item \textbf{Évaluation} : 4--5 critères (Structure, Contenu, Langue, Présentation, Oral) notés sur 5. Remédiation prioritaire sur l'oral (simulations répétées).
\end{enumerate}
\end{aretenir}

\coursec{Le CV : texte fonctionnel et éléments constitutifs}

\courssub{Découverte : le CV, un texte fonctionnel à part entière}

\begin{definition}[Texte fonctionnel]
Un \textbf{texte fonctionnel} est un texte \emph{utilitaire}, rédigé dans un but précis (informer, convaincre, prescrire, candidater) et soumis à des \textbf{contraintes de forme et de fond} strictes (normes, codes, attentes du destinataire). Il se distingue du texte littéraire par son efficacité communicationnelle immédiate.
\end{definition}

\begin{propriete}[Le Curriculum Vitæ (CV) : définition et finalité]
Le CV (du latin \emph{curriculum vitæ}, « parcours de vie ») est un \textbf{document de synthèse} qui présente, de manière structurée et sélective, le \textbf{parcours scolaire, universitaire, professionnel et personnel} d'un candidat. Sa finalité unique est d'\textbf{obtenir un entretien d'embauche}. Il ne raconte pas une vie : il \emph{vend} un profil pour un poste précis.
\end{propriete}

\begin{aretenir}[Trois principes d'or du CV technique]
\begin{enumerate}
    \item \textbf{Pertinence (Ciblage) :} Un CV \emph{master} existe, mais on envoie toujours une version \emph{ciblée} (mots-clés de l'offre, compétences mises en avant).
    \item \textbf{Lisibilité (Forme) :} Lecture en diagonale < 30 secondes (règle du « F-pattern »). Mise en page aérée, police unique (ex. Calibri, Arial, Roboto 11-12), gras pour les repères.
    \item \textbf{Véracité (Éthique) :} Tout doit être vérifiable (attestations, certificats, portfolio). Le mensonge est éliminatoire et destructeur de réputation.
\end{enumerate}
\end{aretenir}

\courssub{Les éléments constitutifs du CV technique (Norme camerounaise)}

\begin{methode}[Structure standard du CV (Ordres chronologique inversé recommandé)]
\begin{enumerate}
    \item \textbf{En-tête (Identité) :} Nom \textbf{PRÉNOM(S)}, Adresse (Quartier, Ville), Téléphone (WhatsApp pro), Email \textbf{professionnel} (prenom.nom@domaine.com), LinkedIn/Profil pro (optionnel), \textbf{Photo d'identité récente} (fond neutre, tenue pro, format 3,5x4,5 cm) — \emph{obligatoire au Cameroun}.
    \item \textbf{Accroche / Titre ciblé :} Une ligne sous la photo. \emph{Ex : « Technicien Supérieur en Électrotechnique (Terminale F4) — Recherche stage Maintenance Industrielle »}. Pas de « CV de... ».
    \item \textbf{Formation (Diplômes) :} Du plus récent au plus ancien. \textbf{Intitulé exact du diplôme}, Série/Spécialité, Établissement, Ville, Année, \textbf{Mention} (si Assez Bien / Bien / Très Bien). \emph{Ex : Baccalauréat Technique F4 Électrotechnique, Lycée Technique de Douala-Bassa, 2024, Mention Assez Bien}.
    \item \textbf{Expériences Professionnelles / Stages :} Du plus récent au plus ancien. \textbf{Période (Mois/Année)}, \textbf{Intitulé du poste}, \textbf{Structure}, \textbf{Ville}. \textbf{3 à 5 puces d'actions/résultats} (verbes d'action, méthode STAR, vocabulaire métier, chiffrage).
    \item \textbf{Compétences Techniques (Savoir-faire) :} Regroupées par domaine. \emph{Ex : Maintenance (Préventive/Curative HT/BT, GMAO), Dessin (AutoCAD, SolidWorks), Énergie (Solaire, Groupe électrogène)}. Niveau : Maîtrise / Bonnes connaissances / Notions.
    \item \textbf{Compétences Transverses / Soft Skills :} 3 à 5 items illustrés. \emph{Ex : Rigueur (0 erreur GMAO stage), Esprit d'équipe (PFE binôme), Adaptabilité (Stage milieu inconnu)}.
    \item \textbf{Langues \& Informatique :} Niveau CECRL (A1-C2) ou descriptif (Lu/Écrit/Parlé). Logiciels bureautiques (Pack Office), techniques (CAO, GMAO, Programmation), permis de conduire (catégorie).
    \item \textbf{Centres d'intérêt / Vie associative :} Sélectifs, valorisants (Sport compétition, Bénévolat technique, Club robotique, Musique). Éviter : « Lecture, Voyage, Cinéma » (trop vagues).
    \item \textbf{Références (Optionnel) :} « Sur demande » ou 1-2 noms (Tuteur stage, Professeur principal) avec accord préalable.
\end{enumerate}
\end{methode}

\begin{exemplebox}[Extrait CV — Rubrique Expérience (Exemple Terminale F3)]
\textbf{Stage Technicien Maintenance — SOCIM Bafoussam} \hfill \textit{Juin -- Juillet 2023 (4 sem.)}
\begin{itemize}
    \item Maintenance préventive broyeurs à boulets (cimenterie) : relevé jeux paliers, contrôle revêtements, graissage selon gammes.
    \item Saisie interventions sur GMAO (Logiciel \emph{MAINTI4}) : 45 fiches renseignées, 0 erreur de codage.
    \item Participation arrêt annuel Unité N°2 : démontage/remontage réducteur 150 kW (équipe 4 pers., respect consignes verrouillage).
\end{itemize}
\emph{Apports :} Lecture plans hydrauliques/pneumatiques, normes sécurité (EPI, PT), rapport de stage soutenu 16/20.
\end{exemplebox}

\coursec{Mise en page et textes iconiques du CV}

\courssub{La réception des textes iconiques : codes visuels et lecture en « F »}

\begin{propriete}[Sémiologie de la mise en page]
Le CV est un \textbf{texte iconique} (image + texte). Le recruteur le \emph{scanne} visuellement avant de le lire. La mise en page guide l'œil selon un parcours en \textbf{F} : ligne haute (Nom/Titre), barre verticale gauche (Rubriques/Dates), lignes horizontales (Intitulés clés).
\end{propriete}

\begin{methode}[Règles de mise en page professionnelle (Check-list visuelle)]
\begin{enumerate}
    \item \textbf{Format :} A4, portrait, marges 2 cm (min 1,5 cm). \textbf{1 page max} (jeune diplômé).
    \item \textbf{Police :} Sans empattement (Sans-serif) : Calibri, Arial, Helvetica, Roboto, Open Sans. Taille 11 ou 12. \textbf{Une seule police} (gras/italique pour hiérarchiser).
    \item \textbf{Couleur :} Sobriété. 1 couleur d'accent (Bleu marine \#003366, Vert sapin \#005533, Gris anthracite \#333333) pour titres/lignes. Fond blanc. \textbf{Impression N\&B testée}.
    \item \textbf{Alignement :} Dates à gauche (colonne étroite), Contenu à droite (colonne large). Puces \textbullet\ ou -- alignées. Justifié ou aligné gauche.
    \item \textbf{Espace :} Interligne 1,15. Espace avant/après rubrique (6-12 pt). Pas de blocs compacts.
    \item \textbf{Fichier :} Export \textbf{PDF} (nom : \texttt{NOM\_Prenom\_CV\_Poste.pdf}). Pas de .docx, .odt, .jpg.
\end{enumerate}
\end{methode}

\begin{popfigure}
\begin{tikzpicture}[
    mindmap,
    concept color=popblue,
    level 1/.style={level distance=3.5cm, sibling angle=45},
    level 2/.style={level distance=2.8cm, sibling angle=45},
    every node/.style={font=\small, text width=2.2cm, align=center},
    grow cyclic
]
\node[concept, text width=2.5cm] {\textbf{STRUCTURE VISUELLE CV}}
    child[concept color=poporange] { node[concept, align=center]{\textbf{EN-TÊTE}\\(Identité, Photo, Titre, Contact)} }
    child[concept color=popgreen] { node[concept, align=center]{\textbf{FORMATION}\\(Diplômes, Mentions, Projets)} }
    child[concept color=poppurple] { node[concept, align=center]{\textbf{EXPÉRIENCES}\\(Stages, Jobs, PFE, Actions STAR)} }
    child[concept color=poppink] { node[concept, align=center]{\textbf{COMPÉTENCES}\\(Techniques, Transverses, Outils)} }
    child[concept color=popgold] { node[concept, align=center]{\textbf{LANGUES / IT}\\(Niveaux, Logiciels, Permis)} }
    child[concept color=popmuted] { node[concept, align=center]{\textbf{INTÉRÊTS}\\(Sport, Associatif, Bénévolat)} };
\end{tikzpicture}
\legende{Figure 1 : Carte mentale de la structure visuelle type du CV (Lecture en F : En-tête $\rightarrow$ Colonne gauche $\rightarrow$ Contenu droit).}
\end{popfigure}

\begin{attention}[Erreurs de mise en page fréquentes (Terminale Tech)]
\begin{itemize}
    \item Photo : selfie, filtre Snapchat, fond coloré, lunettes de soleil, cadrage mauvais.
    \item Email : \texttt{kingboy99@..., princess\_love@..., footfan@...} $\rightarrow$ \textbf{Interdit}. Créez \texttt{prenom.nom@gmail.com}.
    \item Police fantaisie (Comic Sans, Chiller), taille 9 ou 14, 3 polices différentes.
    \item Blocs de texte justifiés sans puces (pavé illisible).
    \item Fichier nommé \texttt{CV\_final\_v2\_modif\_ok.pdf}.
    \item CV sur 2 pages pour un élève de Terminale (sauf PFE exceptionnel + portfolio joint).
\end{itemize}
\end{attention}

\coursec{Langue française : Valeurs du subjonctif et du conditionnel}

\courssub{Le conditionnel : politesse, hypothèse, atténuation}

\begin{propriete}[Valeurs du conditionnel (Présent / Passé) en contexte professionnel]
\begin{itemize}
    \item \textbf{Politesse / Atténuation :} Adoucir une demande, un avis, un refus. \emph{« Je \textbf{souhaiterais} postuler... », « Pourriez-vous m'indiquer... », « Cela \textbf{me semblerait} pertinent... »}.
    \item \textbf{Hypothèse / Condition :} Fait conditionné (souvent avec \emph{si} + imparfait). \emph{« Si j'étais retenu, je \textbf{m'investirais} pleinement. »}
    \item \textbf{Information rapportée / Prudence journalistique :} \emph{« L'entreprise \textbf{recruterait} 10 techniciens cette année. »} (Non confirmé).
    \item \textbf{Souhait / Désir :} \emph{« J'\textbf{aimerais} intégrer votre service maintenance. »}
\end{itemize}
\end{propriete}

\courssub{Le subjonctif : subjectivité, nécessité, volonté, doute}

\begin{propriete}[Valeurs du subjonctif (Présent / Passé) en contexte professionnel]
Le subjonctif s'emploie généralement dans une \textbf{subordonnée complétive} introduite par \emph{que}, après un verbe ou une expression marquant la subjectivité.
\begin{center}
\begin{tabularx}{\linewidth}{|l|X|l|}\hline
\rowcolor{popblueL}
\textbf{Valeur} & \textbf{Structure déclenchante (Exemples)} & \textbf{Exemple entretien / CV} \\ \hline
\textbf{Volonté / Souhait / Ordre} & \emph{Vouloir que, souhaiter que, demander que, exiger que, préférer que, il faut que, il est indispensable que...} & « Je \textbf{veux que} ce stage \textbf{me permette} d'acquérir... » / « Il \textbf{faut que} je \textbf{maîtrise} la GMAO. » \\ \hline
\textbf{Nécessité / Obligation} & \emph{Il faut que, il est nécessaire que, il est impératif que, il convient que...} & « Il \textbf{est impératif que} le candidat \textbf{soit} rigoureux sur la sécurité. » \\ \hline
\textbf{Doute / Incertitude / Négation} & \emph{Douter que, nier que, ne pas croire que, ne pas être sûr que, il est douteux que...} & « Je \textbf{ne suis pas sûr qu'il} \textbf{y ait} une place en maintenance prédictive. » \\ \hline
\textbf{Émotion / Sentiment / Jugement} & \emph{Être content que, regretter que, craindre que, apprécier que, il est normal que, il est dommage que...} & « J'\textbf{apprécie que} l'entreprise \textbf{forme} ses stagiaires. » / « Je \textbf{regrette de ne pas avoir} \textbf{pu} faire plus. » \\ \hline
\textbf{Finalité / But} & \emph{Afin que, pour que, de peur que, de sorte que...} & « Je prends des notes \textbf{afin que} je \textbf{n'oublie} rien. » \\ \hline
\textbf{Concession} & \emph{Bien que, quoique, même si + subjonctif...} & « \textbf{Bien que} je \textbf{sois} débutant, je \textbf{suis} opérationnel vite. » \\ \hline
\textbf{Superlatif / Adjectif + que} & \emph{Le seul / le premier / le meilleur / le dernier + que...} & « C'est le \textbf{seul} stage qui \textbf{m'ait permis} de voir de la HT. » \\ \hline
\end{tabularx}
\end{center}
\end{propriete}

\begin{exemplebox}[Exercice d'application : Choix du mode (Conditionnel / Subjonctif / Indicatif)]
Mettez le verbe entre parenthèses au mode et au temps convenables.
\begin{enumerate}
    \item Il est indispensable que le candidat \_\_\_\_\_\_\_\_ (connaître) les normes ISO 9001.
    \item Je \_\_\_\_\_\_\_\_ (aimer) savoir si le poste est toujours vacant.
    \item Bien qu'il \_\_\_\_\_\_\_\_ (être) jeune, il a déjà géré une équipe.
    \item Si j'avais le temps, je \_\_\_\_\_\_\_\_ (suivre) une formation en anglais technique.
    \item Le directeur demande que le rapport \_\_\_\_\_\_\_\_ (être) rendu vendredi.
    \item Je ne pense pas qu'il \_\_\_\_\_\_\_\_ (y avoir) beaucoup de candidats qualifiés.
\end{enumerate}
\textbf{Corrigé :}
\begin{enumerate}
    \item \textbf{connaisse} (Subjonctif présent — Nécessité / Il est indispensable que)
    \item \textbf{aimerais} (Conditionnel présent — Politesse / Souhait)
    \item \textbf{soit} (Subjonctif présent — Concession / Bien que)
    \item \textbf{suivrais} (Conditionnel présent — Hypothèse / Si + imparfait)
    \item \textbf{soit} (Subjonctif présent — Volonté / Demander que)
    \item \textbf{y ait} (Subjonctif présent — Doute / Ne pas penser que)
\end{enumerate}
\end{exemplebox}

\coursec{Production guidée : rédiger son CV ciblé}

\courssub{De l'analyse d'offre au CV final : démarche pas à pas}

\begin{methode}[Procédé de rédaction en 5 étapes]
\begin{enumerate}
    \item \textbf{Analyser l'offre (Mots-clés) :} Surligner : Compétences techniques requises, Soft skills, Diplômes, Verbes d'action, Valeurs entreprise. \emph{Ex : « Maintenance préventive HT/BT », « GMAO », « Rigueur », « Esprit d'équipe », « Bac F4 », « Dynamique »}.
    \item \textbf{Faire l'inventaire « Master » :} Lister TOUTES ses expériences, formations, compétences, projets, outils, langues, centres d'intérêt. Ne rien censurer.
    \item \textbf{Sélectionner \& Hiérarchiser (Matching) :} Ne garder que ce qui \emph{répond} à l'offre. Ordonner par pertinence (Expérience avant Formation si stage pertinent, sinon Formation en premier).
    \item \textbf{Rédiger les contenus (STAR + Verbes d'action) :} Transformer « J'ai fait de la maintenance » en « \textbf{Réalisé} la maintenance préventive de 5 postes HT/BT (GMAO), \textbf{réduit} les arrêts de 10\% ». Verbes : \emph{Conçu, Programmé, Diagnostiqué, Optimisé, Coordonné, Formé, Géré, Résolu}.
    \item \textbf{Mettre en page \& Relire (Zéro faute) :} Appliquer règles visuelles (Section 3). Relire à l'envers (chasse aux fautes). Faire relire par un pair / prof. Exporter PDF.
\end{enumerate}
\end{methode}

\begin{exoresolu}[Adaptation CV : Offre $\rightarrow$ Rubrique Compétences]
\textbf{Énoncé :} Offre : « Technicien Maintenance F4. Missions : Maintenance préventive/curative lignes automatisées (API Siemens S7-1200), GMAO, Amélioration continue. Profil : Bac F4, Rigueur, Force de proposition. »
\tcblower
\textbf{Compétences Techniques (Version Ciblée) :}
\begin{itemize}
    \item \textbf{Automatismes / API :} Siemens S7-1200 (TIA Portal) — \emph{Notions (TP/PFE)} ; Grafcet, Ladder — \emph{Maîtrise}.
    \item \textbf{Maintenance Industrielle :} Préventive (Gammes, GMAO \emph{MAINTI4}) / Curative (Diagnostique pannes capteurs/actionneurs, lecture schémas électrique/pneumatique) — \emph{Bonne pratique (Stage 2 mois)}.
    \item \textbf{Électrotechnique :} Puissance (Transformateurs, Moteurs asynchrones, Variateurs) / Contrôle-commande — \emph{Maîtrise (Cours + TP)}.
    \item \textbf{Amélioration Continue :} Démarche 5S (Stage), Analyse causes racines (5P / Ishikawa) — \emph{Notions}.
\end{itemize}
\textbf{Compétences Transverses :} Rigueur (Zéro erreur GMAO stage), Force de proposition (PFE : optimisation changement revêtements -15\% tps arrêt), Esprit sécurité (Habilitation électrique B2V essais, consignes verrouillage).
\end{exoresolu}

\coursec{Présenter son CV lors d'un entretien d'embauche}

\courssub{L'entretien d'embauche : définition, enjeux et déroulement type}

\begin{definition}[Entretien d'embauche]
L'entretien d'embauche est un \textbf{échange oral structuré} entre un candidat et un recruteur (ou un jury) visant à \textbf{vérifier l'adéquation} entre le profil du candidat (compétences, savoir-être, motivation) et les exigences du poste. C'est une situation de communication professionnelle à part entière, codifiée, où l'enjeu est double : \emph{valider les compétences techniques} (savoir-faire) et \emph{évaluer le potentiel relationnel} (savoir-être).
\end{definition}

Un entretien suit généralement un scénario en quatre phases. Le schéma ci-dessous en synthétise la dynamique temporelle.

\begin{popfigure}
\begin{tikzpicture}[
    node distance=1.2cm and 2.5cm,
    phase/.style={rectangle, rounded corners, draw=popblue, thick, fill=popblueL, minimum width=3.5cm, minimum height=1.8cm, text width=3.2cm, align=center, font=\small},
    arrow/.style={-Stealth, thick, popdark},
    label/.style={font=\footnotesize, text width=3cm, align=center, popmuted}
]
% Nodes
\node[phase, align=center] (accueil){\textbf{1. Accueil \& Mise en condition}\\(2--3 min)\\\textit{Installation, civilités, présentation du jury}};
\node[phase, right=of accueil, align=center] (presentation){\textbf{2. Présentation du candidat}\\(3--5 min)\\\textit{Pitch 2 min + questions de clarification}};
\node[phase, right=of presentation, align=center] (echange){\textbf{3. Échange approfondi}\\(10--15 min)\\\textit{Compétences, motivation, situations, défauts/qualités}};
\node[phase, right=of echange, align=center] (conclusion){\textbf{4. Conclusion \& Perspectives}\\(2--3 min)\\\textit{Questions du candidat, délais, au revoir}};
% Arrows
\draw[arrow] (accueil.east) -- node[above, label] {Transition fluide} (presentation.west);
\draw[arrow] (presentation.east) -- node[above, label] {Rebond sur le CV} (echange.west);
\draw[arrow] (echange.east) -- node[above, label] {Clôture proactive} (conclusion.west);
% Details under nodes
\node[label, below=0.3cm of accueil] {Objectif : \emph{Créer la confiance}};
\node[label, below=0.3cm of presentation] {Objectif : \emph{Accrocher l'attention}};
\node[label, below=0.3cm of echange] {Objectif : \emph{Prouver l'adéquation}};
\node[label, below=0.3cm of conclusion] {Objectif : \emph{Laisser une trace positive}};
\end{tikzpicture}
\legende{Figure 2 : Les 4 phases types d'un entretien d'embauche et leurs objectifs stratégiques.}
\end{popfigure}

\begin{aretenir}[Enjeux de l'entretien pour le candidat]
\begin{itemize}
    \item \textbf{Valider le CV :} Prouver que la réalité correspond au papier (cohérence).
    \item \textbf{Valoriser le potentiel :} Montrer ce que le CV ne dit pas (réactivité, aisance relationnelle, intelligence de situation).
    \item \textbf{Démontrer la motivation :} Elle se lit dans la préparation, les questions posées, l'énergie déployée.
    \item \textbf{Négocier implicitement :} Salaire, évolution, conditions de travail (souvent en fin de processus).
\end{itemize}
\end{aretenir}

\courssub{La présentation de soi en 2 minutes : la méthode « Présent $\rightarrow$ Passé $\rightarrow$ Futur »}

C'est le moment clé (phase 2). On ne vous demande pas de \emph{réciter} votre CV, mais de le \emph{raconter} avec une logique narrative.

\begin{methode}[Se présenter en 2 minutes (Structure en 3 temps)]
\begin{enumerate}
    \item \textbf{Le Présent — « Qui je suis » (30 s) :} Identité, formation actuelle/dernière, spécialité technique, statut (étudiant, demandeur d'emploi, salarié en reconversion). \emph{Exemple : « Bonjour, je m'appelle Marie Ngo, je suis en Terminale Technique série F4 (Électrotechnique) au Lycée Technique de Douala-Bassa. »}
    \item \textbf{Le Passé — « Ce que j'ai fait / Ce que je sais faire » (60 s) :} Sélectionnez \textbf{2 à 3 expériences ou projets phares} (stages, travaux pratiques, projets de fin d'études, jobs d'été) en lien direct avec le poste. Utilisez la méthode \textbf{STAR} (Situation, Tâche, Action, Résultat) pour chacun. \emph{Exemple : « Lors de mon stage de 2 mois à ENEO, j'ai participé à la maintenance préventive des postes HT/BT (Situation/Tâche). J'ai notamment diagnostiqué un défaut d'isolement sur un transformateur 30 kVA en appliquant la méthode de mesure de résistance de terre (Action), ce qui a permis d'éviter une coupure sur le quartier (Résultat). »}
    \item \textbf{Le Futur — « Ce que je veux faire / Ma motivation » (30 s) :} Pourquoi \emph{ce} poste, \emph{cette} entreprise, \emph{maintenant} ? Projetez-vous. \emph{Exemple : « Intégrer votre équipe maintenance me permettrait de mettre ma rigueur technique au service de la continuité de service, tout en développant mon expertise sur les énergies renouvelables, axe stratégique de votre structure. »}
\end{enumerate}
\end{methode}

\begin{attention}[Piège majeur : Réciter son CV]
Ne lisez pas votre CV ! Le jury l'a sous les yeux. \textbf{Commentez-le, hiérarchisez-le, donnez du sens.} Le CV est la \emph{partition} ; l'entretien est l'\emph{interprétation}. Si vous dites : « En 2022 j'ai fait ci, en 2023 j'ai fait ça », vous êtes un robot. Si vous dites : « Mon stage chez X a confirmé mon goût pour la maintenance curative, c'est pourquoi j'ai choisi mon projet de fin d'études sur... », vous êtes un \emph{futur professionnel}.
\end{attention}

\begin{exoresolu}[Exemple de pitch 2 minutes (Candidat Terminale F3 - Maintenance Industrielle)]
\textbf{Énoncé :} Rédiger le script oral d'un élève de Terminale F3 postulant pour un stage de fin d'études dans une cimenterie.

\tcblower
\textbf{Solution (Script annoté) :}
\begin{itemize}
    \item \textbf{[Présent]} « Bonjour Monsieur le Directeur, Mesdames, Messieurs. Je m'appelle Paul Biya, j'ai 18 ans. Je suis en Terminale F3 Maintenance des Systèmes Automatisés au Lycée Technique Industriel de Bafoussam. Je suis titulaire du Probatoire mention Assez Bien. » $\leftarrow$ \textit{Identité claire, niveau validé.}
    \item \textbf{[Passé - Expérience 1]} « L'an dernier, pendant mon stage de 4 semaines à la SOCIM (Société Cimentière), j'ai été affecté à l'atelier mécanique. Ma mission : assister le chef d'équipe sur la maintenance préventive des broyeurs à boulets. Concrètement, j'ai relevé les jeux de paliers, vérifié l'état des revêtements et participé au graissage programmé. J'ai appris à lire les gammes de maintenance et à remplir les fiches d'intervention GMAO. » $\leftarrow$ \textit{Contexte technique précis, vocabulaire métier (GMAO, gamme, palier).}
    \item \textbf{[Passé - Projet PFE]} « Cette expérience a nourri mon Projet de Fin d'Études : « Optimisation du changement des revêtements du broyeur N°2 ». J'ai modélisé la procédure sous MS Project et proposé une réorganisation qui réduit le temps d'arrêt de 15\%. » $\leftarrow$ \textit{Preuve d'autonomie, d'esprit d'amélioration, outil bureautique.}
    \item \textbf{[Futur]} « Revenir chez vous pour mon stage de fin d'études (8 semaines) serait la suite logique. Je connais vos normes sécurité, votre parc machine, et je souhaite approfondir la maintenance prédictive (analyse vibratoire) que vous développez. Je suis disponible, rigoureux sur la sécurité (port EPI, consignes verrouillage) et prêt à m'investir pleinement. » $\leftarrow$ \textit{Connaissance de l'entreprise, compétence transversale (sécurité), disponibilité.}
\end{itemize}
\textbf{Durée estimée :} 1 min 50 s. \textbf{Ton :} Posé, regard franc, débit modéré.
\end{exoresolu}

\courssub{Communication verbale : clarté, vocabulaire et modes (Conditionnel / Subjonctif)}

\begin{propriete}[Registre de langue professionnel : indicateurs linguistiques]
L'entretien exige le \textbf{registre soutenu / formel} adapté au contexte technique.
\begin{itemize}
    \item \textbf{Clarté \& Structure :} Phrases courtes, connecteurs logiques (\emph{donc, par conséquent, en effet, toutefois, par ailleurs, afin que}).
    \item \textbf{Vocabulaire métier :} Utilisez les termes techniques exacts (ex. : \emph{« plan de maintenance préventive », « schéma unifilaire », « tolérance dimensionnelle », « norme ISO 9001 »} et non « les papiers pour réparer »).
    \item \textbf{Conditionnel de politesse :} Indispensable pour formuler des demandes, des souhaits, des hypothèses sans agressivité.
    \item \textbf{Subjonctif de subjectivité :} Marque l'engagement, la nécessité, la volonté. \emph{« Il faut que je \textbf{maîtrise}... », « Je veux que ce stage \textbf{m'apporte}... », « Bien que je \textbf{sois} débutant... »}.
\end{itemize}
\end{propriete}

\begin{center}
\begin{tabularx}{\linewidth}{|l|X|X|}\hline
\rowcolor{popblueL}
\textbf{Fonction} & \textbf{Indicatif (Trop direct / Familier)} & \textbf{Conditionnel / Subjonctif (Professionnel)} \\ \hline
Demander une précision & « C'est quoi le salaire ? » / « Je veux savoir... » & « \textbf{Pourrais}-je connaître la grille salariale ? » / « \textbf{Serait}-il possible de préciser... » \\ \hline
Exprimer son souhait & « Je veux ce poste. » / « J'aimerais bien... » & « Je \textbf{souhaiterais} intégrer votre équipe. » / « Ce poste \textbf{m'intéresserait} vivement. » \\ \hline
Proposer / S'engager & « Je fais ça. » / « Je peux le faire. » & « Je \textbf{pourrais} m'occuper de... » / « Je \textbf{serais} ravi de contribuer à... » \\ \hline
Émettre une réserve & « Je ne sais pas faire ça. » & « Cela \textbf{nécessiterait} une formation que je \textbf{suivrais} avec motivation. » \\ \hline
Exprimer une nécessité & « Il faut que je pars. » / « Je dois y aller. » & « Il \textbf{est indispensable que} je \textbf{maîtrise} cet outil. » (Subjonctif) \\ \hline
Marquer la concession & « Même si je suis débutant... » & « \textbf{Bien que} je \textbf{sois} débutant, je suis opérationnel. » (Subjonctif) \\ \hline
\end{tabularx}
\end{center}

\begin{exemplebox}[Exercice d'application : Transposition aux modes professionnels]
Transformez ces phrases « brutes » en formulations professionnelles (Conditionnel ou Subjonctif selon le sens).
\begin{enumerate}
    \item Je veux faire ce stage chez vous.
    \item Je ne connais pas ce logiciel.
    \item Quand est-ce que je commence ?
    \item Je suis le meilleur pour ce poste.
    \item Il faut que j'apprends l'anglais technique.
    \item Je ne suis pas sûr qu'il y a une place.
\end{enumerate}
\textbf{Corrigé :}
\begin{enumerate}
    \item Je \textbf{souhaiterais} / \textbf{aimerais} effectuer ce stage au sein de votre structure. (Conditionnel - Politesse)
    \item Je ne \textbf{maîtrise} pas encore ce logiciel, mais je \textbf{l'apprendrais} rapidement / une formation me \textbf{permettrait} de l'appréhender. (Conditionnel - Hypothèse/Politesse)
    \item \textbf{Pourriez}-vous m'indiquer la date prévisionnelle de prise de fonction ? (Conditionnel - Politesse)
    \item Je \textbf{pense} que mon profil \textbf{correspondrait} bien aux attentes du poste. (Conditionnel - Atténuation)
    \item Il \textbf{est indispensable que} j'\textbf{apprenne} l'anglais technique. (Subjonctif - Nécessité)
    \item Je \textbf{ne suis pas sûr qu'il y ait} une place. (Subjonctif - Doute)
\end{enumerate}
\end{exemplebox}

\courssub{Communication non verbale : le langage du corps}

Les travaux en sciences de la communication (notamment A. Mehrabian) soulignent l'impact majeur du non-verbal.

\begin{exemplebox}[Exemple chiffré : La règle des 3 V (Mehrabian)]
Dans la communication des sentiments et attitudes (cohérence message/émotion) :
\begin{itemize}
    \item \textbf{Verbal (mots) :} 7 \% de l'impact.
    \item \textbf{Vocal (voix, débit, ton) :} 38 \% de l'impact.
    \item \textbf{Visuel (non verbal) :} 55 \% de l'impact.
\end{itemize}
\emph{Conséquence :} Si votre discours dit « Je suis motivé » mais que votre regard fuit, vos épaules sont affaissées et votre voix monocorde, le jury croira \textbf{votre corps}, pas vos mots. La cohérence \textbf{Verbal / Vocal / Visuel} est la clé de la crédibilité.
\end{exemplebox}

\begin{methode}[Check-list Non Verbale (Avant / Pendant / Après)]
\begin{enumerate}
    \item \textbf{Tenue vestimentaire :} Adaptée au secteur (Bleu de travail propre pour atelier / Costume ou chemise-pantalon soigné pour bureau/encadrement). \textbf{Règle :} Mieux vaut être \emph{légèrement trop habillé} que pas assez. Chaussures cirées, cheveux coiffés, ongles nets.
    \item \textbf{Arrivée / Poignée de main :} Arrivée 10-15 min en avance. Poignée de main \textbf{ferme, brève, sèche}, regard dans les yeux, sourire léger. Ni « main molle » (manque d'assurance), ni « broyeur » (agressivité).
    \item \textbf{Posture assise :} Dos droit contre le dossier, jambes droites ou croisées discrètement (pas jambe sur genou), mains visibles sur la table ou jointes. Pas de bras croisés (fermeture).
    \item \textbf{Regard :} Contact visuel \textbf{franc et mobile} (regarder chaque membre du jury, pas seulement celui qui a posé la question). Ne pas fixer le sol, le plafond ou ses chaussures.
    \item \textbf{Gestuelle :} Gestes \emph{illustrateurs} (main ouverte pour donner, compter sur les doigts pour structurer : « Premier point... Second point... »). Éviter gestes \emph{parasites} : toucher cheveux, stylo qui clique, pied qui bat, bras ballants.
    \item \textbf{Voix :} Volume audible (ni chuchotement, ni cri), débit \textbf{ralenti} (l'oral demande du temps de traitement), articulation soignée, intonation \textbf{variée} (monotonie = ennui).
    \item \textbf{Départ :} Remercier, poignée de main de sortie, regard, « Au revoir Monsieur/Madame », sortir calmement.
\end{enumerate}
\end{methode}

\begin{savaistu}[Ponctualité au Cameroun]
Arriver \textbf{15 minutes en avance} à un entretien est une marque de respect et de sérieux très appréciée des recruteurs camerounais. Cela laisse le temps de se présenter à l'accueil, de souffler, d'observer l'environnement (affichage sécurité, ambiance) et de relire ses notes. Arriver « à l'heure pile » (ou en retard) est perçu comme un manque d'organisation ou de motivation. \emph{Anticipez les embouteillages (Douala, Yaoundé) ou les pannes de transport !}
\end{savaistu}

\courssub{Questions fréquentes du jury : anticiper, structurer, illustrer}

Le jury teste la cohérence CV/Discours, la personnalité et la motivation. Voici les « grands classiques » et la stratégie de réponse.

\begin{center}
\begin{tabularx}{\linewidth}{|l|X|X|}\hline
\rowcolor{popblueL}
\textbf{Question type} & \textbf{Intention du jury} & \textbf{Stratégie de réponse (Structure)} \\ \hline
« Présentez-vous. » & Vérifier synthèse, aisance, fil conducteur. & \textbf{Méthode 3 temps (Présent/Passé/Futur)} vue plus haut. \emph{Ne pas raconter sa vie privée.} \\ \hline
« Pourquoi ce poste / notre entreprise ? » & Tester la motivation \emph{réelle} et la préparation. & \textbf{Entreprise} (valeurs, projets, réputation) + \textbf{Poste} (missions techniques, défis) + \textbf{Vous} (adéquation compétences/valeurs). \emph{Citez un fait précis sur l'entreprise.} \\ \hline
« Quelles sont vos 3 qualités ? » & Vérifier auto-évaluation et adéquation poste. & Choisir \textbf{qualités techniques + transverses}. \textbf{Illustrer chacune par un fait concret (STAR)}. Ex : « Rigueur (stage : 0 erreur sur fiches GMAO), Esprit d'équipe (projet PFE en binôme), Adaptabilité (stage en milieu inconnu). » \\ \hline
« Quels sont vos 3 défauts ? » & Tester lucidité, capacité d'amélioration, honnêteté. \textbf{Pas de faux défauts} (« Je suis trop perfectionniste »). & \textbf{Vrai défaut mineur, non rédhibitoire + Action de correction.}
Ex : « Je parle peu en public $\rightarrow$ Je me suis inscrit au club débat / je force la parole en TP. » / « Anglais technique perfectible $\rightarrow$ Je suis une formation en ligne / je lis des normes en anglais. » \\ \hline
« Où vous voyez-vous dans 5 ans ? » & Tester projection, fidélité potentielle, ambition réaliste. & \textbf{Évolution technique cohérente} : « Technicien confirmé $\rightarrow$ Chef d'équipe / Spécialiste maintenance prédictive / Formateur interne. » Éviter : « Patron », « À l'étranger », « Je ne sais pas ». \\ \hline
« Avez-vous des questions ? » & Tester intérêt, curiosité, maturité. & \textbf{TOUJOURS OUI.} 2-3 questions préparées : « Comment se passe l'intégration des stagiaires ? », « Quels sont les défis techniques actuels du service ? », « Quelle est la place de la formation continue ? ». \textbf{Pas de questions sur salaire/congé au 1er entretien.} \\ \hline
\end{tabularx}
\end{center}

\begin{exoresolu}[Répondre à « Quels sont vos défauts ? » — Cas concret]
\textbf{Énoncé :} Un élève de Terminale G2 (Génie Civil) répond : « Je suis têtu. »
\tcblower
\textbf{Analyse :} « Têtu » a une connotation négative (rigidité, refus d'écoute). Risque de rejet.
\textbf{Reformulation positive (Assertivité / Persévérance) :}
\begin{quote}
« Je peux faire preuve de \textbf{persévérance} (qualité) : quand je bloque sur un problème technique (ex. calcul de ferraillage), je cherche la solution jusqu'au bout. Parfois, cela m'a fait perdre du temps (défaut associé : \emph{difficulté à lâcher prise / demander de l'aide assez tôt}). Aujourd'hui, j'ai appris à \textbf{chronométrer mes recherches} (ex: 30 min max seul) et à \textbf{solliciter mon tuteur ou un pair} si je n'ai pas trouvé. Cela m'a fait gagner en efficacité lors de mon dernier stage. »
\end{quote}
\textbf{Bilan :} Le « défaut » devient une \emph{zone de progression identifiée et gérée}. C'est ce que le jury attend.
\end{exoresolu}

\courssub{Gestion du stress : préparation, respiration, simulation}

Le stress est normal (signe d'enjeu). Il devient handicapant s'il n'est pas géré.

\begin{methode}[Protocol anti-stress (J-1 / Jour J / Pendant)]
\begin{enumerate}
    \item \textbf{J-1 (Préparation cognitive) :} Relire l'offre, son CV, la fiche entreprise. Préparer ses 3 questions. Choisir/préparer sa tenue. Visualiser l'entretien qui se \emph{bien passe} (imagerie mentale positive). Coucher tôt.
    \item \textbf{Jour J (Physiologique) :} Petit déjeuner protéiné (évite hypoglycémie). Éviter caféine/excès sucre. Arriver \textbf{15 min avant}. Aux toilettes : \textbf{Respiration carrée} (Inspir 4 s / Bloque 4 s / Expire 4 s / Bloque 4 s $\times$ 4 cycles) $\rightarrow$ Baisse cortisol, recentrage.
    \item \textbf{Pendant (Ancrage) :} Si blanc / trou de mémoire : \textbf{Buvez une gorgée d'eau} (temps de pause légitime). Dites : « C'est une excellente question, permettez-moi de réfléchir deux secondes. » $\rightarrow$ Montre maîtrise, pas faiblesse. Posez les pieds au sol (ancrage physique), mains à plat sur la table.
\end{enumerate}
\end{methode}

\begin{attention}[Erreurs fatales de gestion du stress]
\begin{itemize}
    \item \textbf{Mensonge / Invention :} Le jury vérifie (tests techniques, appel tuteur stage). Perte de crédibilité totale.
    \item \textbf{Critiquer ancien employeur / tuteur / école :} Signe d'immaturité, risque conflictuel.
    \item \textbf{Monopoliser la parole / Couper la parole :} Manque d'écoute, manque de respect.
    \item \textbf{Téléphone qui sonne / vibre :} Éteignez-le \textbf{avant} d'entrer dans le bâtiment. Pas de mode vibreur dans la poche (bruité distrayant).
\end{itemize}
\end{attention}

\coursec{Intégration, compte rendu et remédiation}

\courssub{TP / Activité d'intégration : Simulation d'entretien devant jury de classe}

C'est l'activité sommative de l'unité. Elle mobilise \textbf{tous les savoirs} : CV rédigé, pitch oral, langue (conditionnel/subjonctif), non-verbal, argumentation.

\begin{experience}[Simulation d'entretien d'embauche — Dispositif complet]
\textbf{Objectif :} Produire et présenter un CV adapté à une offre réelle/fictive, face à un jury, en situation quasi-réelle.
\textbf{Durée :} 2 séances (Préparation 1h + Passages 1h30).
\textbf{Matériel :} Offres d'emploi imprimées (ou projetées), CV élèves imprimés (3 ex : jury x2 + candidat), Grilles d'évaluation (ci-dessous), Chronomètre, Salle aménagée (bureau jury, chaise candidat, eau).

\textbf{Protocole :}
\begin{enumerate}
    \item \textbf{Phase 1 - Préparation (Individuel + Pair) :} Chaque élève choisit une offre. Révise son CV en fonction. Répète son pitch 2 min avec un pair (feedback croisé : « As-tu illustré ? Conditionnel/Subjonctif ? Regard ? »).
    \item \textbf{Phase 2 - Constitution des Jurys :} Groupes de 3-4 élèves. Rôles tournants : \textbf{Président} (anime, gère temps), \textbf{Technicien} (questions métier), \textbf{RH/Comportement} (soft skills, motivation), \textbf{Observateur} (remplit grille, ne parle pas). Le professeur circule, note, ne note \emph{pas} les élèves (évaluation formative).
    \item \textbf{Phase 3 - Passage (15 min / candidat) :}
        \begin{itemize}
            \item Accueil / Installation (1 min).
            \item Pitch 2 min (chronométré).
            \item Questions Jury (8-10 min) : Basées sur CV + questions types.
            \item Questions Candidat (2 min).
            \item Conclusion / Débrief immédiat oral (2 min) : Jury donne 1 point fort / 1 axe progression.
        \end{itemize}
    \item \textbf{Phase 4 - Compte rendu écrit (Devoir) :} Chaque élève rédige une \textbf{fiche d'auto-évaluation} (1 page) : « Ce que j'ai réussi / Ce que je rate / Ce que je change pour le vrai entretien ». Remis au prof pour remédiation individuelle.
\end{enumerate}
\end{experience}

\begin{popfigure}
\begin{center}
\begin{tabularx}{\linewidth}{|>{\bfseries}l|X|c|c|c|}\hline
\rowcolor{popblueL}
\textbf{Critère} & \textbf{Indicateurs observables (Niveau Terminale Tech)} & \textbf{/10} & \textbf{/5} & \textbf{/5} \\ \hline\hline
\multicolumn{5}{|c|}{\cellcolor{popblueL}\textbf{A. CONTENU \& FOND (sur 10)}} \\ \hline
Adéquation Profil/Poste & Compétences techniques (stages, PFE, TP) explicitées et liées aux missions de l'offre. & & & \\ \hline
Argumentation / Illustration & Méthode STAR utilisée. Exemples concrets, chiffrés, techniques. Pas de généralités. & & & \\ \hline
Motivation & Connaissance entreprise/secteur. Projet professionnel cohérent. Questions pertinentes posées. & & & \\ \hline
\multicolumn{5}{|c|}{\cellcolor{poporangeL}\textbf{B. FORME \& LANGUE (sur 5)}} \\ \hline
Expression orale & Clarté, débit maîtrisé, articulation, vocabulaire technique précis, phrases complètes. & & & \\ \hline
Registre / Politesse & Conditionnel/Subjonctif de politesse/nécessité systématiques. Tutoiement absent. Reformulations correctes. & & & \\ \hline
Structure discours & Transitions logiques. Respect temps (Pitch 2 min $\pm$ 30s). Synthèse efficace. & & & \\ \hline
\multicolumn{5}{|c|}{\cellcolor{popgreenL}\textbf{C. COMPORTEMENT \& NON VERBAL (sur 5)}} \\ \hline
Présence / Posture & Tenue pro. Ponctualité. Poignée de main. Assis droit. Mains visibles. Pas gestes parasites. & & & \\ \hline
Regard / Écoute & Contact visuel distribué (tout jury). Écoute active (hochement, note, ne coupe pas). & & & \\ \hline
Gestion émotionnelle & Calme apparent. Résistance pression (questions pièges). Rebond si erreur. Sourire. & & & \\ \hline\hline
\multicolumn{2}{|r|}{\textbf{TOTAL}} & \textbf{/10} & \textbf{/5} & \textbf{/5} \\ \hline
\multicolumn{2}{|r|}{\textbf{NOTE GLOBALE /20}} & \multicolumn{3}{c|}{} \\ \hline
\end{tabularx}
\end{center}
\legende{Grille d'évaluation orale : Fond (10) / Forme (5) / Comportement (5). Utilisée par l'observateur lors de la simulation.}
\end{popfigure}

\begin{aretenir}[Bilan de l'unité « CV \& Entretien »]
\begin{itemize}
    \item Le CV est un \textbf{objet négocié} : on l'adapte à chaque offre (mots-clés, mise en avant sélective).
    \item L'entretien est une \textbf{conversation professionnelle}, pas un interrogatoire. Vous évaluez l'entreprise autant qu'elle vous évalue.
    \item \textbf{Cohérence} est le maître-mot : CV $\leftrightarrow$ Pitch $\leftrightarrow$ Réponses $\leftrightarrow$ Non-verbal $\leftrightarrow$ Projet pro.
    \item La \textbf{préparation} (recherche entreprise, répétition pitch, anticipation questions, logistique) représente 80\% de la réussite.
    \item Le \textbf{conditionnel de politesse}, le \textbf{subjonctif de nécessité/volonté} et le \textbf{vocabulaire technique précis} sont les marqueurs linguistiques de votre professionnalisme naissant.
\end{itemize}
\end{aretenir}

\begin{exercice}[Entraînement individuel : « Mon entretien idéal »]
\begin{enumerate}
    \item Choisissez une offre réelle (site emploi.cm, LinkedIn, journal) correspondant à votre série. Imprimez-la.
    \item Adaptez votre CV (version « Master » $\rightarrow$ version « Cible »).
    \item Rédigez votre \textbf{script complet du pitch 2 minutes} (méthode 3 temps + STAR). Chronométrez-vous à voix haute.
    \item Préparez \textbf{3 questions} à poser au recruteur (pas salaire, pas congés).
    \item Identifiez \textbf{1 défaut réel} et sa \textbf{parade concrète} (méthode : Défaut $\rightarrow$ Action correction).
    \item Simulez seul face à votre webcam (enregistrement). Re-regardez : Posture ? Regard ? Modes verbaux ? Temps ? Corrigez. Recommencez.
\end{enumerate}
\textit{Remettez : Offre + CV Cible + Script Pitch + 3 Questions + Fiche Défaut/Parade + Lien vidéo (optionnel).}
\end{exercice}

\begin{corrige}[Exercice n°1 - Attendus pédagogiques]
Pas de corrigé unique. Évaluation par le professeur sur grille simplifiée :
\begin{itemize}
    \item \textbf{Pertinence :} CV/Offre match (mots-clés présents).
    \item \textbf{Qualité script :} Structure 3 temps respectée, 2 exemples STAR, conditionnel/subjonctif utilisés, temps 1'30''-2'15''.
    \item \textbf{Questions :} Ouvertes, intelligentes, montrent recherche entreprise.
    \item \textbf{Défaut/Parade :} Authentique, non bloquant, action concrète datée.
    \item \textbf{Vidéo :} Progression visible entre prise 1 et prise finale (regard, fluidité, gestuelle).
\end{itemize}
\textbf{Remédiation type :} Élèves « lecture CV » $\rightarrow$ Exercice « Raconte ton stage à un ami en 1 min » (oral spontané) $\rightarrow$ Re-structuration. Élèves « monosyllabes » $\rightarrow$ Jeu « Interdit de dire Oui/Non, phrase complète obligatoire ». Élèves « stress extrême » $\rightarrow$ Passage en binôme « candidat / coach bienveillant » avant jury.
\end{corrige}

\coursec{Intégration, compte rendu et remédiation}

Après avoir simulé l'entretien d'embauche dans la section précédente, nous entrons dans la phase décisive : l'\emph{intégration}. C'est le moment où l'élève transforme l'essai, où les savoirs épars (structure du CV, codes iconiques, valeurs modales, posture orale) se soudent en une compétence unique : \textbf{se présenter avec efficacité et justesse devant un jury}. Cette séance de synthèse, de retour critique et de remédiation ciblée garantit que chaque candidat quitte la classe avec un CV abouti et une aisance orale validée.

\courssub{Synthèse de la leçon : les quatre piliers du CV réussi}

Avant tout retour individuel, posons la \cle{carte mentale} récapitulative. Elle structure la révision et sert de grille de lecture pour l'auto-évaluation finale.

\begin{popfigure}
\begin{tikzpicture}[
  mindmap,
  concept color=popblue!80,
  text=white,
  font=\small\bfseries,
  level 1/.style={level distance=4.5cm, sibling angle=90},
  level 2/.style={level distance=3cm, sibling angle=45},
  every node/.style={scale=0.85}
]
\node[concept, align=center]{Le CV\\en Terminale Tech}
  child[concept color=poporange] {
    node[concept, align=center]{CONTENU\\ (5 rubriques)}
    child { node[concept color=poporange!70, align=center]{1. Identité\\ et coordonnées} }
    child { node[concept color=poporange!70, align=center]{2. Titre\\ et projet pro.} }
    child { node[concept color=poporange!70, align=center]{3. Formation\\ et diplômes} }
    child { node[concept color=poporange!70, align=center]{4. Expériences\\ et stages} }
    child { node[concept color=poporange!70, align=center]{5. Compétences\\ et centres d'intérêt} }
  }
  child[concept color=popgreen] {
    node[concept, align=center]{FORME\\ (Mise en page)}
    child { node[concept color=popgreen!70, align=center]{Police unique\\ (Arial, Calibri)} }
    child { node[concept color=popgreen!70, align=center]{Icônes discrètes\\ (tél., mail, lieu)} }
    child { node[concept color=popgreen!70, align=center]{Alignement parfait\\ (tabulations)} }
    child { node[concept color=popgreen!70, align=center]{1 page max\\ (format PDF)} }
  }
  child[concept color=poppurple] {
    node[concept, align=center]{LANGUE\\ (Morphosyntaxe)}
    child { node[concept color=poppurple!70, align=center]{Subjonctif :\\ nécessité, souhait} }
    child { node[concept color=poppurple!70, align=center]{Conditionnel :\\ hypothèse, politesse} }
    child { node[concept color=poppurple!70, align=center]{Verbes d'action\\ (passé composé)} }
    child { node[concept color=poppurple!70, align=center]{Zéro faute\\ (relecture croisée)} }
  }
  child[concept color=poppink] {
    node[concept, align=center]{ORAL\\ (Entretien)}
    child { node[concept color=poppink!70, align=center]{Pitch 2 min\\ (qui, quoi, pourquoi)} }
    child { node[concept color=poppink!70, align=center]{Preuves STAR\\ (Situation, Tâche,\\ Action, Résultat)} }
    child { node[concept color=poppink!70, align=center]{Gestion du stress\\ (respiration, regard)} }
    child { node[concept color=poppink!70, align=center]{Questions finales\\ (intérêt, veille)} }
  };
\end{tikzpicture}
\legende{Carte mentale récapitulative : les quatre dimensions indissociables d'un CV performant en Terminale technique.}
\end{popfigure}

\begin{aretenir}[Les 5 rubriques obligatoires du CV technique]
\textbf{1. En-tête :} nom, prénom, adresse, téléphone, adresse électronique professionnelle, permis de conduire (si pertinent pour le poste). \textbf{2. Titre et projet :} intitulé du poste visé + phrase d'accroche (ex. \emph{« Technicien en maintenance industrielle, spécialité électricité, je recherche un stage de fin d'études pour mettre en œuvre mes compétences en automatisme »}). \textbf{3. Formation :} ordre antéchronologique (du plus récent au plus ancien). Diplôme, établissement, année, mention, options. \textbf{4. Expériences :} stages, jobs de vacances, bénévolat. Structure : dates | structure | poste | 3 à 4 puces \emph{action + résultat chiffré}. \textbf{5. Compétences :} techniques (logiciels, machines, normes), linguistiques (niveau), transversales (rigueur, travail en équipe), centres d'intérêt (sport collectif, engagement associatif).
\end{aretenir}

\begin{aretenir}[Règles d'or de la mise en page]
\begin{itemize}
  \item \textbf{Sobriété :} fond blanc, police sans empattement (Calibri, Arial), taille 10 à 11 points, interligne 1,15.
  \item \textbf{Hiérarchie visuelle :} rubriques en MAJUSCULES gras (taille 12 à 14), sous-rubriques en minuscules gras.
  \item \textbf{Alignement :} utiliser les tabulations (et non la barre d'espace) pour aligner dates, lieux et descriptions.
  \item \textbf{Icônes :} uniquement pour le contact, la localisation, les langues et l'informatique ; taille modeste, couleur unique (bleu foncé ou gris foncé).
  \item \textbf{Export :} \textbf{PDF uniquement} ; nom du fichier normé : \texttt{Nom\_Prenom\_PosteVise.pdf}.
\end{itemize}
\end{aretenir}

\begin{aretenir}[Valeurs du subjonctif et du conditionnel au service du CV]
\begin{center}
\begin{tabularx}{\linewidth}{|l|X|X|}\hline
\rowcolor{popblueL} \textbf{Mode} & \textbf{Valeur principale au CV / à l'entretien} & \textbf{Exemple d'application} \\ \hline
\textbf{Subjonctif} & Nécessité, obligation, souhait, doute, sentiment (après \emph{il faut que, il est indispensable que, je souhaite que}) & \emph{« Il faut que je \textbf{maîtrise} l'anglais technique. »} \emph{« Le directeur souhaite que je \textbf{sois} opérationnel dès lundi. »} \\ \hline
\textbf{Conditionnel} & Hypothèse, projet futur, politesse, atténuation (\emph{si} + imparfait $\rightarrow$ conditionnel dans la principale) & \emph{« Je \textbf{souhaiterais} intégrer votre service maintenance. »} \emph{« Si j'\textbf{obtenais} le poste, je \textbf{m'investirais} pleinement. »} \\ \hline
\end{tabularx}
\end{center}
\end{aretenir}

\courssub{Compte rendu des simulations d'entretien : analyse collective}

Les simulations menées en séance 6 ont révélé des tendances communes. Le jury (enseignant + pairs) a noté chaque candidat sur une grille de 20 points (5 points : contenu du CV ; 5 points : forme et orthographe ; 5 points : oral et argumentation ; 5 points : langue et registre). Voici la synthèse anonymisée.

\begin{exemplebox}[Profil type « Candidat A » — Note : 14/20]
\textbf{Points forts :} CV structuré, rubriques complètes, verbes d'action variés (\emph{concevoir, programmer, diagnostiquer}). Oral fluide, regard franc, argumentation STAR maîtrisée sur le stage en usine. Subjonctif correct : \emph{« Il est primordial que je m'adapte... »}.
\textbf{Axes d'amélioration :} une faute d'accord (\emph{« stages effectuées »} au lieu de \emph{« stages effectués »}) dans la rubrique Expériences. Icône de téléphone mal alignée. Conditionnel hésitant : \emph{« Je voudrais »} (trop direct) au lieu de \emph{« Je souhaiterais »}. Gestion du temps : pitch de 2 min 30 s (dépassement des 2 minutes prévues).
\end{exemplebox}

\begin{exemplebox}[Profil type « Candidat B » — Note : 9/20 — Seuil de remédiation]
\textbf{Points forts :} projet professionnel clair, centres d'intérêt pertinents (capitaine de l'équipe de football = preuve de leadership).
\textbf{Axes majeurs :} CV sur 1,5 page (marges trop larges). Rubrique « Compétences » vide. Orthographe : 4 fautes (accords, homophones). Oral : lecture du CV, absence de regard, voix monocorde. Subjonctif absent, conditionnel mal formé (\emph{« Je serais voulu »} au lieu de \emph{« J'aurais voulu »} ou mieux \emph{« Je souhaiterais »}). Aucune question posée au jury.
\end{exemplebox}

\begin{exemplebox}[Grille de remédiation — Seuils d'action]
\begin{center}
\begin{tabularx}{\linewidth}{|c|l|X|}\hline
\rowcolor{popblueL} \textbf{Note simulation} & \textbf{Statut} & \textbf{Actions obligatoires} \\ \hline
$\geq$ 14/20 & Validé & Affinage seul (relecture, polissage du PDF). \\ \hline
10 à 13,5/20 & Consolidation & Atelier ciblé (orthographe OU oral OU langue) + nouvelle simulation dans 15 jours. \\ \hline
$<$ 10/20 & \textbf{Remédiation intensive} & \textbf{Reprise individuelle complète :} réécriture du CV + 3 séances d'entraînement oral + exercices de conjugaison quotidiens. Nouvelle évaluation sous 3 semaines. \\ \hline
\end{tabularx}
\end{center}
\legende{Grille décisionnelle : tout élève sous la barre des 10/20 bénéficie d'un accompagnement individualisé obligatoire.}
\end{exemplebox}

\courssub{Remédiation ciblée : les valeurs du subjonctif et du conditionnel}

Les erreurs de morphosyntaxe (subjonctif/conditionnel) sont le marqueur le plus visible du niveau de langue. Nous travaillons ici sur la \cle{reconnaissance} et la \cle{production} en contexte professionnel.

\begin{methode}[Transformer une phrase selon l'intention : subjonctif ou conditionnel]
\emph{Objectif :} adapter le mode à la nuance voulue (obligation, souhait, hypothèse, politesse).
\begin{enumerate}
  \item \textbf{Identifier le verbe principal} et son sujet.
  \item \textbf{Repérer le déclencheur :} \emph{il faut que, je veux que, il est nécessaire que} $\rightarrow$ subjonctif ; \emph{si} + imparfait, \emph{je voudrais, je souhaiterais} $\rightarrow$ conditionnel.
  \item \textbf{Conjuguer :} au subjonctif présent, radical de la 3\up{e} personne du pluriel de l'indicatif présent + terminaisons \emph{-e, -es, -e, -ions, -iez, -ent} ; au conditionnel présent, radical du futur simple + terminaisons \emph{-ais, -ais, -ait, -ions, -iez, -aient}.
  \item \textbf{Vérifier} l'accord avec le sujet et la cohérence du sens.
\end{enumerate}
\end{methode}

\begin{exoresolu}[Entraînement de conjugaison : contexte entretien]
\textbf{Énoncé :} conjuguez les verbes entre parenthèses au mode et au temps adaptés (subjonctif présent ou conditionnel présent).
\begin{enumerate}
  \item Le recruteur exige que le candidat \_\_\_\_\_\_\_\_\_\_ (maîtriser) l'anglais technique.
  \item Je \_\_\_\_\_\_\_\_\_\_ (apprécier) de pouvoir vous présenter mon projet de vive voix.
  \item Il est indispensable que tu \_\_\_\_\_\_\_\_\_\_ (être) ponctuel le jour de l'entretien.
  \item Si j'\_\_\_\_\_\_\_\_\_\_ (obtenir) ce poste, je \_\_\_\_\_\_\_\_\_\_ (s'investir) pleinement dans la maintenance préventive.
  \item Nous voudrions que vous \_\_\_\_\_\_\_\_\_\_ (expliquer) votre expérience en automatisme.
\end{enumerate}
\tcblower
\textbf{Solution détaillée :}
\begin{enumerate}
  \item \textbf{maîtrise} (subjonctif présent — \emph{exiger que} + nécessité). Radical \emph{maîtris-} + terminaison \emph{-e}.
  \item \textbf{apprécierais} (conditionnel présent — politesse, souhait atténué). Radical du futur \emph{apprécier-} + \emph{-ais}.
  \item \textbf{sois} (subjonctif présent — \emph{il est indispensable que} + obligation). Verbe \emph{être}, irrégulier : \emph{que je sois, que tu sois, qu'il soit, que nous soyons, que vous soyez, qu'ils soient}.
  \item \textbf{obtenais / m'investirais} (hypothèse : \emph{si} + imparfait $\rightarrow$ conditionnel dans la principale). \emph{Obtenir} $\rightarrow$ \emph{j'obtenais} ; \emph{s'investir} $\rightarrow$ \emph{je m'investirais}.
  \item \textbf{expliquiez} (subjonctif présent — \emph{voudrions que} + souhait). Radical \emph{expliqu-} + terminaison \emph{-iez} (2\up{e} personne du pluriel, \emph{vous}).
\end{enumerate}
\end{exoresolu}

\begin{exercice}[Entraînement individuel — À rendre corrigé]
Conjuguez au mode approprié (subjonctif ou conditionnel présent) :
\begin{enumerate}
  \item Il faut que je \_\_\_\_\_\_\_\_\_\_ (partir) tôt pour arriver à l'heure. (Indice : nécessité)
  \item Le directeur aimerait que nous \_\_\_\_\_\_\_\_\_\_ (préparer) un dossier de veille technologique.
  \item Si vous \_\_\_\_\_\_\_\_\_\_ (avoir) le baccalauréat technique, vous \_\_\_\_\_\_\_\_\_\_ (pouvoir) intégrer le BTS.
  \item Je \_\_\_\_\_\_\_\_\_\_ (souhaiter) que mon tuteur \_\_\_\_\_\_\_\_\_\_ (me confier) des responsabilités.
  \item Nous \_\_\_\_\_\_\_\_\_\_ (aimer) que le stage \_\_\_\_\_\_\_\_\_\_ (se dérouler) dans de bonnes conditions.
\end{enumerate}
\end{exercice}

\begin{corrige}[Exercice n°1]
1. \textbf{parte} (subjonctif après \emph{il faut que}) — 2. \textbf{préparions} (subjonctif après \emph{aimerait que}) — 3. \textbf{aviez / pourriez} (hypothèse : \emph{si} + imparfait $\rightarrow$ conditionnel) — 4. \textbf{souhaiterais / me confie} (conditionnel de politesse + subjonctif de souhait) — 5. \textbf{aimerions / se déroule} (conditionnel de politesse + subjonctif de souhait).
\end{corrige}

\courssub{Remédiation rédactionnelle : correction collective d'un CV anonyme}

La correction entre pairs, guidée par l'enseignant, est le levier le plus puissant pour ancrer les normes. Nous utilisons un extrait de CV anonymisé, projeté au tableau, volontairement truffé d'erreurs représentatives.

\begin{exemplebox}[Extrait de CV anonyme — Version « brute » (erreurs à corriger)]
\textbf{TITRE :} Technicien Maintenance Industrielle

\textbf{EXPÉRIENCES}
\begin{itemize}
  \item Juin 2023 -- Août 2023 : stagiaire chez Socametal, Douala.
  \item Taches : surveillance machines, réparation pannes, nettoyage atelier.
  \item J'ai appris a soudé a l'arc et utiliser une fraiseuse.
\end{itemize}

\textbf{COMPÉTENCES}
\begin{itemize}
  \item Logiciels : Word, Excel, PowerPoint.
  \item Langues : français (maternel), anglais (niveau B1).
  \item Permis B.
\end{itemize}

\textbf{CENTRES D'INTÉRÊTS} : football, musique, lecture.
\end{exemplebox}

\begin{methode}[Correction collective en 4 passes (méthode du « zoom »)]
\begin{enumerate}
  \item \textbf{Passe 1 — Structure et complétude :} les 5 rubriques sont-elles présentes ? Le titre est-il un poste réel ? Les dates sont-elles alignées ? $\rightarrow$ Ici : la rubrique « Formation » manque ; « Taches » doit devenir « Missions » ; les puces sont trop vagues.
  \item \textbf{Passe 2 — Verbes d'action et chiffres :} remplacer les noms et infinitifs par des participes passés ou des verbes d'action + résultat chiffré. $\rightarrow$ \emph{« Surveillance machines »} devient \emph{« Assuré la surveillance de 5 lignes de production »} ; \emph{« J'ai appris à souder »} devient \emph{« Réalisé des soudures à l'arc sur structures métalliques »}.
  \item \textbf{Passe 3 — Langue et orthographe :} chasse aux fautes (accords, homophones, majuscules). $\rightarrow$ \emph{« appris a soudé »} doit devenir \emph{« appris à souder »} (préposition \emph{à} + infinitif) ; \emph{« et utiliser »} devient \emph{« et à utiliser »} (reprise de la préposition).
  \item \textbf{Passe 4 — Mise en page et icônes :} alignement par tabulations, police unique, icônes cohérentes, export PDF.
\end{enumerate}
\end{methode}

\begin{exemplebox}[Extrait de CV — Version corrigée collectivement]
\textbf{TITRE :} \textsc{Technicien de maintenance industrielle} — \emph{spécialité électromécanique}

\textbf{FORMATION}

2021--2024 \quad \textsc{Baccalauréat technique, série Maintenance des équipements industriels} — Lycée technique de Douala — Mention Assez Bien

\textbf{EXPÉRIENCES PROFESSIONNELLES}

Juin--Août 2023 \quad \textsc{Stagiaire en maintenance} — Socametal, Douala
\begin{itemize}
  \item \textbf{Assuré} la maintenance préventive de premier niveau sur 3 presses hydrauliques (graissage, contrôles hebdomadaires).
  \item \textbf{Diagnostiqué} et \textbf{réparé} 2 pannes sur variateurs de fréquence, réduisant le temps d'arrêt de production.
  \item \textbf{Réalisé} des soudures à l'arc (électrode enrobée) sur châssis de convoyeurs.
\end{itemize}

\textbf{COMPÉTENCES TECHNIQUES}

Automatisme : automates programmables (initiation), Grafcet \quad Soudure : arc, MIG (initiation) \quad Bureautique : Word, Excel, PowerPoint

\textbf{LANGUES}

Français : langue maternelle \quad Anglais : niveau B1 (lecture de documents techniques)

\textbf{PERMIS / MOBILITÉ}

Permis B \quad Disponible sur Douala et sa région

\textbf{ENGAGEMENTS}

Capitaine de l'équipe de football du lycée (organisation de tournois) \quad Bénévole à la Croix-Rouge camerounaise (logistique des collectes)
\end{exemplebox}

\courssub{Trace écrite finale : la fiche mémo du candidat (à copier dans le cahier)}

Cette trace condense tout le module. Elle sert de \cle{check-list} avant tout dépôt de candidature.

\begin{definition}[Le CV technique — Définition synthétique]
Document normé (1 page, format PDF), qui présente \textbf{par ordre antéchronologique} le parcours de formation, les expériences professionnelles (stages, emplois) et les compétences techniques et transversales d'un candidat, en vue d'obtenir un entretien d'embauche.
\end{definition}

\begin{propriete}[Structure type en 5 blocs — Ordre immuable]
\begin{enumerate}
  \item \textbf{En-tête :} identité + contacts professionnels (adresse électronique sérieuse de type \texttt{prenom.nom@domaine.com}, téléphone).
  \item \textbf{Titre et projet :} poste visé + 1 phrase d'accroche (compétence clé + valeur ajoutée).
  \item \textbf{Formation :} diplômes (Probatoire, Baccalauréat technique, certificats), établissement, année, mention, options techniques.
  \item \textbf{Expériences :} stages et emplois. Pour chacun : dates | entreprise | poste | 3 à 4 puces \textbf{verbe d'action + contexte + résultat chiffré}.
  \item \textbf{Compétences et vie associative :} techniques (machines, logiciels, normes), langues (niveau), permis et mobilité, centres d'intérêt (preuve de qualités humaines).
\end{enumerate}
\end{propriete}

\begin{propriete}[Les 4 commandements de la mise en page]
\begin{enumerate}
  \item \textbf{Police unique} sans empattement (Calibri ou Arial), taille 10 à 11, interligne 1,15.
  \item \textbf{Alignement parfait} par tabulations (dates à gauche, contenu à droite).
  \item \textbf{Hiérarchie visuelle :} rubriques en MAJUSCULES GRAS (13 points), sous-rubriques en minuscules gras (11 points).
  \item \textbf{Zéro ornement :} pas de cadre, pas de couleur de fond ; icônes monochromes (bleu foncé ou gris) uniquement pour le contact, le lieu, les langues et l'informatique.
\end{enumerate}
\end{propriete}

\begin{propriete}[Subjonctif ou conditionnel — Mémento rapide]
\begin{center}
\begin{tabularx}{\linewidth}{|X|l|X|}\hline
\rowcolor{popblueL} \textbf{Déclencheur} & \textbf{Mode} & \textbf{Exemple CV / entretien} \\ \hline
\emph{il faut que, il est nécessaire que, je veux que, je souhaite que} & \textbf{Subjonctif} & \emph{« Il faut que je \textbf{maîtrise} la CAO. »} \\ \hline
\emph{je voudrais, je souhaiterais, j'aimerais} (politesse, souhait atténué) & \textbf{Conditionnel} & \emph{« Je \textbf{souhaiterais} rejoindre votre équipe. »} \\ \hline
\emph{si} + imparfait (hypothèse) & \textbf{Conditionnel} & \emph{« Si j'\textbf{obtenais} le poste, je \textbf{m'investirais} pleinement. »} \\ \hline
\end{tabularx}
\end{center}
\end{propriete}

\begin{attention}[Piège fatal : la faute d'inattention]
Une \textbf{seule faute d'orthographe, de grammaire ou de typographie} (accord manquant, homophone mal choisi, majuscule oubliée) peut faire basculer un CV dans la pile des refusés lors d'une lecture diagonale de quelques secondes. \textbf{La relecture croisée (pair + enseignant) est obligatoire.} Ne vous fiez pas au correcteur automatique : il ne signale ni \emph{« stages effectuées »} (accord fautif) ni \emph{« j'ai appris a soudé »} (confusion \emph{a}/\emph{à} et infinitif mal construit).
\end{attention}

\begin{methode}[S'auto-évaluer avec la grille du jury — Dernière ligne droite avant dépôt]
\begin{enumerate}
  \item \textbf{Imprimez votre CV} (sur papier, le regard change).
  \item \textbf{Lisez-le à voix haute :} entendez-vous les phrases lourdes, les répétitions ?
  \item \textbf{Cochez chaque item de la grille ci-dessous} (objectif : 10/10 items validés).
  \item \textbf{Corrigez} chaque item non coché $\rightarrow$ nouvelle version PDF $\rightarrow$ nouvelle vérification.
  \item \textbf{Nommez le fichier :} \texttt{NOM\_Prenom\_IntitulePoste.pdf}.
\end{enumerate}
\end{methode}

\begin{center}
\begin{tabularx}{\linewidth}{|c|X|}\hline
\rowcolor{popblueL} \textbf{OK} & \textbf{Critère de validation (grille du jury)} \\ \hline
$\square$ & Titre précis (poste + spécialité) + phrase d'accroche percutante \\ \hline
$\square$ & 5 rubriques présentes, ordre respecté, antéchronologique \\ \hline
$\square$ & Expériences : verbes d'action variés + résultats chiffrés (au moins 1 par stage) \\ \hline
$\square$ & Compétences techniques : logiciels, machines, normes cités nommément \\ \hline
$\square$ & Langues : niveau précisé (A2, B1, B2...) \\ \hline
$\square$ & 0 faute d'orthographe, de grammaire ou de syntaxe (relecture à voix haute faite) \\ \hline
$\square$ & Subjonctif et conditionnel justes dans la phrase d'accroche et la lettre de motivation \\ \hline
$\square$ & Mise en page : police unique, alignement par tabulations, 1 page maximum \\ \hline
$\square$ & Icônes cohérentes, discrètes, utiles (contact, lieu, langues, informatique) \\ \hline
$\square$ & Fichier PDF testé (ouvert sur téléphone et sur ordinateur), nom de fichier normé \\ \hline
\end{tabularx}
\end{center}

\courssub{Auto-évaluation : suis-je prêt pour le Baccalauréat technique et l'insertion ?}

Chaque élève remplit cette grille en toute honnêteté. Elle n'est pas notée : elle guide le travail personnel restant.

\begin{aretenir}[Grille d'auto-positionnement — Objectifs de la leçon]
\begin{center}
\begin{tabularx}{\linewidth}{|l|X|c|c|c|}\hline
\rowcolor{popblueL} \textbf{Compétence} & \textbf{Indicateur observable} & \textbf{Non acquis} & \textbf{En cours} & \textbf{Maîtrisé} \\ \hline
\textbf{C1. Produire un CV} & Je sais structurer mes 5 rubriques dans l'ordre normé. & $\square$ & $\square$ & $\square$ \\ \hline
\textbf{C2. Rédiger des puces} & J'utilise « verbe d'action + contexte + chiffre » pour chaque expérience. & $\square$ & $\square$ & $\square$ \\ \hline
\textbf{C3. Maîtriser la forme} & Mon CV tient en 1 page, police unique, aligné, 0 faute, PDF soigné. & $\square$ & $\square$ & $\square$ \\ \hline
\textbf{C4. Langue : les modes} & Je choisis le subjonctif (nécessité, souhait) ou le conditionnel (hypothèse, politesse) sans erreur. & $\square$ & $\square$ & $\square$ \\ \hline
\textbf{C5. Présenter à l'oral} & Je tiens un pitch de 2 minutes structuré (qui / quoi / pourquoi), avec regard, voix posée et preuves STAR. & $\square$ & $\square$ & $\square$ \\ \hline
\textbf{C6. Répondre au jury} & J'anticipe 3 questions pièges et j'ai préparé des réponses argumentées. & $\square$ & $\square$ & $\square$ \\ \hline
\end{tabularx}
\end{center}
\end{aretenir}

\begin{savaistu}[CV numérique : votre passeport pour l'emploi]
Conserver son CV à jour en format \textbf{PDF} (avec sa version modifiable en sauvegarde) permet de postuler \textbf{en quelques minutes} à toute offre reçue par courriel, par messagerie ou repérée sur les plateformes d'emploi et les sites d'entreprises. Programmez un rappel \textbf{tous les 3 mois} : « Mise à jour du CV » (nouveau stage, nouvelle compétence, nouveau niveau de langue, nouveau projet). Un CV vivant est un CV qui sert.
\end{savaistu}

\begin{exemplebox}[Bilan personnel — À compléter en bas de la grille]
\textbf{Mes 3 forces validées aujourd'hui :}

1. \dotfill

2. \dotfill

3. \dotfill

\textbf{Mon objectif prioritaire pour la semaine prochaine :}

\dotfill

\textbf{Action concrète (date, ressource, personne ressource) :}

\dotfill
\end{exemplebox}

\emph{Fin de la leçon N24. Le CV est prêt, l'élève est outillé : la porte de l'entretien d'embauche — et celle du Baccalauréat technique — est grande ouverte.}

\newpage

\coursec{Activité d'intégration}

\courssub{Objectifs de l'activité}
Cette activité d'intégration vise à mobiliser l'ensemble des savoirs et savoir-faire travaillés au cours de la leçon dans une situation professionnelle authentique. Elle permet à l'élève de :
\begin{itemize}
    \item Analyser une offre d'emploi réaliste pour en extraire les compétences-clés et les mots-clés attendus par le recruteur.
    \item Produire un CV complet, structuré, mis en page selon les normes professionnelles et adapté à l'offre ciblée (texte fonctionnel).
    \item Mobiliser les connaissances en morphosyntaxe (valeurs du conditionnel de politesse et du subjonctif) dans une production orale formelle.
    \item Présenter son parcours et sa motivation devant un jury simulé en adoptant une posture professionnelle (communication non verbale, élocution, gestion du temps).
    \item Évaluer le travail d'un pair à l'aide d'une grille critériée (fond, forme, comportement) et formuler un compte rendu argumenté.
    \item Identifier collectivement les erreurs récurrentes pour engager une remédiation ciblée.
\end{itemize}

\courssub{Situation}
\textbf{Contexte :} La \textbf{SABC} (Société Anonyme des Brasseries du Cameroun), leader du secteur des boissons au Cameroun, lance une campagne de recrutement pour ses sites de Douala et de Yaoundé. L'annonce officielle, parue dans \emph{Cameroon Tribune} du 15 octobre 2026 et relayée sur le site de l'entreprise, mentionne le profil suivant :

\begin{exemplebox}[Extrait de l'offre d'emploi SABC -- Réf. SABC/TECH/2026-042]
\textbf{Poste :} Technicien de Maintenance Industrielle (H/F) -- 3 postes \\
\textbf{Localisation :} Douala (2 postes) / Yaoundé (1 poste) \\
\textbf{Type de contrat :} CDI (période d'essai de 3 mois renouvelable) \\
\textbf{Niveau d'études requis :} Baccalauréat Technique (séries F3, F4, G) ou BTS Maintenance / Électrotechnique / Mécanique. \\
\textbf{Expérience :} 0 à 2 ans (stages académiques inclus). Jeunes diplômés encouragés. \\
\textbf{Compétences techniques attendues :}
\begin{itemize}
    \item Maintenance préventive et curative sur lignes d'embouteillage ou lignes de brassage.
    \item Lecture de schémas électriques, pneumatiques et hydrauliques.
    \item Diagnostic de pannes sur automates programmables industriels (API) -- \emph{notions de base exigées}.
    \item Connaissance des règles HSE (Hygiène, Sécurité, Environnement) et des procédures de consignation.
    \item Utilisation d'un logiciel de GMAO (Gestion de Maintenance Assistée par Ordinateur) -- un atout.
\end{itemize}
\textbf{Compétences transverses :} rigueur, réactivité, esprit d'équipe, sens du service, communication technique (français impératif, anglais technique lu et écrit apprécié). \\
\textbf{Dossier de candidature :} CV (1 page maximum) + lettre de motivation + copies des diplômes + attestations de stage. \\
\textbf{Date limite :} 30 novembre 2026. \\
\textbf{Processus :} présélection sur dossier $\rightarrow$ tests techniques écrits $\rightarrow$ entretien oral devant un jury technique et RH.
\end{exemplebox}

\courssub{Consignes}
Chaque élève traite la situation individuellement puis participe à l'évaluation collective. Les étapes sont chronométrées pour simuler la pression réelle.

\begin{enumerate}
    \item \textbf{Analyse de l'offre (15 min) :} relevez dans un tableau à deux colonnes les \textbf{compétences techniques obligatoires} et les \textbf{compétences transverses / atouts}. Identifiez 5 mots-clés à faire apparaître dans votre CV pour passer les filtres de tri automatique des candidatures (logiciels ATS).
    \item \textbf{Rédaction du CV (45 min) :} produisez un CV \textbf{strictement tenu sur une page A4}, structuré selon la norme étudiée (en-tête, accroche/profil, formation, expériences/stages, compétences techniques, compétences transverses, centres d'intérêt, références). Adaptez le vocabulaire à l'offre SABC. Soignez la mise en page (alignements, police unique, espaces réguliers). Imprimez 4 exemplaires.
    \item \textbf{Préparation de l'oral (20 min) :} rédigez le script d'une présentation de \textbf{2 minutes maximum} (chronométrée) structurée ainsi : \emph{Qui suis-je ? Que sais-je faire ? Pourquoi ce poste ? Pourquoi moi ?} Intégrez \textbf{au moins deux conditionnels de politesse} (ex. \emph{Je souhaiterais..., Je voudrais mettre en avant...}) et \textbf{au moins un subjonctif} (ex. \emph{Il est essentiel que je maîtrise..., Avant que je n'intègre l'équipe...}).
    \item \textbf{Simulation d'entretien (10 min par candidat) :} la classe est divisée en jurys de 4 personnes (3 élèves + 1 enseignant observateur). Chaque candidat présente son CV (2 min) puis répond aux questions du jury (3 min) : une question technique (ex. \emph{Décrivez la procédure de consignation avant une intervention}), une question comportementale (ex. \emph{Racontez une panne critique gérée en stage}), une question de motivation.
    \item \textbf{Évaluation et compte rendu (10 min par jury) :} le jury note sur 20 selon la grille (Fond /10, Forme /5, Comportement /5). Il rédige un \textbf{compte rendu argumenté} (forces, axes de progression, décision : \emph{Reçu / Reçu avec réserves / Non reçu}).
    \item \textbf{Remédiation collective (20 min) :} mise en commun des erreurs fréquentes relevées par les jurys (fond : inadéquation offre/CV ; forme : fautes, mise en page ; langue : absence de conditionnel/subjonctif, tics de langage ; comportement : regard fuyant, voix basse). L'enseignant synthétise et propose des exercices de reprise.
\end{enumerate}

\courssub{Ressources à mobiliser}
\begin{aretenir}[Structure type du CV technique (1 page)]
\begin{tabularx}{\linewidth}{|l|X|}
\hline
\textbf{Rubrique} & \textbf{Contenu attendu pour un élève de Terminale technique} \\
\hline
En-tête & Nom, prénom, adresse (ville, quartier), téléphone (format +237 6XX XX XX XX), adresse électronique professionnelle (prenom.nom@email.com). \textbf{Photo non obligatoire} ; si photo : format identité, fond neutre, tenue professionnelle. \\
\hline
Accroche / Profil & 2 à 3 lignes : titre visé + spécialité + atout majeur + objectif. Ex. : \emph{« Technicien de maintenance (Bac F4), spécialisé en électrotechnique, stage de 6 mois sur ligne de conditionnement. Recherche un poste de maintenance préventive et curative. Rigoureux et réactif. »} \\
\hline
Formation & Ordre antéchronologique (le plus récent en premier). Diplôme exact, spécialité, établissement, ville, année, \textbf{mention} (si Assez Bien ou supérieur). Projet de fin d'études : titre + une ligne de description technique. \\
\hline
Expériences / Stages & Ordre antéchronologique. Dates (mois/année), entreprise, ville, intitulé du poste. \textbf{3 à 4 puces d'actions concrètes} commençant par un verbe d'action (réalisé, diagnostiqué, programmé, remplacé, participé). Quantifier les résultats (ex. \emph{« maintenance préventive sur 12 moteurs asynchrones »}). \\
\hline
Compétences techniques & Regroupées par domaines : \textbf{électricité/automatisme}, \textbf{mécanique/pneumatique}, \textbf{outils/logiciels}, \textbf{HSE}. Préciser le niveau honnêtement : \emph{maîtrise / bonnes notions / notions}. \\
\hline
Compétences transverses & Langues (français : langue d'usage ; anglais : technique lu/écrit, niveau précisé), informatique (traitement de texte, tableur), qualités (rigueur, réactivité, esprit d'équipe) \textbf{illustrées par un exemple court}. \\
\hline
Centres d'intérêt & 2 à 3 items montrant l'équilibre et des qualités : sport collectif, bricolage/électronique, bénévolat associatif. Éviter les mentions trop vagues (« lecture, voyage, musique »). \\
\hline
Références & \emph{« Disponibles sur demande »} ou nom/fonction/téléphone d'un tuteur de stage ou d'un enseignant référent (accord préalable obtenu). \\
\hline
\end{tabularx}
\end{aretenir}

\begin{aretenir}[Morphosyntaxe : conditionnel (politesse) et subjonctif (nécessité, souhait)]
\begin{tabularx}{\linewidth}{|l|X|X|}
\hline
\textbf{Mode / Valeur} & \textbf{Formulation type à l'entretien} & \textbf{Exemple contextualisé} \\
\hline
\textbf{Conditionnel présent} & \emph{Je voudrais / Je souhaiterais / J'aimerais / Je pourrais} + infinitif & \emph{« Je \textbf{souhaiterais} vous présenter mon expérience sur les automates programmables. »} \\
\textit{Politesse, atténuation, souhait} & \emph{Permettez-moi de / Il me semblerait utile de} & \emph{« \textbf{Permettez-moi de} préciser ma pratique de la GMAO. »} \\
\hline
\textbf{Subjonctif présent} & \textbf{Il faut que / Il est essentiel que / Il est important que / Avant que / Pour que / Afin que} + subjonctif & \emph{« Il est \textbf{essentiel que je maîtrise} la consignation avant d'intervenir. »} \\
\textit{Nécessité, obligation, finalité, antériorité} & \emph{Je veux que / Je préfère que / Avant que je ne} & \emph{« \textbf{Avant que je n'intègre} l'équipe, je révise la documentation technique des lignes. »} \\
\hline
\end{tabularx}
\end{aretenir}

\begin{attention}[Pièges grammaticaux fréquents]
\begin{itemize}
    \item Après \emph{il faut que}, \emph{il est important que}, \emph{avant que}, le verbe est au \textbf{subjonctif}, jamais à l'indicatif : on écrit \emph{« il faut que je \textbf{fasse} »} et non \emph{« il faut que je fais »}.
    \item Le conditionnel de politesse ne s'emploie pas pour les faits réels : on dit \emph{« je souhaiterais intégrer votre équipe »} (souhait) mais \emph{« j'ai réalisé ce stage »} (indicatif, fait accompli).
    \item Après \emph{avant que}, le \emph{ne} explétif est facultatif mais élégant au registre soutenu : \emph{« avant que je (ne) vienne »}.
\end{itemize}
\end{attention}

\begin{aretenir}[Grille d'évaluation (sur 20)]
\begin{tabularx}{\linewidth}{|l|X|c|}
\hline
\textbf{Critère} & \textbf{Indicateurs observables} & \textbf{Points} \\
\hline
\textbf{FOND (10)} & Adéquation profil/offre (mots-clés, compétences techniques citées) ; cohérence parcours/projet ; pertinence des exemples (méthode STAR : Situation, Tâche, Action, Résultat) ; connaissance de l'entreprise. & /10 \\
\hline
\textbf{FORME (5)} & CV : mise en page (1 page, lisibilité, alignements, police unique), orthographe et grammaire (0 faute), vocabulaire technique précis. Oral : élocution (débit, articulation, volume), structure du discours (accroche, arguments, conclusion), \textbf{utilisation effective du conditionnel et du subjonctif}. & /5 \\
\hline
\textbf{COMPORTEMENT (5)} & Ponctualité, tenue professionnelle, contact visuel, posture, écoute active (reformulation de la question), gestion du stress, respect du temps (2 min $\pm$ 15 s), politesse (formules d'entrée et de sortie). & /5 \\
\hline
\textbf{TOTAL} & & \textbf{/20} \\
\hline
\end{tabularx}
\end{aretenir}

\courssub{Démarche et corrigé commenté}

\textbf{Étape 1 -- Analyse de l'offre (modélisation attendue)}
\begin{center}
\begin{tabularx}{\linewidth}{|X|X|}
\hline
\textbf{Compétences techniques obligatoires (mots-clés)} & \textbf{Compétences transverses / atouts} \\
\hline
Maintenance préventive / curative & Rigueur, réactivité \\
Lignes d'embouteillage et de brassage & Esprit d'équipe, sens du service \\
Lecture de schémas électriques & Communication technique (français) \\
Pneumatique / hydraulique & Anglais technique (lu/écrit) \\
Automates programmables industriels (API) & GMAO (atout) \\
Consignation, règles HSE & \\
\hline
\end{tabularx}
\end{center}
\emph{Commentaire :} l'élève doit extraire \textbf{exactement} ces termes pour les réinjecter dans son CV (rubrique Compétences, description des stages). Les 5 mots-clés prioritaires : \cle{maintenance préventive}, \cle{automates programmables}, \cle{consignation}, \cle{schémas électriques}, \cle{GMAO}.

\textbf{Étape 2 -- Production du CV (exemple commenté pour un élève de Bac F4 Électrotechnique, 19 ans, Douala)}

\begin{exemplebox}[CV -- NGUEMA David (modèle adapté à l'offre)]
\textbf{NGUEMA David} \hfill +237 6 99 88 77 66 \hfill david.nguema@email.com \\
Douala, Bonabéri \hfill Permis B \\
\rule{\linewidth}{0.5pt}
\textbf{TECHNICIEN DE MAINTENANCE INDUSTRIELLE -- SPÉCIALITÉ ÉLECTROTECHNIQUE} \\
\emph{Titulaire du Bac F4 (mention Assez Bien) -- stage de 6 mois en maintenance sur ligne de conditionnement. Recherche un poste de maintenance préventive et curative.}
\rule{\linewidth}{0.5pt}
\textbf{FORMATION} \\
\textbf{2023--2026} \hfill \textbf{Baccalauréat Technique F4 -- Électrotechnique} -- mention Assez Bien \\
Lycée technique de Douala-Bonabéri \\
\textit{Projet de fin d'études :} « Automatisation d'un système de convoyage par automate programmable » (programmation, câblage d'armoire, mise en service). \\
\textbf{2020--2023} \hfill \textbf{Probatoire technique F4} -- Lycée technique de Douala.
\rule{\linewidth}{0.5pt}
\textbf{EXPÉRIENCES PROFESSIONNELLES ET STAGES} \\
\textbf{Fév. -- juil. 2026 (6 mois)} \hfill \textbf{Stagiaire technicien de maintenance} -- \textbf{SOCAVER (Douala)} -- verrerie industrielle \\
\begin{itemize}
    \item \textbf{Maintenance préventive hebdomadaire} sur 4 fours de refusion (brûleurs, régulation de température, chaînes de convoyage) -- \emph{aucun arrêt imprévu sur la période}.
    \item \textbf{Diagnostic et réparation} de variateurs de fréquence et de moteurs asynchrones triphasés.
    \item \textbf{Lecture de schémas électriques} et pneumatiques pour la localisation des pannes (capteurs, électrovannes).
    \item \textbf{Initiation à la GMAO} : saisie des bons d'intervention, consultation des historiques de matériel.
    \item \textbf{Application stricte de la consignation} avant toute intervention (procédure HSE interne).
\end{itemize}
\textbf{Juin -- juil. 2025 (2 mois)} \hfill \textbf{Stage d'ouvrier électricien} -- \textbf{ENEO (Douala)} \\
Raccordement de branchements basse tension, pose de compteurs, respect des normes de sécurité.
\rule{\linewidth}{0.5pt}
\textbf{COMPÉTENCES TECHNIQUES} \\
\textbf{Automatisme / électricité :} automates programmables (bonnes notions), variateurs de vitesse, lecture de schémas électriques, câblage d'armoires. \\
\textbf{Mécanique / fluides :} pneumatique et hydraulique (notions), maintenance de pompes et réducteurs. \\
\textbf{Outils / logiciels / HSE :} GMAO (notions), multimètre, pince ampèremétrique, consignation, port des EPI, analyse des risques.
\rule{\linewidth}{0.5pt}
\textbf{COMPÉTENCES TRANSVERSES ET LANGUES} \\
\begin{itemize}
    \item \textbf{Français :} langue d'usage -- expression technique écrite et orale structurée.
    \item \textbf{Anglais :} technique lu/écrit (lecture de manuels et de fiches techniques).
    \item \textbf{Rigueur / réactivité :} respect des délais d'intervention (panne de four résolue en 45 minutes).
    \item \textbf{Esprit d'équipe :} travail en binôme avec un technicien senior, transmission des consignes à l'équipe suivante.
\end{itemize}
\rule{\linewidth}{0.5pt}
\textbf{CENTRES D'INTÉRÊT} \\
Football (club de quartier -- \emph{esprit collectif, endurance}) -- \textbf{électronique} (montages personnels : automatisation d'un portail) -- bénévolat (premiers secours).
\rule{\linewidth}{0.5pt}
\textbf{RÉFÉRENCES} \\
M. \textbf{KAMGA Joseph}, chef de service maintenance, SOCAVER Douala -- tél. 6 77 66 55 44 (accord donné) \\
M. \textbf{FOUDA Pierre}, professeur d'électrotechnique, Lycée technique de Douala -- tél. 6 55 44 33 22.
\end{exemplebox}

\emph{Commentaire détaillé du correcteur :}
\begin{itemize}
    \item \textbf{Fond (9/10) :} l'accroche cite le poste exact visé. Le projet de fin d'études (automate programmable) et le stage SOCAVER (fours, variateurs, GMAO, consignation) répondent point par point à l'offre. Les mots-clés sont présents. Les résultats sont quantifiés (« aucun arrêt imprévu », « 45 minutes »). L'anglais technique est mentionné.
    \item \textbf{Forme (5/5) :} une page exacte, police unique, alignements réguliers, zéro faute, vocabulaire précis.
    \item \textbf{Adaptation :} la rubrique « Compétences techniques » regroupe les savoir-faire par domaine (et non en liste brute). Les niveaux sont honnêtes (\emph{bonnes notions}, \emph{notions}) : exagérer son niveau est un piège classique démasqué aux tests techniques.
\end{itemize}

\textbf{Étape 3 -- Script oral de 2 minutes (marqueurs grammaticaux mis en évidence)}

\begin{exemplebox}[Script de présentation orale -- David NGUEMA]
\textit{(Entre, serre la main, contact visuel, sourire)} \\
\emph{« Bonjour Mesdames, Messieurs. \textbf{Je me permets de me présenter} : David Nguema, 19 ans, titulaire d'un baccalauréat technique F4 Électrotechnique, mention Assez Bien, au Lycée technique de Douala.}

\emph{\textbf{Je souhaiterais} d'abord mettre en avant mon stage de six mois à la SOCAVER, verrerie industrielle. J'y ai assuré la maintenance préventive sur des fours de refusion et des convoyeurs, dans le respect de la consignation -- \textbf{il est essentiel que je respecte} ces procédures pour la sécurité de tous.}

\emph{J'ai également programmé un \textbf{automate programmable} pour mon projet de fin d'études : l'automatisation d'un convoyage. \textbf{Il est important que je poursuive} cette montée en compétence sur les équipements utilisés sur vos lignes.}

\emph{\textbf{Je voudrais} insister sur ma rigueur : lors d'une panne critique de variateur à la SOCAVER, j'ai diagnostiqué et remplacé le composant en 45 minutes, évitant deux heures d'arrêt de ligne. \textbf{Avant que je n'intègre} votre équipe de maintenance, je m'engage à réviser en autonomie la documentation technique de vos installations.}

\emph{\textbf{J'aimerais} conclure en disant que rejoindre la SABC, fleuron industriel camerounais, représente pour moi l'opportunité d'appliquer ma formation sur des installations de pointe, dans le respect de vos standards HSE. Je vous remercie de votre attention et \textbf{je reste à votre disposition} pour toute question technique. »}
\end{exemplebox}

\emph{Analyse grammaticale :}
\begin{itemize}
    \item \textbf{Conditionnels de politesse :} \emph{je souhaiterais}, \emph{je voudrais}, \emph{j'aimerais} ; s'ajoutent les formules atténuées \emph{je me permets de} et \emph{je reste à votre disposition}. Objectif ($\geq 2$) atteint.
    \item \textbf{Subjonctifs :} \emph{il est essentiel que je respecte}, \emph{il est important que je poursuive}, \emph{avant que je n'intègre} (subjonctif après « avant que », avec \emph{ne} explétif). Objectif ($\geq 1$) atteint.
    \item \textbf{Structure :} accroche (identité + diplôme) $\rightarrow$ preuve du stage (compétences techniques + HSE) $\rightarrow$ preuve du projet (automatisme) $\rightarrow$ preuve comportementale (méthode STAR) $\rightarrow$ motivation pour l'entreprise $\rightarrow$ ouverture aux questions. Durée mesurée : \textbf{1 min 55 s} (dans la fenêtre 1 min 45 s -- 2 min 15 s).
\end{itemize}

\textbf{Étape 4 -- Simulation d'entretien (exemples de questions du jury et réponses attendues)}

\begin{tabularx}{\linewidth}{|l|X|}
\hline
\textbf{Question du jury} & \textbf{Réponse attendue (éléments clés)} \\
\hline
\textit{Technique :} « Décrivez la procédure de consignation que vous appliquiez en stage. » & 1. Identifier les sources d'énergie (électrique, pneumatique, hydraulique, gravité). 2. Couper et isoler (disjoncteur, vanne). 3. \textbf{Verrouiller} (cadenas personnel + étiquette nom/date). 4. \textbf{Dissiper les énergies résiduelles} (condensateurs, pression). 5. \textbf{Vérifier l'absence de tension ou de pression}. 6. Intervenir. 7. \textbf{Déconsigner} dans l'ordre inverse (retrait du cadenas par son poseur uniquement). \\
\hline
\textit{Comportementale (STAR) :} « Racontez une panne critique gérée en stage. » & \textbf{S :} arrêt brutal du four n\textdegree{}3 à 14 h (pic de production). \textbf{T :} remettre en route en moins d'une heure. \textbf{A :} diagnostic du variateur (défaut de surintensité), vérification du moteur, remplacement du variateur par la pièce de rechange, reparamétrage, essai à vide puis en charge. \textbf{R :} redémarrage à 14 h 42, rapport rédigé dans la GMAO, cause identifiée (encrassement du ventilateur) $\rightarrow$ action préventive de nettoyage planifiée. \\
\hline
\textit{Motivation :} « Pourquoi la SABC et pas une autre entreprise ? » & Leader national, installations modernes et réputées, politique de formation continue, ancrage local, fierté des marques. Éviter les réponses génériques (« pour le salaire », « c'est près de chez moi »). \\
\hline
\end{tabularx}

\textbf{Étape 5 -- Grille d'évaluation remplie (exemple du candidat David)}

\begin{center}
\begin{tabularx}{\linewidth}{|l|X|c|}
\hline
\textbf{Critère} & \textbf{Commentaire du jury} & \textbf{Note} \\
\hline
\textbf{FOND (10)} & Profil parfaitement calibré sur l'offre. Mots-clés intégrés. Exemples concrets et chiffrés. Connaissance de l'entreprise. Projet de fin d'études pertinent. & \textbf{9,5} \\
\hline
\textbf{FORME (5)} & CV impeccable (1 page, 0 faute, mise en page professionnelle). Oral fluide, structuré, vocabulaire précis. \textbf{3 conditionnels et 3 subjonctifs} utilisés naturellement. Diction claire. & \textbf{5} \\
\hline
\textbf{COMPORTEMENT (5)} & Tenue correcte (chemise, pantalon). Contact visuel constant. Posture ouverte. Écoute active (reformule la question technique). Temps respecté (1 min 55 s). Politesse à l'entrée et à la sortie. & \textbf{5} \\
\hline
\textbf{TOTAL} & \textbf{Reçu -- profil solide, prêt pour les tests techniques écrits.} & \textbf{19,5 / 20} \\
\hline
\end{tabularx}
\end{center}

\textbf{Étape 6 -- Compte rendu argumenté du jury (modèle rédigé par les élèves jurés)}

\begin{exemplebox}[Compte rendu du jury -- Candidat NGUEMA David]
\textbf{Décision : REÇU (19,5/20)} \\
\textbf{Forces majeures :} adéquation parfaite entre le profil et le poste. CV « miroir » de l'offre. Stage SOCAVER exploité intelligemment (vocabulaire professionnel, quantification). Maîtrise grammaticale orale (conditionnel, subjonctif) naturelle, non forcée. Posture professionnelle mature. \\
\textbf{Axes de progression (mineurs) :} préciser le niveau d'anglais par une certification ou un niveau de référence ; lier explicitement le centre d'intérêt « électronique » à l'automatisme industriel (capteurs, actionneurs). \\
\textbf{Recommandation :} convoquer en priorité aux tests techniques écrits.
\end{exemplebox}

\textbf{Étape 7 -- Remédiation collective (synthèse de l'enseignant après l'activité)}

\begin{attention}[Erreurs fréquentes relevées par les jurys -- pistes de remédiation]
\begin{enumerate}
    \item \textbf{Fond -- CV « passe-partout » :} environ 40 \% des CV ne contenaient pas les mots-clés \emph{consignation}, \emph{automates}, \emph{GMAO}, \emph{schémas électriques}. \textbf{Remédiation :} exercice « mots-clés » : à partir de 3 offres différentes, lister 10 mots-clés techniques et les placer dans la rubrique Compétences.
    \item \textbf{Fond -- expériences vagues :} « stage maintenance » sans verbe d'action ni quantification. \textbf{Remédiation :} atelier de réécriture des puces selon la méthode STAR : \emph{verbe d'action + contexte technique + résultat chiffré}. Ex. : \emph{« Remplacé 3 roulements sur une pompe centrifuge, alignement vérifié, vibration résiduelle nulle. »}
    \item \textbf{Forme -- dépassement d'une page / mise en page instable :} marges variables, polices multiples. \textbf{Remédiation :} modèle verrouillé (styles de titre, tabulations fixes), impression en PDF obligatoire.
    \item \textbf{Langue -- absence du conditionnel et du subjonctif à l'oral :} environ 60 \% des élèves ont utilisé l'indicatif (\emph{« je veux dire que... », « il faut que je fais... »}). \textbf{Remédiation :} exercices de transformation : \emph{« je veux que tu viennes »} $\rightarrow$ \emph{« je souhaiterais que tu viennes »} / \emph{« il faut que tu viennes »}. Entraînement oral en binôme (chronomètre de 2 min).
    \item \textbf{Comportement -- regard fuyant, voix basse, temps non géré :} environ 30 \% des candidats ont parlé plus de 3 min ou moins d'1 min. \textbf{Remédiation :} simulation « ascenseur » (présentation éclair de 30 s) puis présentation de 2 min ; enregistrement vidéo et auto-évaluation avec la grille de comportement.
\end{enumerate}
\end{attention}

\courssub{Bilan de l'activité}
Cette activité d'intégration a permis de valider la \textbf{transférabilité des acquis} de la leçon vers une situation professionnelle réaliste et exigeante. Les points saillants mis en évidence sont :

\begin{enumerate}
    \item \textbf{L'analyse de l'offre comme boussole :} les élèves ayant réussi sont ceux qui ont traité le CV comme une \textbf{réponse ciblée} à une commande (l'offre), et non comme une autobiographie. La compétence « extraire les mots-clés » est désormais comprise comme une technique professionnelle essentielle.
    \item \textbf{L'interdépendance de l'écrit et de l'oral :} le CV structure le discours oral (l'accroche orale reprend l'accroche écrite, les puces des stages deviennent des preuves STAR). La maîtrise de la mise en page libère l'esprit du candidat pour le fond lors de l'entretien.
    \item \textbf{La grammaire au service de la relation professionnelle :} le conditionnel de politesse et le subjonctif de nécessité ne sont plus des exercices scolaires mais des \textbf{marqueurs de professionnalisme} (distance respectueuse, engagement sur la sécurité). Leur absence a été perçue par les jurys d'élèves comme un manque de maturité.
    \item \textbf{L'évaluation par les pairs comme levier d'apprentissage :} le fait d'être juré a obligé les élèves à expliciter les critères de qualité (grille), développant leur \textbf{métacognition}. Le compte rendu argumenté a structuré leur pensée critique.
    \item \textbf{La remédiation ciblée :} l'analyse collective des erreurs a permis de prioriser les reprises : \textbf{verbes d'action + quantification} pour le CV, \textbf{transformation indicatif $\rightarrow$ conditionnel/subjonctif} pour l'oral, \textbf{gestion du temps + regard} pour le comportement.
\end{enumerate}

L'élève de Terminale technique sort de cette séquence avec un \textbf{CV prêt à envoyer} (fichier PDF nommé par exemple \texttt{Nguema\_David\_CV\_TechnicienMaintenance.pdf}), un \textbf{pitch oral rodé} et une \textbf{grille d'auto-évaluation} réutilisable pour tout futur entretien. C'est l'aboutissement concret de la formation en communication professionnelle du cycle terminal technique.

\newpage

\coursec{Méthodes et exercices résolus}

\courssub{Méthode 1 : Analyser un CV comme texte fonctionnel}

\begin{methode}[Analyser un texte fonctionnel : le CV]
\textbf{Objectif :} Identifier la nature, la finalité et la structure d'un CV pour en comprendre la logique communicationnelle.
\begin{enumerate}
    \item \textbf{Identifier la situation de communication :} Qui écrit ? (Le candidat) À qui ? (Le recruteur) Dans quel but ? (Obtenir un entretien). \cle{Finalité persuasive}.
    \item \textbf{Repérer les conventions du genre :} Document normé, concis (1 page idéalement), structuré en rubriques standardisées, langage télégraphique (noms, verbes d'action, pas de phrases complètes), mise en page aérée.
    \item \textbf{Analyser l'organisation de l'information :} Ordre antéchronologique (du plus récent au plus ancien) ou thématique (par compétences). Hiérarchisation : coordonnées \rightarrow accroche/profil \rightarrow expériences \rightarrow formation \rightarrow compétences \rightarrow centres d'intérêt.
    \item \textbf{Évaluer l'adéquation poste/profil :} Vérifier la présence de \cle{mots-clés} de l'offre d'emploi (compétences techniques, logiciels, langues, soft skills).
    \item \textbf{Conclure sur l'efficacité :} Lisibilité (police, interlignes), cohérence du parcours, absence de fautes, impact visuel.
\end{enumerate}
\end{methode}

\begin{exoresolu}[Analyse d'un extrait de CV]
\textbf{Énoncé :} Vous trouvez ci-dessous un extrait du CV de M. Kamga, candidat au poste de \emph{Technicien de maintenance industrielle} dans une agro-industrie à Douala. Analysez la pertinence de cet extrait par rapport au poste visé.

\begin{center}
\begin{tabularx}{\linewidth}{|l|X|}
\hline
\textbf{Rubrique} & \textbf{Contenu} \\ \hline
\textbf{Expériences} & 2022--2024 : \textbf{Opérateur de production} -- Usine SOSUCAM (Nkoteng).\\
& \textbullet\ Surveillance ligne d'embouteillage (cadence 36 000 b/h).\\
& \textbullet\ Maintenance de 1er niveau : changement courroies, graissage, capteurs.\\
& \textbullet\ Respect normes HACCP / Sécurité. \\ \hline
\textbf{Formation} & 2021 : \textbf{Baccalauréat F3 (Électrotechnique)} -- Lycée Technique de Bafoussam.\\
& 2024 : \textbf{CQP Maintenance Industrielle} -- CFA Douala (en cours). \\ \hline
\textbf{Compétences} & \textbullet\ Lecture plans électriques / schémas pneumatiques.\\
& \textbullet\ GMAO (logiciel MAINTI4).\\
& \textbullet\ Soudure à l'arc / TIG (notions).\\
& \textbullet\ Anglais technique (niveau B1). \\ \hline
\end{tabularx}
\end{center}

\tcblower
\textbf{Solution détaillée :}

\textbf{1. Situation de communication :} L'émetteur est M. Kamga (candidat), le destinataire est le DRH / Responsable maintenance de l'agro-industrie. La finalité est professionnelle : convaincre de sa capacité à assurer la maintenance curative et préventive.

\textbf{2. Conventions du genre respectées :}
\begin{itemize}
    \item \textbf{Concision :} Phrases nominales, verbes d'action (\emph{surveillance, maintenance, respect}), pas de "je".
    \item \textbf{Structure claire :} Rubriques distinctes, puces pour la lecture diagonale.
    \item \textbf{Antéchronologie :} Expérience la plus récente (2022-2024) en premier.
\end{itemize}

\textbf{3. Adéquation Poste/Profil (Analyse critique) :}
\begin{itemize}
    \item \textbf{Points forts (Matching fort) :}
    \begin{itemize}
        \item \textbf{Secteur d'activité :} Expérience en agro-industrie (SOSUCAM) = connaissance des normes \cle{HACCP}, cadences élevées, contraintes hygiène/sécurité. \textbf{Atout majeur}.
        \item \textbf{Technique :} Maintenance 1er niveau effectuée (courroies, capteurs, graissage) correspond au cœur du poste.
        \item \textbf{Outils :} Maîtrise \cle{GMAO (MAINTI4)} = autonomie immédiate sur la gestion des interventions.
        \item \textbf{Formation :} Bac F3 (base électrotechnique solide) + CQP Maintenance en cours = double compétence théorie/pratique validée.
    \end{itemize}
    \item \textbf{Points de vigilance / Améliorations :}
    \begin{itemize}
        \item \emph{Soudure TIG} : Notions seulement. Préciser "initiée" ou "en formation" pour ne pas sur-vendre.
        \item \emph{Anglais B1} : Suffisant pour lire des notices techniques, mais préciser "lu, écrit, technique" à l'oral.
        \item \emph{Accroche/Profil manquante :} Une phrase de 2 lignes en haut du CV résumant : "Technicien maintenance, 2 ans agro-industrie, CQP en cours, GMAO, HACCP, disponible immédiatement" aurait renforcé l'impact immédiat.
    \end{itemize}
\end{itemize}

\textbf{4. Conclusion :} Cet extrait est \textbf{très pertinent}. Il démontre une \cle{transférabilité immédiate} des compétences (secteur, outils, normes). Le candidat a tout intérêt à ajouter une rubrique "Profil" en en-tête et à valider son CQP à l'oral.
\end{exoresolu}

\courssub{Méthode 2 : Structurer les éléments constitutifs d'un CV technique}

\begin{methode}[Construire la structure standard d'un CV technique]
\textbf{Objectif :} Assembler les 6 blocs obligatoires + optionnels dans l'ordre logique pour un profil technique.
\begin{enumerate}
    \item \textbf{En-tête (Identité) :} Nom, Prénom, Adresse (Ville, Quartier), Téléphone (WhatsApp pro), Email (prénom.nom@domaine.pro), \cle{LinkedIn / Portfolio} (si à jour). \textbf{Pas de photo obligatoire} au Cameroun, mais recommandée (format identité, fond neutre, tenue pro).
    \item \textbf{Accroche / Titre / Profil (Le "Hook") :} 2--3 lignes. \cle{Titre du poste visé} + \cle{Années d'expérience} + \cle{Compétence phare} + \cle{Disponibilité}. Ex: \emph{"Technicien Maintenance Industrielle (Bac F3 + CQP), 2 ans en agro-industrie (SOSUCAM). Expertise GMAO, HACCP, pneumatique. Disponible immédiatement."}
    \item \textbf{Expériences Professionnelles (Cœur du CV technique) :} Ordre antéchronologique. Pour chaque poste : Dates (Mois/Année), Intitulé exact, Entreprise (Ville), \textbf{3--5 puces "Action + Contexte + Résultat/Chiffre"}. Verbes d'action : \emph{Assurer, Diagnostiquer, Réparer, Programmer, Optimiser, Former}.
    \item \textbf{Formation / Diplômes :} Ordre antéchronologique. Diplôme le plus haut en premier. Intitulé exact, Établissement, Ville, Année, \cle{Mention/Spécialité} (ex: Bac F3 Mention Assez Bien). Stages significatifs intégrés ici ou dans Expériences.
    \item \textbf{Compétences Techniques (Hard Skills) :} Regroupées par catégories : \emph{Maintenance (Préventive, Curative, TPM)}, \emph{Électricité/Automatisme (Schneider, Siemens, Câblage, Lecture plans)}, \emph{Outils/Logiciels (GMAO, CAO, Bureautique)}, \emph{Langues (Niveau CECRL : A1--C2)}.
    \item \textbf{Compétences Transverses / Soft Skills (3 max, prouvées) :} \emph{Rigueur (normes ISO/HACCP)}, \emph{Esprit d'équipe (équipe de 5 techniciens)}, \emph{Réactivité (astreintes week-end)}.
    \item \textbf{Centres d'intérêt / Vie associative (Optionnel mais conseillé) :} Éviter "Lecture, Voyage". Préférer : \emph{Bénévolat Croix-Rouge (secourisme)}, \emph{Football (capitaine équipe)}, \emph{Bricolage électronique (réparation Arduino)}. Montre engagement, leadership, passion technique.
\end{enumerate}
\end{methode}

\begin{exoresolu}[Compléter et ordonner un CV désorganisé]
\textbf{Énoncé :} Mme Ngo Bibaa, titulaire d'un BTS Génie Civil (2023), a effectué un stage de 6 mois chez RAZEL (chef de chantier adjoint) et un job d'été comme géomètre-topographe stagiaire chez GEOCAM. Elle maîtrise AutoCAD, Covadis, MS Project, le français et l'anglais (B2). Elle vise un poste de \emph{Chargée d'études / Dessinatrice projeteuse} chez un promoteur immobilier à Yaoundé. Les éléments ci-dessous sont mélangés. Classez-les dans la structure standard (Méthode 2) et rédigez l'\textbf{Accroche}.

\begin{itemize}
    \item Permis B, voiture personnelle.
    \item Stage RAZEL (Jan--Juin 2023) : Suivi avancement VRD, métrés, relation sous-traitants.
    \item BTS Génie Civil -- Lycée Technique de Douala -- 2023 -- Mention Bien.
    \item AutoCAD (expert), Covadis (avancé), MS Project (intermédiaire), Pack Office.
    \item Job d'été GEOCAM (Juillet--Août 2022) : Relevés topographiques, implantation bornes, traitement données.
    \item Ngo Bibaa Marie, Yaoundé (Mvog-Ada), 6 90 00 00 00, marie.ngo@email.com.
    \item Anglais B2 (TOEIC 780), Français langue maternelle.
    \item Trésorière Association Anciens Élèves (gestion budget 500k FCFA).
    \item Rigueur, Organisation, Communication technique.
\end{itemize}

\tcblower
\textbf{Solution détaillée :}

\textbf{CV Structuré de Mme Ngo Bibaa}

\begin{center}
\begin{tabularx}{\linewidth}{|l|X|}
\hline
\textbf{EN-TÊTE} & \textbf{Ngo Bibaa Marie} \\
& Yaoundé, Mvog-Ada | 6 90 00 00 00 | marie.ngo@email.com \\
& \textit{Permis B, Véhicule personnel} \\ \hline
\textbf{ACCROCHE} & \textbf{Chargée d'études / Dessinatrice Projeteuse (BTS Génie Civil, Mention Bien)} \\
& Stage de fin d'études (6 mois) chez RAZEL : suivi chantier VRD, métrés, coordination. Maîtrise AutoCAD/Covadis/MS Project. Anglais technique B2. Rigueur, organisation, mobilité immédiate. \\ \hline
\textbf{EXPÉRIENCES} & \textbf{01/2023 -- 06/2023 : Stagiaire Chef de Chantier Adjoint -- RAZEL (Yaoundé)} \\
& \textbullet\ Suivi quotidien avancement travaux VRD (voiries, réseaux divers) sur projet 2 km. \\
& \textbullet\ Réalisation métrés hebdomadaires (quantités béton, enrobés) pour facturation : \textbf{zéro écart} vs planning. \\
& \textbullet\ Interface quotidienne avec 3 sous-traitants : planning, sécurité, qualité. \\
& \textbf{07/2022 -- 08/2022 : Stagiaire Géomètre-Topographe -- GEOCAM (Yaoundé)} \\
& \textbullet\ Relevés topographiques station totale (précision $\pm$ 5 mm) pour implantation bornes foncières. \\
& \textbullet\ Traitement données sous Covadis : génération plans de recollement livrés client sous 48h. \\ \hline
\textbf{FORMATION} & \textbf{2023 : BTS Génie Civil -- Lycée Technique de Douala -- \textbf{Mention Bien}} \\
& \emph{Projet de fin d'études :} "Dimensionnement et dessin d'un bâtiment R+2 sous Robot/Autocad". \\ \hline
\textbf{COMPÉTENCES TECHNIQUES} & \textbf{Logiciels CAO/DAO :} AutoCAD (Expert -- 3 ans pratique), Covadis (Avancé -- VRD/Topo), MS Project (Intermédiaire -- Planning). \\
& \textbf{Bureautique :} Pack Office (Excel avancé : TCD, RechercheV pour suivi coûts). \\
& \textbf{Langues :} Français (Maternel), Anglais Technique B2 (TOEIC 780 -- lecture plans/emails pro). \\ \hline
\textbf{SOFT SKILLS} & \textbullet\ \textbf{Rigueur / Normes :} Respect strict DTU / Cahiers charges chantiers RAZEL. \\
& \textbullet\ \textbf{Organisation / Planning :} Gestion simultanée métrés, planning, relation sous-traitants. \\
& \textbullet\ \textbf{Communication technique :} Rédaction rapports de chantier quotidiens, réunions hebdo. \\ \hline
\textbf{VIE ASSOCIATIVE} & \textbullet\ \textbf{Trésorière Association Anciens Élèves LT Douala (2023--2024)} : Gestion budget 500 000 FCFA, comptabilité rigoureuse, rapport AG. \\
& \textbullet\ Volley-ball (équipe mixte -- esprit équipe, régularité). \\ \hline
\end{tabularx}
\end{center}

\textbf{Justifications des choix :}
\begin{itemize}
    \item \textbf{Accroche :} Intègre le titre visé, le diplôme + mention, l'expérience clé (RAZEL), les 3 hard skills critiques (CAO, Covadis, Project), la langue, la soft skill (rigueur) et la dispo.
    \item \textbf{Expériences :} Stage RAZEL (plus pertinent pour le poste visé) mis en premier malgré antéchronologie stricte (stage 2023 vs job 2022) \rightarrow \textbf{Choix stratégique : Pertinence > Chronologie pure}. Puces quantifiées (2 km, zéro écart, 48h, $\pm$ 5mm).
    \item \textbf{Compétences :} Regroupées par catégorie (Logiciels, Langues). Niveaux honnêtes (Expert/Avancé/Intermédiaire).
    \item \textbf{Vie associative :} Trésorière = gestion budgétaire, confiance, rigueur (transférable gestion coûts chantier). Volley = équipe.
\end{itemize}
\end{exoresolu}

\courssub{Méthode 3 : Optimiser la mise en page et les codes iconiques du CV}

\begin{methode}[Soigner la forme : Mise en page, typographie, icônes]
\textbf{Objectif :} Produire un CV lisible en 6 secondes (lecture diagonale recruteur) et compatible ATS (logiciels de tri automatique).
\begin{enumerate}
    \item \textbf{Format \& Fichier :} \cle{PDF} (fige la mise en page), nom du fichier : \texttt{Nom\_Prenom\_PosteVisé.pdf} (ex: \texttt{Kamga\_Jean\_TechnicienMaintenance.pdf}). Une seule page (sauf profil senior 15+ ans).
    \item \textbf{Police \& Couleurs :} Police \textbf{sans empattement} (Sans-serif) : Calibri, Arial, Roboto, Open Sans, Helvetica. Taille 10--11 pt corps, 13--14 pt rubriques, 16 pt Nom. \textbf{Max 2 couleurs} : Noir/Gris anthracite (texte) + \textbf{1 couleur d'accent} (bleu marine, vert sapin, bordeaux) pour les séparateurs, icônes, titres rubriques. Fond blanc.
    \item \textbf{Mise en page (Grille) :} Marges 1,5--2 cm. Colonnes : \textbf{Colonne étroite (30\%)} à gauche pour \emph{Coordonnées, Compétences, Langues, Centres d'intérêt} \textbf{Colonne large (70\%)} à droite pour \emph{Profil, Expériences, Formation}. Guides de lecture visuels (filets verticaux, encadrés légers).
    \item \textbf{Codes iconiques (Pictogrammes) :} Utiliser \textbf{avec parcimonie} (1 par rubrique max) pour guider l'œil.
    \begin{itemize}
        \item \texttt{\textbackslash{}faPhone} / \texttt{\textbackslash{}faEnvelope} / \texttt{\textbackslash{}faMapMarker} / \texttt{\textbackslash{}faLinkedin} (Coordonnées).
        \item \texttt{\textbackslash{}faBriefcase} / \texttt{\textbackslash{}faGraduationCap} / \texttt{\textbackslash{}faTools} / \texttt{\textbackslash{}faLanguage} / \texttt{\textbackslash{}faHeart} (Rubriques).
        \item \texttt{\textbullet} ou \texttt{\textbackslash{}faChevronRight} pour les puces.
    \end{itemize}
    \item \textbf{Espace blanc (Respiration) :} Interligne 1,15. Espace avant/après rubrique (6--8 pt). Alignement gauche impératif (pas de justifié qui crée des "rivières").
    \item \textbf{Accessibilité ATS :} \textbf{Pas de colonnes complexes, pas de tableaux imbriqués, pas de zones de texte, pas d'images de fond.} Texte sélectionnable. Rubriques nommées standard : "Expérience professionnelle", "Formation", "Compétences".
\end{enumerate}
\end{methode}

\begin{exoresolu}[Correction de défauts de mise en page]
\textbf{Énoncé :} Un élève vous remet son CV imprimé pour un stage de fin d'études en Électrotechnique. Vous observez : (1) Police \emph{Comic Sans MS} 12 pt, justifié. (2) Photo de vacances (plage, lunettes soleil) en grand format (4x5 cm) en haut à droite. (3) Fond de page gris foncé avec texte blanc. (4) Rubriques dans le désordre : "Mes Hobbies", "Mon École", "Ce que j'ai fait". (5) Fichier envoyé en .docx. (6) Puces : "Je suis dynamique", "J'aime le travail en équipe", "Je connais le câblage". Liste 5 corrections prioritaires argumentées.

\tcblower
\textbf{Solution détaillée :}

\begin{center}
\begin{tabularx}{\linewidth}{|c|X|X|}
\hline
\textbf{\#} & \textbf{Défaut constaté} & \textbf{Correction + Argumentation technique} \\ \hline
1 & \textbf{Police Comic Sans MS + Justifié} & \textbf{Police : Calibri/Roboto 10,5 pt, Alignement Gauche.} \\
  & & Comic Sans = image "amateur/enfantin", décrédibilise un technicien. Justifié crée des espaces variables (rivières) qui fatiguent l'œil en lecture diagonale. Gauche = rythme de lecture constant. \\ \hline
2 & \textbf{Photo "Vacances" 4x5 cm} & \textbf{Photo identité 3,5x4,5 cm (format standard), fond neutre, tenue pro (polo/chemise), sourire léger.} \\
  & & Photo = 1er contact visuel. Photo plage = manque de sérieux, jugement social négatif immédiat. Taille trop grande = empiète sur l'info utile. Fond neutre = focus sur le visage. \\ \hline
3 & \textbf{Fond gris foncé / Texte blanc} & \textbf{Fond blanc, texte noir/anthracite, accents couleur unique.} \\
  & & \textbf{Incompatible ATS :} Logiciels de lecture (parsing) échouent sur fond sombre/texte clair. Impression coûteuse encre. Lisibilité réduite (fatigue visuelle). Standard pro = fond clair. \\ \hline
4 & \textbf{Rubriques fantaisistes + Désordre} & \textbf{Rubriques standardisées + Ordre : Profil $\rightarrow$ Expériences $\rightarrow$ Formation $\rightarrow$ Compétences.} \\
  & & "Mes Hobbies" $\rightarrow$ "Centres d'intérêt". "Mon École" $\rightarrow$ "Formation". "Ce que j'ai fait" $\rightarrow$ "Expériences". Ordre standard = repérage instantané par recruteur (schéma mental acquis). Mots-clés ATS reconnus. \\ \hline
5 & \textbf{Format .docx} & \textbf{Export PDF (ISO 19005-1 / PDF/A si possible).} \\
  & & .docx = mise en page cassée selon version Word/OS (polices manquantes, sauts de page). PDF = \textbf{fidélité visuelle garantie} 100\%. Nom fichier : \texttt{Ngo\_Bibaa\_Marie\_Stage\_Electrotechnique.pdf}. \\ \hline
6 & \textbf{Puces "Soft" non prouvées + "Je"} & \textbf{Puces "Action + Contexte + Résultat" (STAR), sans "Je", verbes d'action.} \\
  & & "Dynamique" $\rightarrow$ \emph{"Maintenance préventive sur 15 moteurs asynchrones (planification GMAO, 0 panne/6 mois)"}. "Travail équipe" $\rightarrow$ \emph{"Binôme avec chef électricien pour câblage armoires HT/BT (respect NF C 15-100)"}. "Connais câblage" $\rightarrow$ \emph{"Câblage armoires de commande automatismes (Schneider, Siemens), lecture schémas unifilaires"}. \\ \hline
\end{tabularx}
\end{center}

\textbf{Synthèse :} Ces 6 corrections transforment un document "amateur" en outil professionnel \textbf{lisible, parsable (ATS), crédible} et \textbf{percutant}.
\end{exoresolu}

\courssub{Méthode 4 : Maîtriser les valeurs du Subjonctif et du Conditionnel en contexte professionnel}

\begin{methode}[Identifier et employer le Subjonctif et le Conditionnel au CV et à l'oral]
\textbf{Objectif :} Utiliser les modes verbaux appropriés pour exprimer la volonté, l'hypothèse, la politesse, la probabilité dans la rédaction (lettre, mail) et l'entretien.
\begin{enumerate}
    \item \textbf{Le Subjonctif (Subjectivité / Volonté / Nécessité / Doute / Émotion) :} Déclenché par des structures \cle{Il faut que / Il est indispensable que / Je veux que / Je souhaite que / Bien que / Pour que / Avant que} + \textbf{Sujet différent} (généralement).
    \begin{itemize}
        \item \textbf{Valeur Obligation/Nécessité (CV/Lettre) :} \emph{"Il est indispensable que le candidat \textbf{maîtrise} la GMAO."} \emph{"Il faut que je \textbf{sois} opérationnel dès le premier jour."}
        \item \textbf{Valeur Volonté/Souhait (Lettre de motivation/Entretien) :} \emph{"Je souhaite que mon profil \textbf{retienne} votre attention."} \emph{"Je veux que nous \textbf{puissions} échanger sur mes réalisations."}
        \item \textbf{Valeur Concession (Entretien) :} \emph{"Bien que je \textbf{n'aie} pas 5 ans d'expérience, ma formation CQP compense."}
        \item \textbf{Valeur But (Finalité) :} \emph{"Je me forme pour que je \textbf{puisse} évoluer vers la maintenance prédictive."}
    \end{itemize}
    \item \textbf{Le Conditionnel (Hypothèse / Politesse / Atténuation / Information rapportée) :}
    \begin{itemize}
        \item \textbf{Valeur Hypothétique / Conditionnelle (Projet pro) :} \emph{"Si j'intègre votre équipe, je \textbf{m'investirais} pleinement dans la TPM."} (Si + Imparfait $\rightarrow$ Conditionnel Présent).
        \item \textbf{Valeur Politesse / Atténuation (Entretien/Email) :} \emph{"Je \textbf{souhaiterais} vous rencontrer."} (Plus courtois que "Je veux"). \emph{"\textbf{Pouvez}-vous me préciser... ?"} \emph{"J'\textbf{aimerais} savoir..."}
        \item \textbf{Valeur "Journalistique" / Précaution (Évoquer des rumeurs/projets) :} \emph{"L'usine \textbf{recruterait} 10 techniciens cette année."} (Non confirmé).
        \item \textbf{Valeur Futur dans le passé (Récit parcours) :} \emph{"Je savais que je \textbf{réussirais} ce BTS."}
    \end{itemize}
    \item \textbf{PIÈGE À ÉVITER :} \textbf{Indicatif après "Il faut que" / "Je veux que"} $\rightarrow$ \emph{*Il faut que je suis...} (FAUX). \textbf{Subjonctif après "Si"} $\rightarrow$ \emph{*Si j'aurais su...} (FAUX : Si + Imparfait).
\end{enumerate}
\end{methode}

\begin{exoresolu}[Transformation et choix du mode : Lettre de motivation \& Entretien]
\textbf{Énoncé :} Complétez les phrases suivantes en conjuguant le verbe entre parenthèses au \textbf{Subjonctif Présent} ou au \textbf{Conditionnel Présent}. Justifiez chaque choix grammatical (valeur).

1. Il est primordial que le technicien \_\_\_\_\_\_\_\_\_\_ (connaître) les normes de sécurité HACCP.
2. Si j'obtenais ce poste, je \_\_\_\_\_\_\_\_\_\_ (proposer) des améliorations sur la maintenance préventive.
3. Je souhaiterais que vous \_\_\_\_\_\_\_\_\_\_ (m'accorder) un entretien à votre convenance.
4. Bien que le planning \_\_\_\_\_\_\_\_\_\_ (être) chargé, l'équipe a respecté les délais.
5. Le directeur \_\_\_\_\_\_\_\_\_\_ (rechercher) activement un profil comme le mien, d'après mon réseau.
6. Pour que ma candidature \_\_\_\_\_\_\_\_\_\_ (être) complète, je joins mes attestations.
7. \_\_\_\_\_\_\_\_\_\_ (Pouvoir / Vous) me confirmer la date de l'entretien par email ?

\tcblower
\textbf{Solution détaillée :}

\begin{center}
\begin{tabularx}{\linewidth}{|c|l|X|}
\hline
\textbf{N°} & \textbf{Phrase complète} & \textbf{Analyse grammaticale (Mode + Valeur)} \\ \hline
1 & Il est primordial que le technicien \textbf{connaisse} les normes... & \textbf{Subjonctif Présent} (3ème pers. sg). \\
  & & \textbf{Structure déclenchante :} "Il est primordial que" (Impersonnel + Adjectif nécessité). \\
  & & \textbf{Valeur :} \cle{Nécessité / Obligation objective}. Le respect des normes n'est pas négociable. \\ \hline
2 & Si j'obtenais ce poste, je \textbf{proposerais} des améliorations... & \textbf{Conditionnel Présent} (1ère pers. sg). \\
  & & \textbf{Structure :} Protase "Si + Imparfait" (j'obtenais) $\rightarrow$ Apodose Conditionnel. \\
  & & \textbf{Valeur :} \cle{Hypothèse / Projection future conditionnelle}. Action future dépendante d'une condition non réalisée à l'instant T. \\ \hline
3 & Je souhaiterais que vous \textbf{m'accordiez} un entretien... & \textbf{Subjonctif Présent} (2ème pers. pl). \\
  & & \textbf{Structure :} "Je souhaiterais que" (Verbe de volonté/volonté atténuée par conditionnel) + Sujet différent (vous). \\
  & & \textbf{Valeur :} \cle{Volonté / Souhait poli}. Le conditionnel "souhaiterais" adoucit la demande, mais la subordonnée reste au subjonctif. \\ \hline
4 & Bien que le planning \textbf{soit} chargé, l'équipe a respecté les délais. & \textbf{Subjonctif Présent} (3ème pers. sg - verbe être irrégulier). \\
  & & \textbf{Structure :} "Bien que" (Conjonction de concession) + Indicatif réalité (a respecté) en principale. \\
  & & \textbf{Valeur :} \cle{Concession / Opposition}. Le fait (planning chargé) est réel, mais présenté comme un obstacle surmonté. \\ \hline
5 & Le directeur \textbf{recruterait} activement un profil... & \textbf{Conditionnel Présent} (3ème pers. sg). \\
  & & \textbf{Structure :} Verbe d'énonciation "d'après mon réseau" (source indirecte). \\
  & & \textbf{Valeur :} \cle{Information rapportée / Probabilité / "Journalistique"}. Le locuteur n'en est pas témoin direct, il rapporte une rumeur/une info de seconde main avec prudence. \\ \hline
6 & Pour que ma candidature \textbf{soit} complète, je joins mes attestations. & \textbf{Subjonctif Présent} (3ème pers. sg - verbe être). \\
  & & \textbf{Structure :} "Pour que" (Conjonction de but) + Sujet (candidature). \\
  & & \textbf{Valeur :} \cle{But / Finalité}. Action (joindre pièces) visant à réaliser l'état visé (candidature complète). \\ \hline
7 & \textbf{Pourriez}-vous me confirmer la date... ? & \textbf{Conditionnel Présent} (2ème pers. pl - inversion sujet). \\
  & & \textbf{Structure :} Question directe polie. \\
  & & \textbf{Valeur :} \cle{Politesse / Atténuation de la demande}. "Pourriez-vous" (Conditionnel) est la norme épistolaire/orale formelle vs "Pouvez-vous" (Indicatif) plus direct/brut. \\ \hline
\end{tabularx}
\end{center}

\textbf{Rappel formes irrégulières clés :} \emph{Être : que je sois, que tu sois, qu'il soit, que nous soyons, que vous soyez, qu'ils soient.} \emph{Avoir : que j'aie, que tu aies, qu'il ait, que nous ayons, que vous ayez, qu'ils aient.} \emph{Savoir : que je sache...} \emph{Pouvoir : que je puisse...} \emph{Vouloir : que je veuille...} \emph{Faire : que je fasse...} \emph{Aller : que j'aille...}
\end{exoresolu}

\courssub{Méthode 5 : Rédiger les descriptions d'expérience par la méthode STAR}

\begin{methode}[Rédiger des puces d'expérience percutantes : Méthode STAR]
\textbf{Objectif :} Transformer une liste de tâches en preuves de compétences transférables (Action $\rightarrow$ Résultat).
\begin{enumerate}
    \item \textbf{S -- Situation (Contexte) :} Où ? Quand ? Quel environnement ? (Entreprise, service, effectif, enjeu). \emph{Implicite dans la rubrique "Expérience", explicite dans la puce si nécessaire.}
    \item \textbf{T -- Tâche (Mission/Objectif) :} Quel était votre rôle précis ? Quel problème à résoudre ? Quel objectif chiffré ?
    \item \textbf{A -- Action (Votre contribution personnelle) :} Qu'avez-vous \textbf{VOTRE} fait ? Verbe d'action fort à l'infinitif (ou participe passé accordé si "Réalisation de..."). \cle{Verbes techniques :} \emph{Diagnostiquer, Programmer, Câbler, Calibrer, Optimiser, Former, Documenter, Coordonner, Réduire, Augmenter.}
    \item \textbf{R -- Résultat (Résultat mesurable) :} Quel impact ? \textbf{Chiffres, pourcentages, délais, coûts, qualité, sécurité.} \emph{"Réduction de 15\% des arrêts machine", "Zéro accident", "Gain de 2h/jour", "Livraison J-2", "Budget respecté à 98\%".}
    \item \textbf{Synthèse en puce unique :} \emph{[Verbe Action] + [Objet technique/Contexte] + [Moyen/Outils] + [Résultat Chiffré].}
    \item \textbf{Nombre de puces :} 3 à 5 par expérience significative. 1 à 2 pour stages courts/jobs d'été.
\end{enumerate}
\end{methode}

\begin{exoresolu}[Réécriture STAR d'expériences "brutes"]
\textbf{Énoncé :} M. Essomba (Bac F2 Électronique, BTS Maintenance) a noté ses tâches "brutes" pour son stage de 6 mois chez ENEO (Distribution électricité). Réécrivez chaque point en puce STAR (Action + Contexte + Résultat Chiffré).

\begin{enumerate}
    \item "Je faisais la maintenance des postes de transformation."
    \item "J'aidais les chefs d'équipe pour les pannes câbles souterrains."
    \item "Je notais les interventions sur l'ordinateur."
    \item "Je surveillais le stock de fusibles et isolateurs."
    \item "J'ai appris à utiliser un détecteur de défauts."
\end{enumerate}

\tcblower
\textbf{Solution détaillée :}

\begin{center}
\begin{tabularx}{\linewidth}{|c|X|X|}
\hline
\textbf{Brut} & \textbf{Analyse STAR} & \textbf{Puce STAR finale (CV)} \\ \hline
1 & \textbf{S :} Postes HTA/BT (20 kV / 400 V), parc 50 postes. \\
  & \textbf{T :} Maintenance préventive mensuelle (norme ENEO). \\
  & \textbf{A :} \emph{Effectué} visites, contrôles diélectriques, serrage connexions, nettoyage isolateurs, relevés thermographiques. \\
  & \textbf{R :} 50 postes / mois, \textbf{0 défaillance} sur période, rapport conforme 100\%. & \textbf{\textbullet\ Effectué la maintenance préventive mensuelle de 50 postes de transformation HTA/BT (contrôles diélectriques, thermographie, serrage) : 100\% conformité normative, 0 incident imputable maintenance sur 6 mois.} \\ \hline
2 & \textbf{S :} Réseau souterrain urbain, pannes câbles XLPE. \\
  & \textbf{T :} Localisation défauts + réparation (équipe 4 pers). \\
  & \textbf{A :} \emph{Participé} à la localisation par réflectométrie (TDR), \emph{préparé} jonctions/réparations (gainages thermorétractables), \emph{assisté} remise en tension. \\
  & \textbf{R :} 12 interventions, \textbf{délai moyen remise service < 4h} (cible < 6h). & \textbf{\textbullet\ Participé à la localisation (réflectométrie TDR) et réparation de 12 défauts câbles souterrains XLPE (jonctions thermorétractables) : temps moyen de rétablissement 3h45 (objectif < 6h).} \\ \hline
3 & \textbf{S :} GMAO ENEO (logiciel SAP PM / Maximo). \\
  & \textbf{T :} Traçabilité exhaustive interventions (légalité, historique, coûts). \\
  & \textbf{A :} \emph{Saisi} et \emph{clôturé} ordres de travail (OT) en temps réel : comptes-rendus, pièces changées, temps passé, photos. \\
  & \textbf{R :} \textbf{100\% OT clôturés} à J+1 max, données fiables pour KPI direction. & \textbf{\textbullet\ Assuré la traçabilité GMAO (SAP PM) de l'activité maintenance : saisie/clôture 100\% des Ordres de Travail à J+1 (comptes-rendus, consommations, photos), fiabilisant les KPI de performance.} \\ \hline
4 & \textbf{S :} Magasin maintenance centrale, 500 références. \\
  & \textbf{T :} Disponibilité pièces critiques (fusibles HT, isolateurs, huiles). \\
  & \textbf{A :} \emph{Surveillé} seuils d'alerte (stock mini), \emph{lancé} demandes achat (DA), \emph{réceptionné/vérifié} conformité livraisons. \\
  & \textbf{R :} \textbf{0 rupture stock} pièces critiques, \textbf{-12\%} valeur stock dormant (inventaire tournant). & \textbf{\textbullet\ Géré le stock pièces critiques (fusibles HT, isolateurs) : surveillance seuils, lancement DA, réception qualité $\rightarrow$ 0 rupture sur pièces critiques, optimisation stock dormant -12\% via inventaire tournant.} \\ \hline
5 & \textbf{S :} Formation interne ENEO / Fabricant (Megger/BAUR). \\
  & \textbf{T :} Montée en compétence diagnostic câbles. \\
  & \textbf{A :} \emph{Suivi} formation certifiante "Localisation défauts câbles HTA" (théorie + terrain), \emph{manipulé} générateur d'impulsions / récepteur. \\
  & \textbf{R :} \textbf{Certificat obtenu}, \emph{autonome} sur appareil pour pré-localisation. & \textbf{\textbullet\ Validé la certification interne "Localisation défauts câbles HTA" (théorie + pratique terrain sur générateur d'impulsions Megger) : autonomie acquise sur pré-localisation défauts, certificat joint au dossier.} \\ \hline
\end{tabularx}
\end{center}

\textbf{Apport de la méthode STAR :} Le recruteur voit \textbf{quoi} (technique précis), \textbf{comment} (outils/méthodes), \textbf{combien} (volume), \textbf{bien} (résultat). C'est la \cle{preuve} de compétence.
\end{exoresolu}

\courssub{Méthode 6 : Présenter et défendre son CV à l'oral (Entretien)}

\begin{methode}[Structurer sa prise de parole : "Parlez-moi de vous / Présentez votre parcours" (2--3 min)]
\textbf{Objectif :} Raconter son parcours comme un \textbf{récit cohérent} (Storytelling pro) liant passé $\rightarrow$ présent $\rightarrow$ futur (le poste), et non lire son CV.
\begin{enumerate}
    \item \textbf{Accroche (30 sec) -- "Qui suis-je ?" :} \cle{Identité pro + Cible + Valeur ajoutée}. \emph{"Bonjour, je m'appelle X, je suis Technicien Maintenance (Bac F3 + CQP), 2 ans d'expérience en agro-industrie chez SOSUCAM. Ma force : l'autonomie sur GMAO et la rigueur HACCP. Je suis là aujourd'hui car votre poste correspond exactement à mon cœur de métier."}
    \item \textbf{Parcours clé (1 min 30) -- "Ce que j'ai fait qui prouve que je sais faire" :} Choisir \textbf{2--3 expériences majeures} (pas toutes). Pour chacune : \textbf{Contexte $\rightarrow$ Défi/Projet marquant $\rightarrow$ Mon Action (STAR) $\rightarrow$ Résultat/Apprentissage}. \emph{Transition :} "Cette expérience m'a permis de développer X, que j'ai approfondi ensuite chez Y..."
    \item \textbf{Projet / Motivation (30 sec) -- "Pourquoi VOUS ? Pourquoi MOI ?" :} Lien explicite : \emph{"Votre usine lance la TPM / modernise la ligne X. Mon expérience sur la GMAO / la maintenance préventive / l'automatisme Siemens me permet d'être opérationnel vite. Je veux m'investir durablement dans cette dynamique."}
    \item \textbf{Clôture (15 sec) :} \emph{"Je suis à votre disposition pour détailler tel projet technique ou faire un cas pratique. Merci de votre écoute."}
    \item \textbf{Techniques orales :} Regard (contact visuel), Voix (débit modéré, articulation), Posture (ouvert, stable), \textbf{Silence} (respirer entre les blocs), \textbf{Mots de liaison} (\emph{Par conséquent, En outre, C'est dans ce contexte que...}). \textbf{Interdit :} "Euh", "Du coup", "Genre", lire ses notes, critiquer ex-employeurs.
\end{enumerate}
\end{methode}

\begin{exoresolu}[Simulation : Réponse à "Parlez-moi de vous" pour un poste de Technicien Bureau d'Études Élec.]
\textbf{Énoncé :} M. Fouda (BTS Électrotechnique, 3 ans expé : 1 an câbleur armoires chez SOCAVER, 2 ans dessinateur projeteur chez BET INGECAM) passe un entretien pour un poste de \emph{Technicien BE Électricité} chez un grand promoteur (immeubles R+10). Il a 2 min 30. Rédigez son script oral structuré (Méthode 6) en intégrant : \emph{Logiciels : See Electrical, Caneco, AutoCAD. Normes : NF C 15-100. Projet marquant : Hôpital de district (courants forts/faibles, GTB). Soft skills : Rigueur normative, coordination corps d'état.}

\tcblower
\textbf{Solution détaillée : Script Oral (environ 220 mots = 2 min 15 à l'oral)}

\textbf{1. ACCROCHE (25 sec)}
« Bonjour. Je m'appelle Fouda Jean, 26 ans, Technicien Bureau d'Études Électricité (BTS Électrotechnique). J'ai 3 ans d'expérience terrain et conception. Mon fil conducteur : \textbf{la rigueur normative NF C 15-100} appliquée de la réalisation (câblage) à la conception (schémas, calculs). Aujourd'hui, je vise ce poste chez vous car vos projets d'immeubles de grande hauteur demandent exactement cette double compétence exécution/conception que j'ai construite. »

\textbf{2. PARCOURS CLÉ : 2 EXPÉRIENCES STRUCTURANTES (1 min 30)}

\textbf{Expérience 1 -- Le "Terrain" (1 an -- SOCAVER, Câbleur Armoires) :}
« Départ chez SOCAVER, câbleur armoires HT/BT. \textbf{Contexte :} Série de 50 armoires process pour usine chimique. \textbf{Défi :} Respect délais serrés + normes strictes zone ATEX. \textbf{Mon action :} J'ai \emph{revu le processus de câblage} (repérage fils, torons, tests diélectriques) et \emph{formé 2 intérimaires} aux points de contrôle critiques. \textbf{Résultat :} \textbf{Zéro non-conformité} à la recette client, \textbf{-15\% temps câblage} par armoire. \emph{Apport :} Je connais la réalité du terrain, la place des doigts, les erreurs de lecture de schéma. »

\textbf{Expérience 2 -- La "Conception" (2 ans -- BET INGECAM, Dessinateur Projeteur) :}
« Bascule en BE chez INGECAM. \textbf{Projet phare :} \textbf{Hôpital de district (R+3, 15 000 m²)}. Courants forts (TGBT, groupes, secours) + Courants faibles (SSI, VDI, GTB KNX). \textbf{Mon rôle :} \emph{Conception complète} sous \textbf{See Electrical / Caneco BT} (calculs sections, protections, sélectivité), \emph{plans d'implantation / cheminements} sous AutoCAD, \emph{dossiers consultation/entreprises} (CCTP, DPGF). \textbf{Coordination :} \textbf{Réunions hebdo} avec archi, fluides, structure pour \emph{réserver les gaines/emprises} (conflits résolus en amont). \textbf{Résultat :} \textbf{0 réserve technique} à la réception provisoire, \textbf{validation Consuel} du 1er coup. »

\textbf{3. PROJET / MOTIVATION (20 sec)}
« Votre promoteur lance des tours R+10. La complexité : \textbf{GTB centralisée, secours sécurité, sélectivité verticale}. Ma maîtrise \textbf{Caneco (calculs courts-circuits, sélectivité)} + mon \textbf{reflexe terrain (gainages, accessibilité maintenance)} me permettent de produire des dossiers \textbf{"constructibles" et "maintenables"} dès la phase PRO/DCE. Je veux m'inscrire dans la durée sur ce type de projets. »

\textbf{4. CLÔTURE (10 sec)}
« Je suis prêt à détailler le dimensionnement du TGBT hôpital ou un cas pratique sur la sélectivité. Merci. »

\textbf{Analyse qualité :} \textbf{Temps respecté.} \textbf{Structure claire} (Accroche $\rightarrow$ 2 blocs STAR $\rightarrow$ Projet $\rightarrow$ Clôture). \textbf{Mots-clés techniques} (NF C 15-100, See Electrical, Caneco, GTB, KNX, Consuel, Sélectivité, CCTP). \textbf{Preuves chiffrées} (50 armoires, -15\%, 15 000 m², 0 réserve, 1er coup). \textbf{Lien explicite} avec la cible (R+10, GTB, sélectivité verticale). \textbf{Ton professionnel, assertif, pas arrogant.}
\end{exoresolu}

\courssub{Méthode 7 : Intégration -- Simulation complète : De l'offre à l'entretien}

\begin{methode}[Démarche intégrée : Répondre à une offre réelle (CV + Lettre + Préparation orale)]
\textbf{Objectif :} Traiter une situation complète d'insertion professionnelle en mobilisant toutes les compétences de la leçon.
\begin{enumerate}
    \item \textbf{Analyser l'offre (Déconstruction) :} Surligner : \cle{Mots-clés techniques} (logiciels, machines, normes), \cle{Compétences comportementales} (autonomie, équipe, rapport), \cle{Conditions} (lieu, horaires, permis, salaire). Identifier le \textbf{problème caché} du recruteur (remplacement ? croissance ? compétence manquante ?).
    \item \textbf{Adapter le CV "Master" (Ciblage) :}
    \begin{itemize}
        \item Modifier \textbf{Titre/Accroche} avec mots-clés offre.
        \item \textbf{Réordonner/Seleccionner} expériences : Mettre en 1er celle qui "matche" le plus. Supprimer ou réduire les sans lien.
        \item \textbf{Injecter mots-clés} dans puces STAR (ex: si offre demande "TPM", ajouter "TPM" dans puce maintenance).
        \item Vérifier \textbf{Compétences Techniques} : Niveaux honnêtes, logiciels cités dans l'offre en gras.
    \end{itemize}
    \item \textbf{Rédiger Lettre/Email d'accompagnement (Structure "Vous / Moi / Nous") :}
    \begin{itemize}
        \item \textbf{Vous :} Accroche sur l'entreprise (valeurs, projet, actualité). \emph{"Votre projet d'extension de l'unité de production..."}
        \item \textbf{Moi :} 2--3 arguments \textbf{preuve} (STAR condensé) répondant \textbf{directement} aux besoins clés de l'offre. \emph{"Mon expérience de 2 ans sur ligne d'embouteillage Krones (36 000 b/h) me permet d'être opérationnel immédiatement sur votre nouvelle ligne."}
        \item \textbf{Nous :} Projection + Appel à l'action. \emph{"Je souhaite vous rencontrer pour échanger sur ma contribution à la fiabilisation de votre outil. Disponible pour un entretien à votre convenance."}
        \item \textbf{Grammaire :} Conditionnel politesse (\emph{je souhaiterais, je pourrais}), Subjonctif nécessité (\emph{il est important que je puisse}).
    \end{itemize}
    \item \textbf{Préparer l'entretien (Fiche de révision "Anti-stress") :} 1 page recto :
    \begin{itemize}
        \item Mon \textbf{Pitch 2 min} (Méthode 6).
        \item \textbf{3 Réalisations STAR} détaillées** (prêtes pour questions comportementales : "Parlez-moi d'une panne difficile...", "Un conflit...").
        \item \textbf{3 Questions à poser} au recruteur (intérêt, culture, technique) : \emph{"Comment est organisée la maintenance préventive aujourd'hui ?", "Quels sont les indicateurs de performance (KPI) suivis ?", "Quelles sont les perspectives d'évolution technique ?"}.
        \item \textbf{Points de vigilance CV} (trous, changements, niveau anglais) $\rightarrow$ Réponses préparées honnêtes/positives.
    \end{itemize}
    \item \textbf{Post-entretien :} Email de remerciement J+1 (rappel 1 point fort + motivation confirmée).
\end{enumerate}
\end{methode}

\begin{exoresolu}[Cas intégré : Réponse à une offre réelle -- Technicien Frigoriste]
\textbf{Énoncé :} Offre d'emploi : \emph{"Entreprise de froid industriel/clim (Yaoundé) recrute Technicien Frigoriste H/F. Missions : Maintenance préventive/curative installations NH3 / CO2 / Freons (chambres froides, tunnels congélation, climatisation centrale). Diagnostic pannes, recharges, soudure brasage fort. GMAO. Déplacements sites clients. Permis B. Bac Pro / BTS Froid / CQP. 2 ans expé. Autonome, rigueur sécurité, relation client."}

Profil candidat : M. Atangana, 28 ans, \textbf{BTS Froid Conditionnement d'Air (2020)}. Expé : 2021-2023 \emph{Technicien itinérant} chez FROID SERVICE (clim résidentiel/tertiaire R410A/R32, peu froid industriel). 2023-2024 \emph{Technicien maintenance} chez BRASSERIES DU CAMEROUN (froid industriel NH3/CO2, chambres froides, tunnels, GMAO SAP PM, astreintes). Permis B. Anglais technique notice. Soudeur brasage fort certifié.

\textbf{Travail demandé :}
1. Liste des 5 mots-clés techniques + 2 comportementaux de l'offre.
2. Rédigez l'\textbf{Accroche CV} ciblée.
3. Rédigez \textbf{2 puces STAR} pour l'expérience Brasseries (la plus pertinente).
4. Rédigez le paragraphe \textbf{"Moi" de la lettre} (3 arguments preuve).
5. Donnez \textbf{2 questions techniques} à poser au recruteur.

\tcblower
\textbf{Solution détaillée :}

\textbf{1. Analyse Offre -- Mots-clés :}
\begin{itemize}
    \item \textbf{Techniques :} \cle{NH3 (Ammoniac) / CO2}, \cle{Chambres froides / Tunnels congélation}, \cle{Soudure brasage fort}, \cle{GMAO (SAP PM)}, \cle{Diagnostic / Recharge / Maintenance préventive \& curative}.
    \item \textbf{Comportementaux :} \cle{Autonomie (itinérante/astreintes)}, \cle{Relation client / Rigueur sécurité}.
\end{itemize}

\textbf{2. Accroche CV Ciblée :}
\textbf{Technicien Frigoriste (BTS Froid, 4 ans expé) -- Spécialiste Froid Industriel NH3/CO2 (Brasseries).} Expertise maintenance préventive/curative chambres froides \& tunnels congélation, diagnostic avancé, brasage fort certifié, GMAO SAP PM. Autonome en itinérance/astreintes, rigueur sécurité (ATEX/NH3), relation client. Permis B. \textbf{Disponible immédiatement.}

\textbf{3. 2 Puces STAR Brasseries (Expérience 2023-2024 -- La plus pertinente) :}
\begin{itemize}
    \item \textbf{\textbullet\ Assuré la maintenance préventive (plan GMAO SAP PM : 120 interventions/mois) et curative d'un parc froid industriel NH3/CO2 (chambres -25°C, tunnels -35°C, clim centrale) : \textbf{Disponibilité installation > 98,5\%}, \textbf{MTTR < 1h30} sur pannes critiques production.}
    \item \textbf{\textbullet\ Réalisé en autonomie diagnostics complexes (analyse vibratoire, thermographie, analyse huile), recharges fluides, \textbf{brasage fort cuivre/acier certifié} (réparation évaporateurs, condenseurs) : \textbf{0 fuite récurrente}, \textbf{conformité réglementaire ICPE/ATEX} validée inspection annuelle.}
\end{itemize}

\textbf{4. Paragraphe "Moi" Lettre de motivation :}
« \textbf{Votre besoin d'autonomie sur installations NH3/CO2 trouve réponse dans mon expérience actuelle chez Brasseries du Cameroun.} Depuis 18 mois, j'assure seul la maintenance d'un parc critique (chambres, tunnels, NH3/CO2) avec une disponibilité > 98,5\% (KPI suivi SAP PM). \textbf{Ma certification brasage fort et ma pratique quotidienne du diagnostic (thermographie, analyse huile)} me permettent d'intervenir en première intention sur tout type de panne, y compris en astreinte week-end. \textbf{Enfin, mon passage antérieur en clim tertiaire (Froid Service) m'a forgé la relation client et la gestion d'interventions multi-sites,} atout pour votre activité itinérante chez les clients. »

\textbf{5. 2 Questions techniques à poser (Montrent expertise + projection) :}
\begin{enumerate}
    \item \emph{"Quelle est la répartition de votre parc entre fluides NH3, CO2 et HFC (R404A/R448A) ? Avez-vous des installations en transcritique CO2 ?"} (Montre connaissance technique pointue, vocabulaire pro).
    \item \emph{"Comment est gérée la traçabilité réglementaire (registres fluides, contrôles périodiques ICPE, ATEX) dans votre GMAO ? Est-ce intégré aux gammes de maintenance préventive ?"} (Montre rigueur normative, connaissance GMAO avancée, anticipation charge admin).
\end{enumerate}
\end{exoresolu}

\courssub{Les 5 erreurs qui coûtent des points au Bac (et en vrai)}

\begin{attention}[Les 5 erreurs fatales au Probatoire / Bac Technique -- Français : CV \& Entretien]
\begin{enumerate}
    \item \textbf{CV "Passe-partout" non ciblé :} Envoyer le \textbf{même CV} pour toutes les offres sans modifier l'Accroche, l'ordre des expériences, ni injecter les \cle{mots-clés} de l'offre. \textbf{Conséquence :} Rejet ATS (logiciel) ou lecture humaine < 6 sec = poubelle. \textbf{Remède :} 1 CV Master $\rightarrow$ 1 CV Ciblé par offre (Méthode 7, étape 2).
    \item \textbf{Puces "Tâches" au lieu de "Réalisations STAR" :} Écrire \emph{"Maintenance machines", "Gestion stock", "Rédaction rapports"} sans \textbf{verbe d'action fort, contexte technique précis, chiffre/résultat}. \textbf{Conséquence :} Le recruteur ne voit \textbf{pas votre valeur ajoutée}, ni votre niveau d'autonomie. \textbf{Remède :} Systématiser \textbf{Action + Contexte + Résultat Chiffré} (Méthode 5).
    \item \textbf{Fautes de français / Syntaxe / Modes verbaux :} Fautes d'orthographe (impardonnable), phrases mal construites, \textbf{mauvais choix Subjonctif/Conditionnel} (ex: *\emph{Il faut que je suis}..., *\emph{Si j'aurais su}...). \textbf{Conséquence :} Signal fort de \textbf{manque de rigueur} et de niveau langue insuffisant pour la communication pro (rapports, mails, sécurité). \textbf{Remède :} Relire à voix haute, correcteur, réviser Méthode 4 (valeurs des modes).
    \item \textbf{Présentation orale = Lecture du CV :} Réciter les dates et intitulés ligne par ligne. \textbf{Conséquence :} Ennui du jury, absence de \textbf{storytelling}, pas de lien avec le poste, temps dépassé. \textbf{Remède :} Structure \textbf{Accroche $\rightarrow$ 2 Expériences STAR choisies $\rightarrow$ Projet/Motivation $\rightarrow$ Clôture} (Méthode 6). S'entraîner à voix haute (chrono).
    \item \textbf{Impréparation à l'entretien (Zéro question, Zéro connaissance entreprise) :} Ne pas savoir citer un chiffre clé de l'entreprise, son actualité, ses produits. Ne pas avoir de questions techniques pertinentes à poser. \textbf{Conséquence :} Perçu comme \textbf{passif, peu motivé, peu curieux}. \textbf{Remède :} Fiche "Anti-stress" 1 page (Méthode 7, étape 4) : Pitch, 3 STAR détaillés, 3 Questions pro, Réponses points faibles.
\end{enumerate}
\end{attention}

\newpage

\coursec{Exercices}

\coursec{Leçon 24 : Présenter son CV dans le cadre d'un entretien d'embauche}

\courssub{Je vérifie mes connaissances}

\begin{exercice}[QCM : Les éléments du CV]
Parmi les propositions suivantes, cochez celle qui correspond à la rubrique \textbf{obligatoire} dans un CV moderne pour un poste technique au Cameroun.
\begin{enumerate}
    \item La photo d'identité en costume traditionnel.
    \item L'état civil complet (nom, prénoms, date et lieu de naissance, nationalité, situation matrimoniale, adresse, téléphone, email).
    \item La liste détaillée des loisirs (football, musique, lecture).
    \item Le salaire mensuel souhaité en FCFA.
\end{enumerate}
\end{exercice}

\begin{exercice}[Vrai ou Faux justifié : Subjonctif et Conditionnel]
Déterminez si les affirmations suivantes sont vraies ou fausses. Justifiez votre réponse par une phrase explicative.
\begin{enumerate}
    \item Dans une lettre de motivation, « Je \textbf{veuille} que vous \textbf{examiniez} ma candidature » : le verbe \emph{examiniez} est au subjonctif présent car il exprime un fait certain.
    \item « Je \textbf{souhaiterais} obtenir un entretien » : le verbe \emph{souhaiterais} est au conditionnel présent et marque la politesse et l'atténuation.
    \item « Il faut que tu \textbf{sois} ponctuel » : \emph{sois} est à l'indicatif présent car c'est une obligation réelle.
    \item « Si j'\textbf{avais} le temps, je \textbf{réviserais} mon CV » : cette phrase exprime une hypothèse réalisable dans l'avenir immédiat.
\end{enumerate}
\end{exercice}

\begin{exercice}[Définitions et concepts clés]
Donnez une définition précise (2 à 3 lignes) pour chacun des termes suivants dans le contexte de la recherche d'emploi :
\begin{enumerate}
    \item Le \cle{CV (Curriculum Vitae)}.
    \item La \cle{lettre de motivation}.
    \item L'\cle{entretien d'embauche}.
    \item Le \cle{texte iconique} appliqué au CV (mise en page, police, icônes).
\end{enumerate}
\end{exercice}

\begin{exercice}[Application directe : Morphosyntaxe]
Conjuguez les verbes entre parenthèses au mode et au temps appropriés (Subjonctif présent ou Conditionnel présent) pour compléter cet extrait de lettre de motivation.
\begin{enumerate}
    \item Je souhaite que vous \rule{3cm}{0.4pt} (examiner) ma candidature avec attention.
    \item Il est indispensable que le candidat \rule{3cm}{0.4pt} (maîtriser) les logiciels comptables.
    \item Je \rule{3cm}{0.4pt} (apprécier) de pouvoir vous exposer mes motivations de vive voix.
    \item Il faut que je \rule{3cm}{0.4pt} (joindre) mes diplômes à ce courrier.
    \item Nous \rule{3cm}{0.4pt} (être) honorés de collaborer avec votre structure.
\end{enumerate}
\end{exercice}

\courssub{Je m'entraîne}

\begin{exercice}[Complétion guidée : Lettre de motivation à trous]
Complétez la lettre ci-dessous en conjuguant les verbes entre parenthèses au \textbf{subjonctif} ou au \textbf{conditionnel} selon le sens (volonté, nécessité, politesse, hypothèse). \textit{10 verbes à conjuguer.}

\textbf{Objet : Candidature au poste de Technicien en Maintenance Industrielle}

Monsieur le Directeur,

Par la présente, je vous adresse ma candidature pour le poste de Technicien en Maintenance publié dans \emph{Le Jour} du 10 octobre. Je \rule{2cm}{0.4pt} (souhaiter) vivement intégrer votre équipe. Titulaire d'un BTS en Maintenance des Systèmes Automatisés, je \rule{2cm}{0.4pt} (aimer) mettre mes compétences à votre service.

Il est primordial que je \rule{2cm}{0.4pt} (pouvoir) démontrer mon savoir-faire lors d'un entretien. Je \rule{2cm}{0.4pt} (être) reconnaissant de me recevoir à votre convenance. Il faut que vous \rule{2cm}{0.4pt} (recevoir) mon dossier complet. Afin que mon dossier \rule{2cm}{0.4pt} (être) complet, je joins mes certificats de travail.

Je \rule{2cm}{0.4pt} (vouloir) vous rencontrer pour détailler mon parcours. Nous \rule{2cm}{0.4pt} (pouvoir) convenir d'une date ensemble. Dans l'attente d'une réponse que j'\rule{2cm}{0.4pt} (espérer) favorable, je vous prie d'agréer, Monsieur le Directeur, l'expression de mes salutations distinguées.

\textit{Signature}
\end{exercice}

\begin{exercice}[Réception / Analyse comparative : Deux CV pour un poste de Secrétaire Comptable]
Voici les profils synthétiques de deux candidates pour un poste de Secrétaire Comptable à Yaoundé (exigeant : rigueur, maîtrise d'Sage/Excel, discrétion, 2 ans exp.).

\textbf{Candidate A (Marie K.)} : BTS Comptabilité (2021). Expérience : 3 ans comme Comptable Unique dans une PME (gestion facturation, déclarations fiscales, paie). Maîtrise parfaite Sage 100, Excel avancé (TCD). Loisirs : Lecture, Bénévolat association caritative. Mise en page : sobre, police Calibri 11, icônes discrètes, photo pro.

\textbf{Candidate B (Grace N.)} : DUT Gestion Administrative (2022). Expérience : 6 mois Stage Secrétariat (accueil, standard, classement). 1 an Vendeuse en boutique. Maîtrise Word, Excel débutant. Connaissances comptables théoriques. Loisirs : Danse, Sorties entre amis, Réseaux sociaux. Mise en page : colorée, police fantaisie, photo selfie, icônes nombreuses.

\textbf{Consigne :} Rédigez un texte argumenté d'environ \textbf{15 lignes} justifiant le choix de la candidate à retenir pour le poste. Appuyez-vous sur les critères du poste (compétences techniques, expérience, savoir-être, présentation du CV).
\end{exercice}

\begin{exercice}[Correction et réécriture : Un CV truffé d'erreurs]
Le CV ci-dessous contient \textbf{10 erreurs} (orthographe, grammaire modes/temps, chronologie inversée, mise en page, informations inutiles/manquantes). Identifiez-les, corrigez-les et réécrivez le CV proprement sur une page.

\begin{center}
\begin{tabularx}{\linewidth}{|X|}
\hline
\textbf{CURRICULUM VITAE} \\
\hline
\textbf{NOM :} mbarga jean-pierre \quad \textbf{AGE :} 25 ans \\
\textbf{ADRESSE :} Quartier Mvog-Ada, Yaoundé \quad \textbf{Tel :} 6XX XX XX XX \\
\textbf{EMAIL :} jeanpierre@yahoo.fr \quad \textbf{SITUATION :} Célibataire, 2 enfants \\
\hline
\textbf{FORMATION :} \\
2023 -- 2021 : BTS Électrotechnique (Lycée Technique de Douala) \textit{(Mention AB)} \\
2018 -- 2015 : Baccalauréat F3 (Lycée de Nkolbisson) \\
\hline
\textbf{EXPÉRIENCES PROFESSIONNELLES :} \\
2024 -- Présent : Stagiaire Maintenance à SONEL (Douala) -- \textit{Si j'aurais su, je ferait mieux.} \\
2022 -- 2023 : Apprenti Électricien chez M. Toto (Yaoundé) -- \textit{Il faut que je suis sérieux.} \\
\hline
\textbf{COMPÉTENCES :} \\
Câblage, Dépannage, Lecture plans, \textbf{Anglais : Lu, Écrit, Parlé (Niveau BAC)} \\
\textbf{LOISIRS :} Football, Boire des verres, Dormir, Politique \\
\hline
\textbf{RÉFÉRENCES :} M. Le Directeur de SONEL (Tel: 6XX XX XX XX) \\
\hline
\end{tabularx}
\end{center}
\end{exercice}

\begin{exercice}[Simulation orale : « Présentez-vous » (2 minutes)]
Préparez la structure de votre réponse orale à la question classique d'entretien : \emph{« Bonjour, présentez-vous en deux minutes. »}
\begin{enumerate}
    \item Rédigez le \textbf{plan détaillé} de votre intervention (Accroche, Parcours formation, Expériences clés, Compétences techniques/transverses, Projet professionnel, Conclusion/Call to action).
    \item Indiquez pour chaque partie le \textbf{temps approximatif} alloué (total 120 secondes).
    \item Listez \textbf{3 pièges à éviter} (ex: réciter le CV, vie privée, hésitations).
\end{enumerate}
\textit{Cet exercice sera évalué par les pairs avec une grille d'évaluation (Voix, Regard, Structure, Pertinence, Langue).}
\end{exercice}

\courssub{Je me prépare au Bac}

\begin{exercice}[Sujet Type Baccalauréat Technique : Production d'un CV complet]
\textbf{Situation :} Vous êtes \textbf{Paul BIKOI}, titulaire d'un \textbf{BTS en Génie Civil (Option Bâtiment)} obtenu en 2023 au Lycée Technique de Bafoussam. Vous avez effectué un stage de 6 mois (Janv.--Juin 2024) à l'entreprise \textbf{``BTP EXCELLENCE''} (Douala) comme Chef de Chantier Adjoint (suivi exécution, gestion matériaux, sécurité). Vous parlez couramment Français et Anglais (niveau B2), maîtrisez AutoCAD, ArchiCAD, MS Project. Permis B. Célibataire, 24 ans, résident à Douala (Bonabéri).

Vous répondez à l'offre ci-dessous parue dans \emph{Mutations} du 15 Novembre 2024 :

\begin{quote}
\textbf{ENTREPRISE ``CONSTRUCT FUTURE'' (Yaoundé) RECRUTE}
\textbf{POSTE :} Conducteur de Travaux Junior (H/F)
\textbf{PROFIL :} BTS/DUT Génie Civil. Expérience stage/premier emploi acceptée. Maîtrise DAO (AutoCAD) et Planning (MS Project). Anglais technique lu/écrit. Rigueur, leadership, mobilité nationale.
\textbf{DOSSIER :} CV + LM + Copies diplômes. Email : rh@constructfuture.cm
\end{quote}

\textbf{Travail demandé :}
\begin{enumerate}
    \item Rédigez le \textbf{CV complet, tenu sur une page}, adapté à cette offre (Mise en page soignée, rubriques ordonnées, mots-clés de l'offre intégrés, photo mentionnée, pas de fautes).
    \item Soulignez dans la marge (ou par annotation) \textbf{3 choix stratégiques} que vous avez faits pour cibler ce poste précis (ex: ordre des rubriques, vocabulaire technique, mise en valeur stage).
\end{enumerate}
\end{exercice}

\begin{exercice}[Sujet Type Probatoire Technique : Grammaire en contexte + Production]
\textbf{Partie A : Transformation grammaticale (8 points)}
Transformez les phrases suivantes de l'indicatif vers le \textbf{conditionnel de politesse} (ou forme atténuée) adaptées à un entretien d'embauche formel.
\begin{enumerate}
    \item Je \textbf{veux} savoir si le poste est toujours vacant.
    \item Je \textbf{peux} commencer dès lundi prochain.
    \item Il \textbf{faut} que je parte tôt vendredi pour un examen.
    \item Je \textbf{pense} que je suis le meilleur candidat.
    \item Nous \textbf{allons} signer le contrat la semaine prochaine.
    \item Vous \textbf{devez} me donner une réponse rapide.
    \item Je \textbf{sais} que l'entreprise a de grands projets.
    \item Je \textbf{viens} pour l'entretien à 10h.
\end{enumerate}

\textbf{Partie B : Rédaction de la lettre de motivation (12 points)}
En vous servant des verbes transformés en Partie A et du profil de l'Exercice 9 (Paul Bikoi, Conducteur de Travaux Junior), rédigez la \textbf{lettre de motivation} formelle accompagnant le CV. Respectez la structure administrative (En-tête, Objet, Formules d'appel, Corps structuré en 3 paragraphes : Qui suis-je ? Pourquoi vous ? Pourquoi moi ? Proposition entretien, Formules de politesse, Signature). \textit{Longueur : 20 à 25 lignes.}
\end{exercice}

\begin{exercice}[Sujet Type Bac : Intégration - Analyse de jury et Synthèse]
Après l'entretien de Paul Bikoi (Exercice 9), le jury de recrutement (DRH, Chef Projet, Conducteur Principal) remplit une grille d'évaluation orale sur 20 points.

\textbf{Grille d'évaluation (Extrait) :}
\begin{center}
\begin{tabularx}{\linewidth}{|l|X|c|}
\hline
\textbf{Critères} & \textbf{Indicateurs observés pour Paul} & \textbf{Note /20} \\
\hline
Présentation / Apparence & Costume propre, posture droite, contact visuel franc, sourire & \\
\hline
Structure du discours & Plan clair (Parcours $\to$ Stage $\to$ Compétences $\to$ Projet), transitions fluides, temps respecté (1'55) & \\
\hline
Compétences Techniques & Réponses précises sur DAO, Béton, Planning, Sécurité. Exemple concret stage (gestion conflit fournisseur). & \\
\hline
Compétences Transverses & Anglais technique lu (notice), aisance relationnelle, gestion stress (silence 3s avant réponse difficile). & \\
\hline
Motivation / Adéquation & Connaissance entreprise (projets cités), mobilité acceptée, questions pertinentes posées au jury (formation continue). & \\
\hline
Langue / Expression & Français standard soutenu, vocabulaire technique précis, pas de tics de langage, conditionnel de politesse utilisé. & \\
\hline
\end{tabularx}
\end{center}

\textbf{Travail demandé :}
\begin{enumerate}
    \item En vous mettant à la place du \textbf{DRH}, rédigez le \textbf{compte rendu synthétique de l'entretien} (environ 15 lignes) destiné au Directeur Général pour avis de recrutement. Structure : Introduction (candidat, poste), Points forts (3), Points de vigilance / Axes de progression (2), Avis motivé (Oui/Non/Réserve).
    \item Sur le plan linguistique, citez \textbf{3 exemples précis} tirés de l'oral supposé de Paul qui illustrent l'usage correct du \textbf{subjonctif} (nécessité, souhait) et du \textbf{conditionnel} (politesse, hypothèse) lors de l'échange.
\end{enumerate}
\end{exercice}

\newpage

\coursec{Corrigés des exercices}

\coursec{Corrigés des exercices — Leçon 24 : Présenter son CV dans le cadre d'un entretien d'embauche}

\begin{corrige}[Exercice 1 — QCM : Les éléments du CV]
\textbf{Bonne réponse : 2.} L'état civil complet (nom, prénoms, date et lieu de naissance, nationalité, situation matrimoniale, adresse, téléphone, email) est la rubrique \cle{obligatoire} : elle permet au recruteur d'identifier et de contacter le candidat.

\textbf{Pourquoi les autres propositions sont fausses :}
\begin{itemize}
\item \textbf{1.} La photo n'est pas obligatoire ; si elle est jointe, elle doit être professionnelle (fond neutre, tenue correcte), pas nécessairement en costume traditionnel.
\item \textbf{3.} Les loisirs sont une rubrique \emph{facultative} ; ils ne figurent que s'ils apportent une information valorisante (esprit d'équipe, engagement).
\item \textbf{4.} Le salaire souhaité ne se mentionne jamais dans un CV ; il se négocie éventuellement lors de l'entretien.
\end{itemize}

\textbf{Point de vigilance :} Seules les rubriques « État civil », « Formation », « Expérience professionnelle », « Compétences » sont strictement indispensables. Les autres (Photo, Loisirs, Références, Objectif) sont à adapter selon l'offre et la culture de l'entreprise.
\end{corrige}

\begin{corrige}[Exercice 2 — Vrai ou Faux justifié : Subjonctif et Conditionnel]
\begin{enumerate}
\item \textbf{Faux.} \emph{Examiniez} est bien au subjonctif présent, mais la justification est fausse : après « je veux que », le subjonctif exprime la \cle{volonté} (un souhait, un fait non réalisé), et non un fait certain. Le fait certain s'exprime à l'indicatif.
\item \textbf{Vrai.} \emph{Souhaiterais} est au conditionnel présent ; dans une lettre de motivation, le conditionnel a ici une valeur d'\cle{atténuation} et de politesse : il rend la demande moins directe, donc plus courtoise.
\item \textbf{Faux.} Après « il faut que », le verbe se met obligatoirement au \cle{subjonctif} (\emph{sois} = subjonctif présent de \emph{être}), car « il faut que » exprime la nécessité, qui appartient au domaine du non-réel. L'indicatif serait une faute de grammaire.
\item \textbf{Faux.} « Si j'avais le temps, je réviserais » (imparfait + conditionnel présent) exprime une hypothèse \cle{irréelle ou peu probable dans le présent} (l'énonciateur n'a pas le temps maintenant). Une hypothèse réalisable dans l'avenir s'exprime avec le présent + futur : « Si j'ai le temps, je réviserai mon CV ».
\end{enumerate}

\textbf{Point de vigilance :} Distinguer systématiquement : \emph{Si} + indicatif présent $\to$ futur (réel) ; \emph{Si} + imparfait $\to$ conditionnel présent (irréel présent) ; \emph{Si} + plus-que-parfait $\to$ conditionnel passé (irréel passé).
\end{corrige}

\begin{corrige}[Exercice 3 — Définitions et concepts clés]
\begin{enumerate}
\item \textbf{Le CV (Curriculum Vitae)} : texte fonctionnel bref (une page en général) qui présente de manière structurée l'état civil, la formation, les expériences professionnelles, les compétences et les centres d'intérêt d'un candidat, afin d'obtenir un entretien d'embauche.
\item \textbf{La lettre de motivation} : courrier administratif qui accompagne le CV ; le candidat y explique, de façon argumentée et polie (conditionnel de politesse), pourquoi il postule à l'offre et ce qu'il peut apporter à l'entreprise.
\item \textbf{L'entretien d'embauche} : échange oral formel entre un candidat et un ou plusieurs recruteurs (jury), destiné à vérifier l'adéquation entre le profil du candidat (compétences, savoir-être, motivation) et le poste à pourvoir.
\item \textbf{Le texte iconique appliqué au CV} : ensemble des éléments visuels qui accompagnent le texte et facilitent sa lecture : mise en page aérée, rubriques en gras, police lisible et unique (Calibri, Arial), puces, icônes discrètes, photo professionnelle éventuelle. L'image doit servir le contenu, jamais le masquer.
\end{enumerate}
\end{corrige}

\begin{corrige}[Exercice 4 — Application directe : Morphosyntaxe]
\begin{enumerate}
\item Je souhaite que vous \textbf{examiniez} ma candidature avec attention. \textit{(Subjonctif présent : après « souhaiter que », valeur de volonté.)}
\item Il est indispensable que le candidat \textbf{maîtrise} les logiciels comptables. \textit{(Subjonctif présent : tournure impersonnelle de nécessité.)}
\item Je \textbf{apprécierais} de pouvoir vous exposer mes motivations de vive voix. \textit{(Conditionnel présent : politesse, atténuation.)}
\item Il faut que je \textbf{joigne} mes diplômes à ce courrier. \textit{(Subjonctif présent : après « il faut que », nécessité.)}
\item Nous \textbf{serions} honorés de collaborer avec votre structure. \textit{(Conditionnel présent : politesse, formule de courtoisie.)}
\end{enumerate}

\textbf{Point de vigilance :} Ne jamais mettre l'indicatif après « il faut que » ni après « je souhaite que ». Attention à la confusion entre \emph{serions} (conditionnel) et \emph{soyons} (subjonctif) : ici « nous serions honorés » est une formule de politesse, pas une subordonnée après « que ».
\end{corrige}

\begin{corrige}[Exercice 5 — Complétion guidée : Lettre de motivation à trous]
\begin{enumerate}
\item Je \textbf{souhaiterais} vivement intégrer votre équipe. \textit{(Conditionnel de politesse.)}
\item Je \textbf{aimerais} mettre mes compétences à votre service. \textit{(Conditionnel de politesse.)}
\item Il est primordial que je \textbf{puisse} démontrer mon savoir-faire. \textit{(Subjonctif après tournure impersonnelle de nécessité.)}
\item Je \textbf{serais} reconnaissant de me recevoir à votre convenance. \textit{(Conditionnel de politesse ; formulation attendue : « je vous serais reconnaissant de bien vouloir me recevoir ».)}
\item Il faut que vous \textbf{receviez} mon dossier complet. \textit{(Subjonctif après « il faut que ».)}
\item Afin que mon dossier \textbf{soit} complet, je joins mes certificats. \textit{(Subjonctif après « afin que », but.)}
\item Je \textbf{voudrais} vous rencontrer pour détailler mon parcours. \textit{(Conditionnel de politesse.)}
\item Nous \textbf{pourrions} convenir d'une date ensemble. \textit{(Conditionnel : suggestion atténuée.)}
\item Une réponse que j'\textbf{espère} favorable. \textit{(Indicatif présent : « espérer » exprime un fait tenu pour probable, il se construit avec l'indicatif — piège classique !)}
\end{enumerate}

\textbf{Point de vigilance :} « J'espère que » se construit à l'\cle{indicatif}, contrairement à « je souhaite que » et « je veux que » qui demandent le subjonctif. C'est une erreur fréquente au Baccalauréat technique.
\end{corrige}

\begin{corrige}[Exercice 6 — Analyse comparative : Deux CV pour un poste de Secrétaire Comptable]
\textbf{Texte argumenté attendu (corrigé type) :}

Pour ce poste de Secrétaire Comptable, c'est la candidate A, Marie K., qui doit être retenue, car son profil répond point par point aux exigences de l'offre.

Sur le plan des compétences techniques, le poste exige la maîtrise de Sage et d'Excel : or Marie maîtrise parfaitement Sage 100 et Excel avancé (tableaux croisés dynamiques), tandis que Grace n'est que débutante sur Excel et ne possède que des connaissances comptables théoriques. Le diplôme de Marie (BTS Comptabilité) correspond exactement au métier, alors que le DUT de Gestion administrative de Grace est plus éloigné du cœur du poste.

Sur le plan de l'expérience, l'offre demande deux ans de pratique : Marie justifie de trois ans comme Comptable unique, avec des tâches directement transférables (facturation, déclarations fiscales, paie). Grace, elle, n'a que six mois de stage en secrétariat et une année de vente, sans lien avec la comptabilité.

Sur le plan du savoir-être, la rigueur et la discrétion exigées se lisent jusque dans la présentation des CV : celui de Marie est sobre, professionnel (police Calibri, photo pro), ce qui témoigne de son sérieux ; celui de Grace, coloré, avec police fantaisie et photo selfie, dénote une méconnaissance des codes professionnels. Enfin, les loisirs de Marie (bénévolat) valorisent son engagement, quand ceux de Grace n'apportent aucune information utile au recruteur.

En conclusion, par la formation, l'expérience, les compétences logicielles et la présentation professionnelle de sa candidature, Marie K. est la candidate la plus adaptée au poste.

\textbf{Barème indicatif (sur 10) :} choix clairement annoncé et argumenté (2 pts) ; mobilisation des 4 critères du poste (4 pts) ; comparaison explicite A/B (2 pts) ; organisation et correction de la langue (2 pts).
\end{corrige}

\begin{corrige}[Exercice 7 — Correction et réécriture : Un CV truffé d'erreurs]
\textbf{Les 10 erreurs identifiées et corrigées :}
\begin{enumerate}
\item « mbarga jean-pierre » $\to$ \textbf{MBARGA Jean-Pierre} (nom en majuscules, prénom avec majuscule).
\item Chronologie de la formation inversée : « 2023--2021 » $\to$ \textbf{2021--2023} (dates dans l'ordre croissant ; rubrique en ordre antichronologique : le diplôme le plus récent en premier).
\item « 2018--2015 » $\to$ \textbf{2015--2018}.
\item « Si j'aurais su, je ferait mieux » $\to$ double faute : on n'écrit jamais « si j'aurais » (conditionnel après \emph{si}) et « ferait » sans s. Formulation correcte et professionnelle : \textbf{« Si j'avais su, je ferais mieux »} — mais une telle phrase n'a pas sa place dans un CV : la remplacer par les tâches réalisées (ex. : \emph{maintenance préventive des lignes}).
\item « Il faut que je suis sérieux » $\to$ \textbf{« Il faut que je sois sérieux »} (subjonctif obligatoire après « il faut que ») ; là encore, à remplacer par une description de tâches concrètes.
\item « Anglais : Lu, Écrit, Parlé (Niveau BAC) » $\to$ mentionner un \textbf{niveau précis} (ex. : niveau B1, anglais technique lu) ; « niveau BAC » ne veut rien dire pour un recruteur.
\item Loisirs inadaptés : « Boire des verres, Dormir, Politique » $\to$ à supprimer (nuisibles ou à risque) ; ne garder que des loisirs valorisants (football = esprit d'équipe).
\item « Célibataire, 2 enfants » : information de la vie privée non pertinente ; on peut indiquer simplement la situation matrimoniale, sans détail.
\item Rubrique \textbf{OBJECTIF / titre du poste visé manquante} en tête de CV : ajouter par exemple « Électrotechnicien — Technicien de maintenance ».
\item Email peu professionnel ou rubrique incomplète : prévoir un email sobre (\emph{jp.mbarga@...}) et vérifier que le téléphone est complet ; la référence « M. Le Directeur de SONEL » doit être précise (nom, fonction, contact) et donnée avec son accord.
\end{enumerate}

\textbf{CV réécrit (modèle) :}
\begin{center}
\begin{tabularx}{\linewidth}{|X|}
\hline
\textbf{CURRICULUM VITAE — Technicien en Électrotechnique} \\
\hline
\textbf{MBARGA Jean-Pierre}, 25 ans \\
Quartier Mvog-Ada, Yaoundé — Tél. : 6XX XX XX XX — jp.mbarga@mail.fr \\
\hline
\textbf{FORMATION} \\
2021 -- 2023 : BTS Électrotechnique, Lycée Technique de Douala (Mention Assez Bien) \\
2015 -- 2018 : Baccalauréat F3, Lycée de Nkolbisson \\
\hline
\textbf{EXPÉRIENCE PROFESSIONNELLE} \\
2024 -- présent : Stagiaire en maintenance, SONEL Douala — maintenance préventive, dépannage des installations, lecture de schémas électriques. \\
2022 -- 2023 : Apprenti électricien, Atelier M. Toto, Yaoundé — câblage d'installations domestiques, pose de tableaux. \\
\hline
\textbf{COMPÉTENCES} \\
Câblage, dépannage, lecture de plans électriques ; Anglais technique lu (niveau B1). \\
\textbf{CENTRES D'INTÉRÊT} : Football (esprit d'équipe). \\
\textbf{RÉFÉRENCE} : M. N., Responsable maintenance, SONEL Douala (avec son accord). \\
\hline
\end{tabularx}
\end{center}
\end{corrige}

\begin{corrige}[Exercice 8 — Simulation orale : « Présentez-vous » (2 minutes)]
\textbf{1. Plan détaillé et répartition du temps (120 secondes) :}
\begin{center}
\begin{tabularx}{\linewidth}{|l|X|c|}
\hline
\rowcolor{popblueL} \textbf{Étape} & \textbf{Contenu} & \textbf{Temps} \\
\hline
Accroche & Nom, diplôme, poste visé, une phrase de positionnement (« ce que je peux vous apporter ») & 15 s \\
\hline
Parcours de formation & Diplôme principal, établissement, mention éventuelle, en lien avec le poste & 20 s \\
\hline
Expériences clés & 1 ou 2 expériences (stages) avec tâches concrètes et un résultat chiffré si possible & 30 s \\
\hline
Compétences & Techniques (logiciels, gestes métier) puis transverses (rigueur, travail en équipe) & 25 s \\
\hline
Projet professionnel & Ce que je veux apprendre et apporter dans cette entreprise précise & 20 s \\
\hline
Conclusion & Remerciement, ouverture (« je me tiens à votre disposition pour vos questions ») & 10 s \\
\hline
\end{tabularx}
\end{center}

\textbf{2. Trois pièges à éviter :}
\begin{itemize}
\item \textbf{Réciter le CV} mot à mot : l'entretien doit sélectionner et mettre en valeur, non répéter.
\item \textbf{Parler de sa vie privée} (famille, religion, opinions politiques) : rester sur le plan professionnel.
\item \textbf{Les hésitations et tics de langage} (« euh », « en fait », « voilà quoi ») : préparer, s'entraîner à voix haute, marquer des silences plutôt que combler.
\end{itemize}

\textbf{Conseils de langue :} employer le conditionnel de politesse (« je souhaiterais », « je pourrais ») et le subjonctif après « il faut que », « je souhaite que » ; parler à la première personne, phrases courtes, débit posé.
\end{corrige}

\begin{corrige}[Exercice 9 — Sujet Type Baccalauréat Technique : Production d'un CV complet]
\textbf{1. CV attendu (modèle de réponse) :}
\begin{center}
\begin{tabularx}{\linewidth}{|X|}
\hline
\textbf{Paul BIKOI — Conducteur de Travaux Junior (Génie Civil / Bâtiment)} \hfill \textit{(photo professionnelle)} \\
24 ans — Bonabéri, Douala — Tél. : 6XX XX XX XX — paul.bikoi@mail.cm — Permis B, mobile nationalement \\
\hline
\textbf{FORMATION} \\
2021 -- 2023 : BTS Génie Civil, option Bâtiment, Lycée Technique de Bafoussam. \\
\hline
\textbf{EXPÉRIENCE PROFESSIONNELLE} \\
Janvier -- Juin 2024 : Chef de chantier adjoint (stage), BTP EXCELLENCE, Douala. \\
— Suivi de l'exécution des travaux et contrôle de conformité aux plans. \\
— Gestion et approvisionnement des matériaux (réduction des ruptures de stock). \\
— Application et contrôle des règles de sécurité sur chantier. \\
\hline
\textbf{COMPÉTENCES TECHNIQUES} \\
DAO : AutoCAD, ArchiCAD — Planification de chantier : MS Project — Lecture de plans, métrés, suivi de planning. \\
\hline
\textbf{LANGUES} : Français (courant), Anglais technique lu et écrit (niveau B2). \\
\textbf{QUALITÉS} : Rigueur, sens du leadership, esprit d'équipe, disponibilité et mobilité. \\
\textbf{CENTRES D'INTÉRÊT} : Football (capitaine d'équipe), lecture technique. \\
\hline
\end{tabularx}
\end{center}

\textbf{2. Trois choix stratégiques (annotations attendues) :}
\begin{enumerate}
\item \textbf{Titre du CV calqué sur l'intitulé de l'offre} (« Conducteur de Travaux Junior ») : le recruteur voit immédiatement l'adéquation candidat/poste.
\item \textbf{Mots-clés de l'offre repris explicitement} : AutoCAD, MS Project, « anglais technique lu/écrit », « rigueur », « leadership », « mobilité nationale » — le CV répond point par point au profil recherché.
\item \textbf{Mise en valeur du stage} placé avant les compétences et détaillé en tâches concrètes (suivi d'exécution, matériaux, sécurité), car l'offre accepte « expérience stage/premier emploi » : c'est l'atout majeur d'un jeune diplômé.
\end{enumerate}

\textbf{Barème indicatif (sur 20) :} rubriques complètes et ordonnées (état civil, formation antichronologique, expérience, compétences, langues) : 6 pts ; ciblage de l'offre (mots-clés, titre) : 4 pts ; mise en page soignée, une page, photo mentionnée : 4 pts ; correction linguistique (orthographe, accords) : 4 pts ; annotations stratégiques pertinentes : 2 pts.

\textbf{Point de vigilance :} Dates dans l'ordre antichronologique (le plus récent d'abord) ; aucune information privée inutile ; pas de faute d'orthographe sur les noms propres et les logiciels ; le CV tient sur une page.
\end{corrige}

\begin{corrige}[Exercice 10 — Sujet Type Probatoire Technique : Grammaire en contexte + Production]
\textbf{Partie A : Transformation au conditionnel de politesse (8 points — 1 pt par phrase)}
\begin{enumerate}
\item Je \textbf{voudrais} savoir si le poste est toujours vacant.
\item Je \textbf{pourrais} commencer dès lundi prochain.
\item Je \textbf{souhaiterais} partir tôt vendredi pour un examen, si cela ne vous incommode pas. \textit{(« il faut que » ne s'atténue pas au conditionnel ; on reformule la demande.)}
\item Je \textbf{penserais} / j'\textbf{aurais tendance à penser} que mon profil correspond bien au poste. \textit{(Atténuation de l'affirmation trop sûre d'elle.)}
\item Nous \textbf{pourrions} signer le contrat la semaine prochaine.
\item \textbf{Souhaiteriez-vous} / \textbf{pourriez-vous} me donner une réponse rapide ? \textit{(« vous devez » est un ordre : on le transforme en demande polie.)}
\item Je \textbf{crois savoir} / j'\textbf{aurais cru comprendre} que l'entreprise a de grands projets.
\item Je \textbf{me présenterais} pour l'entretien à 10 h, si l'horaire vous convient. \textit{(Forme attendue : « je me présenterais ».)}
\end{enumerate}

\textbf{Point de vigilance :} Le conditionnel de politesse s'applique aux verbes de volonté et de pouvoir (\emph{vouloir, pouvoir, devoir, savoir}). Quand le verbe ne s'y prête pas (« il faut que », « vous devez »), on reformule la phrase entière : c'est la reformulation courtoise qui est évaluée, pas seulement la conjugaison.

\textbf{Partie B : Lettre de motivation (corrigé type, 12 points)}

\begin{quote}
Paul BIKOI \hfill Douala, le 20 novembre 2024\\
Bonabéri, Douala\\
Tél. : 6XX XX XX XX\\
À l'attention du Directeur des Ressources Humaines\\
Entreprise CONSTRUCT FUTURE, Yaoundé

\textbf{Objet : Candidature au poste de Conducteur de Travaux Junior}

Monsieur le Directeur,

Suite à votre annonce parue dans \emph{Mutations} du 15 novembre 2024, je voudrais vous proposer ma candidature au poste de Conducteur de Travaux Junior au sein de votre entreprise.

Titulaire d'un BTS de Génie Civil, option Bâtiment, obtenu en 2023 au Lycée Technique de Bafoussam, j'ai consolidé ma formation par un stage de six mois à l'entreprise BTP EXCELLENCE de Douala, en qualité de Chef de chantier adjoint. J'y ai assuré le suivi de l'exécution des travaux, la gestion des matériaux et le contrôle des règles de sécurité. Il faut que je précise que cette expérience m'a permis de maîtriser les réalités du chantier.

Votre entreprise, reconnue pour la qualité de ses réalisations, correspond exactement au cadre dans lequel je souhaiterais développer mes compétences. Je maîtrise les logiciels AutoCAD, ArchiCAD et MS Project, et je lis l'anglais technique (niveau B2). Rigoureux et doté d'un bon sens du leadership, je pourrais m'intégrer rapidement à vos équipes et je suis pleinement mobile sur le plan national.

Je serais honoré de vous exposer mes motivations lors d'un entretien, à votre convenance. Dans cette attente, je vous prie d'agréer, Monsieur le Directeur, l'expression de mes salutations distinguées.

\textit{Signature}\\
Paul BIKOI
\end{quote}

\textbf{Barème indicatif Partie B (12 pts) :} en-tête et objet conformes (2 pts) ; formules d'appel et de politesse (2 pts) ; structure en 3 paragraphes « Qui suis-je ? / Pourquoi vous ? / Pourquoi moi ? » (4 pts) ; emploi correct du conditionnel de politesse et du subjonctif (2 pts) ; correction de la langue et longueur respectée (2 pts).
\end{corrige}

\begin{corrige}[Exercice 11 — Sujet Type Bac : Intégration — Analyse de jury et Synthèse]
\textbf{1. Compte rendu synthétique du DRH (corrigé type, environ 15 lignes) :}

\begin{quote}
\textbf{Compte rendu d'entretien — Candidat : Paul BIKOI — Poste : Conducteur de Travaux Junior}

Monsieur le Directeur Général,

Le jury a reçu ce jour M. Paul Bikoi, candidat au poste de Conducteur de Travaux Junior. Sa présentation générale est excellente : tenue professionnelle, posture droite, contact visuel franc. Son discours, structuré selon un plan clair (parcours, stage, compétences, projet), a été tenu dans le temps imparti (1 min 55 s).

Trois points forts sont à retenir. Premièrement, des compétences techniques solides et vérifiées : réponses précises sur la DAO, le béton, la planification et la sécurité, illustrées par un exemple concret de gestion d'un conflit avec un fournisseur durant son stage. Deuxièmement, de réelles compétences transverses : lecture d'une notice en anglais technique, aisance relationnelle et bonne gestion du stress (silence maîtrisé de trois secondes avant une réponse difficile). Troisièmement, une motivation avérée : le candidat connaît nos projets, accepte la mobilité nationale et a posé des questions pertinentes sur la formation continue.

Deux points de vigilance : le candidat est jeune diplômé et son expérience reste limitée à un stage de six mois ; il devra être encadré par un conducteur principal durant sa première année. Par ailleurs, son anglais, correct à la lecture, devra être consolidé à l'oral.

Au vu de l'adéquation remarquable entre son profil et le poste, le jury émet un \textbf{avis favorable} à son recrutement, assorti d'une période de tutorat.

Le DRH
\end{quote}

\textbf{2. Trois exemples linguistiques tirés de l'oral supposé :}
\begin{itemize}
\item \textbf{Subjonctif de nécessité :} « Il faut que je \emph{sois} encadré la première année pour progresser rapidement. » (subjonctif obligatoire après « il faut que »).
\item \textbf{Subjonctif de souhait :} « Je souhaite que mon équipe \emph{respecte} les règles de sécurité sur le chantier. » (subjonctif après « souhaiter que »).
\item \textbf{Conditionnel d'atténuation (politesse) et hypothèse :} « Je \emph{souhaiterais} évoluer vers un poste de conducteur principal » (conditionnel de politesse) et « Si vous me \emph{confiez} ce chantier, je le \emph{mènerais} à terme dans les délais. » (conditionnel d'atténuation en principale, « si » au présent pour hypothèse réalisable).
\end{itemize}

\textbf{Barème indicatif (sur 20) :} structure du compte rendu respectée (introduction, 3 points forts, 2 points de vigilance, avis motivé) : 8 pts ; fidélité aux indicateurs de la grille : 4 pts ; avis clair et justifié : 2 pts ; exemples linguistiques corrects et bien identifiés : 4 pts ; correction de la langue : 2 pts.

\textbf{Point de vigilance :} Le compte rendu est un texte \cle{fonctionnel} : ton neutre et objectif, phrases factuelles appuyées sur la grille, pas de jugement personnel gratuit. L'avis doit être \emph{motivé} (jamais un simple « oui » ou « non »). Pour la question linguistique, il faut nommer le mode, le temps ET la valeur (nécessité, souhait, politesse, hypothèse).
\end{corrige}

\newpage

\coursec{Fiche bilan}

\courssub{Carte mentale de la leçon}
\begin{popfigure}
\begin{tikzpicture}[
    mindmap,
    concept color=popblue,
    level 1 concept/.append style={level distance=3.5cm, sibling angle=360/7},
    level 2 concept/.append style={level distance=2.5cm},
    text=white,
    font=\small\bfseries,
    every node/.append style={scale=0.85}
]
\node[concept, minimum size=2.8cm, align=center] (centre){Le CV \&\\ Entretien d'embauche}
    child[concept color=poporange] { node[concept, minimum size=2.2cm, align=center] (def){1. Définition \& Rôle\\ \textit{(Texte fonctionnel)}} }
    child[concept color=popgreen] { node[concept, minimum size=2.2cm, align=center] (struct){2. Structure\\ \textit{(Rubriques obligatoires)}} }
    child[concept color=poppurple] { node[concept, minimum size=2.2cm, align=center] (icon){3. Mise en page\\ \textit{(Textes iconiques)}} }
    child[concept color=poppink] { node[concept, minimum size=2.2cm, align=center] (lang){4. Langue\\ \textit{(Subjonctif / Conditionnel)}} }
    child[concept color=popgold] { node[concept, minimum size=2.2cm, align=center] (redac){5. Rédaction\\ \textit{(Production guidée)}} }
    child[concept color=popblue] { node[concept, minimum size=2.2cm, align=center] (oral){6. Présentation orale\\ \textit{(Entretien \& Jury)}} }
    child[concept color=popdark] { node[concept, minimum size=2.2cm, align=center] (int){7. Intégration\\ \textit{(Évaluation \& Remédiation)}} };
\end{tikzpicture}
\legende{Carte mentale récapitulative de la leçon 24 : les 7 piliers du CV et de sa présentation orale (programme MINESEC Terminale Tech).}
\end{popfigure}

\courssub{Synthèse des savoirs et savoir-faire}

\textbf{1. Définitions clés}
\begin{itemize}
    \item \cle{CV (Curriculum Vitae)} : Texte fonctionnel normé présentant le parcours (formation, expériences, compétences) d'un candidat pour une candidature.
    \item \cle{Texte iconique} : Ensemble des signes visuels (mise en page, police, puces, icônes, espacement) qui structurent l'information et guident la lecture rapide (diagonale) du recruteur.
    \item \cle{Entretien d'embauche} : Échange oral structuré où le candidat \emph{présente}, \emph{justifie} et \emph{défend} son CV face à un jury/recruteur.
\end{itemize}

\textbf{2. Structure type du CV (Ordre antéchronologique recommandé)}
\begin{center}
\begin{tabularx}{\linewidth}{|l|X|}\hline
\rowcolor{popblueL} \textbf{Rubrique} & \textbf{Contenu essentiel} \\ \hline
\cle{En-tête} & Nom, Prénom, Contact (tél., email pro), Localisation, Photo (optionnelle mais courante), Titre visé. \\ \hline
\cle{Profil / Objectif} & 2--3 lignes : identité pro + valeur ajoutée + poste ciblé (mots-clés de l'offre). \\ \hline
\cle{Formation} & Diplômes (du + récent au + ancien), Établissement, Année, Mentions, Mémoire/Projet pertinent. \\ \hline
\cle{Expériences} & Poste, Structure, Dates, \textbf{3--4 réalisations chiffrées} (verbes d'action : \textit{géré, conçu, augmenté, réduit}). \\ \hline
\cle{Compétences} & Techniques (logiciels, machines, langues \cle{CECRL}), Transverses (management, gestion de projet), Certifications. \\ \hline
\cle{Centres d'intérêt} & Activités montrant soft skills (sport collectif, bénévolat, veille techno), éviter les listes génériques. \\ \hline
\cle{Références} & « Disponibles sur demande » ou Nom/Fonction/Contact de 1--2 personnes (accord préalable). \\ \hline
\end{tabularx}
\end{center}

\textbf{3. Langue française : Valeurs des modes (Subjonctif vs Conditionnel)}
\begin{center}
\begin{tabularx}{\linewidth}{|l|X|X|}\hline
\rowcolor{popblueL} \textbf{Contexte / Fonction} & \textbf{Subjonctif} (Subjectivité, non-réalisé) & \textbf{Conditionnel} (Hypothèse, atténuation) \\ \hline
\cle{Volonté / Souhait} & \textit{Il faut que je \textbf{parte} à l'heure.} & \textit{Je \textbf{aimerais} partir à l'heure.} \\ \hline
\cle{Doute / Probabilité} & \textit{Je doute qu'il \textbf{soit} libre.} & \textit{Il \textbf{serait} libre (selon lui).} \\ \hline
\cle{Nécessité / Obligation} & \textit{Il est nécessaire que tu \textbf{relises} le CV.} & \textit{Tu \textbf{devrais} relire le CV.} \\ \hline
\cle{Politesse / Demande} & (Rare en principal) & \textit{Pourriez-vous \textbf{regarder} mon CV ?} \\ \hline
\cle{Hypothèse (Si...)} & (Pas en proposition principale) & \textit{Si j'avais le poste, je \textbf{m'investirais} pleinement.} \\ \hline
\end{tabularx}
\end{center}
\cle{Règle d'or entretien} : Conditionnel pour la politesse (\textit{Je souhaiterais, Je pourrais}) et l'hypothèse (\textit{Si j'intègre votre équipe...}) ; Subjonctif après verbes de volonté/sentiment/doute (\textit{Il faut que je maîtrise..., Je veux que ce poste me permette...}).

\textbf{4. Méthodes express}
\begin{methode}[Rédiger son CV en 5 étapes]
\begin{itemize}
    \item \textbf{Analyser l'offre :} Repérer mots-clés (compétences, valeurs, logiciels).
    \item \textbf{Sélectionner \& Hiérarchiser :} Ne garder que le pertinent pour \emph{ce} poste. Ordre antéchronologique.
    \item \textbf{Rédiger les « puces » :} Verbe d'action + Contexte + \textbf{Résultat chiffré} (ex. \textit{Géré un stock de 5000 références, réduit les pertes de 15\%}).
    \item \textbf{Maquetter (Iconique) :} Police unique (sans empattement), alignements parfaits, espace blanc, icônes discrètes (tél, mail, lieu), 1 page max (débutant).
    \item \textbf{Relire \& Valider :} Zéro faute, cohérence dates, format PDF, nom de fichier pro (\texttt{Nom\_Prenom\_CV.pdf}).
\end{itemize}
\end{methode}

\begin{methode}[Présenter son CV à l'oral (Entretien) en 4 temps]
\begin{itemize}
    \item \textbf{Accroche (1 min) :} « Bonjour, je suis [Prénom Nom], [Diplôme/Spécialité], je vise ce poste car [adéquation valeurs/compétences]. »
    \item \textbf{Parcours guidé (3--4 min) :} Suivre le fil du CV \textbf{sans le lire}. Choisir 2--3 expériences phares : \cle{Contexte $\to$ Action $\to$ Résultat (STAR)}.
    \item \textbf{Motivation \& Projection (2 min) :} Pourquoi \emph{cette} entreprise ? Pourquoi \emph{vous} ? (Conditionnel : \textit{J'apporterais..., Je m'investirais...}).
    \item \textbf{Questions / Conclusion (1 min) :} Poser 1 question pertinente. Remercier, serrer la main, sourire.
\end{itemize}
\end{methode}

\begin{savaistu}[Contexte camerounais : CV et insertion professionnelle]
Au Cameroun, le CV est exigé pour tout concours de la fonction publique (ENAM, ENSET, concours administratifs), pour les recrutements dans le secteur privé (banques, industries, ONG) et pour les stages de fin d'études (PFE). Les employeurs (GICAM, MINFOPRA, grandes entreprises) attendent un CV \textbf{clair, sans faute, une page}, avec une photo d'identité récente. Les langues officielles (français/anglais) sont un atout majeur : mentionnez votre niveau \cle{CECRL} (B2, C1). Adaptez le vocabulaire technique à la spécialité (Génie civil, Électrotechnique, Commerce, Secrétariat bureautique, etc.). Le réseau \cle{REJECAM} et les plateformes \textit{Emploi.cm}, \textit{Jobecam} sont des ressources utiles.
\end{savaistu}

\begin{aretenir}[Checklist avant le Bac / Probatoire]
    $\square$ Je sais définir le CV comme \textbf{texte fonctionnel} et \textbf{texte iconique}.
    $\square$ Je connais les \textbf{6 rubriques obligatoires} et leur ordre standard.
    $\square$ Je sais rédiger une \textbf{puce d'expérience} : Verbe d'action + Contexte + \textbf{Résultat chiffré}.
    $\square$ Je maîtrise la \textbf{mise en page} : police unique, alignements, icônes, 1 page, PDF.
    $\square$ Je distingue les valeurs du \textbf{Subjonctif} (volonté, doute, nécessité) et du \textbf{Conditionnel} (politesse, hypothèse, atténuation).
    $\square$ Je sais conjuguer aux temps requis (présent, passé) les verbes usuels d'entretien (\textit{pouvoir, vouloir, devoir, falloir, aimer, espérer}).
    $\square$ Je connais la structure en \textbf{4 temps de l'oral} : Accroche $\to$ Parcours (STAR) $\to$ Motivation $\to$ Questions.
    $\square$ J'applique le \cle{STAR} (Situation, Tâche, Action, Résultat) pour illustrer mes compétences.
    $\square$ J'ai préparé \textbf{2 questions pertinentes} à poser au jury (formation, équipe, projets).
    $\square$ Je gère le \textbf{non-verbal} : contact visuel, posture ouverte, voix audible, sourire.
    $\square$ Mon dossier de candidature (CV + Lettre + Pièces) est \textbf{complet, sans faute, cohérent} et prêt à être déposé.
\end{aretenir}

\findecours

\end{document}
